\documentclass[10pt]{article}
\usepackage[T1]{fontenc}
\usepackage[letterpaper, left=0.65in, textwidth=4.7in, marginparsep=0.3in, marginparwidth=2.45in, top=0.75in, bottom=0.75in]{geometry}
\usepackage{enumitem}
\usepackage{titlesec}
\usepackage{xcolor}
\usepackage{listings}
\usepackage{booktabs}
\usepackage{tabularx}
\usepackage[most]{tcolorbox}
\usepackage{multicol}
\usepackage[titles]{tocloft}
\usepackage{marginnote}
\usepackage{marginfix}
\usepackage{ragged2e}
\usepackage{hyperref}

\hypersetup{colorlinks=true,linkcolor=blue!50!black,urlcolor=blue!50!black}
\setlist{nosep,leftmargin=1.3em,topsep=2pt}
\setlength{\parindent}{0pt}
\setlength{\parskip}{4pt}
\setcounter{secnumdepth}{0}
\setcounter{tocdepth}{1}

% Heading Formats
\titleformat{\part}[block]{\LARGE\bfseries\color{blue!40!black}}{}{0pt}{}
\titlespacing*{\part}{0pt}{10pt}{2pt}
\titleformat{\section}{\Large\bfseries\color{blue!50!black}}{}{0pt}{}
\titlespacing*{\section}{0pt}{14pt}{4pt}
\titleformat{\subsection}{\normalsize\bfseries\itshape\color{black!80}}{}{0pt}{}
\titlespacing*{\subsection}{0pt}{7pt}{1pt}

% Editorial Table of Contents (Book Style)
\renewcommand{\contentsname}{Contents}
\renewcommand{\cfttoctitlefont}{\hfill\Large\bfseries\color{blue!50!black}}
\renewcommand{\cftaftertoctitle}{\hfill}
\setlength{\cftbeforetoctitleskip}{4pt}
\setlength{\cftaftertoctitleskip}{12pt}

% Parts: Structural rubrics without page numbers or dot leaders
\renewcommand{\cftpartfont}{\large\bfseries\color{blue!40!black}}
\setlength{\cftbeforepartskip}{11pt}
\setlength{\cftpartindent}{0pt}
\cftpagenumbersoff{part}
\renewcommand{\cftpartleader}{\hfill}

% Sections/Chapters: Indented, clean whitespace, flush-right page numbers
\renewcommand{\cftsecfont}{\normalsize\normalfont}
\renewcommand{\cftsecpagefont}{\normalsize\normalfont\color{black!75}}
\setlength{\cftsecindent}{1.4em}
\setlength{\cftsecnumwidth}{1.8em}
\setlength{\cftbeforesecskip}{4pt}
\renewcommand{\cftsecleader}{\hfill}
\cftsetpnumwidth{1.6em}
\cftsetrmarg{2.0em}

% Code Listing Settings (adapt language as needed)
\lstset{basicstyle=\ttfamily\footnotesize,columns=fullflexible,keepspaces=true,
  frame=l,framerule=1.5pt,rulecolor=\color{gray!50},xleftmargin=8pt,
  aboveskip=4pt,belowskip=4pt,showstringspaces=false,upquote=true,breaklines=true}
\newcommand{\code}[1]{{\ttfamily\small #1}}

% Collision-Free Sidenote Engine
\setlength{\marginparpush}{7pt}
\newcommand{\marginnotestyle}{\scriptsize\sffamily\color{black!80}\RaggedRight}

\makeatletter
\newcommand{\safemargin}[1]{%
  \ifinner%
    % Suppressed inside tables/boxes to avoid collisions and "Float(s) lost" crashes
  \else%
    \marginpar{\marginnotestyle #1}%
  \fi%
}
\makeatother

\newcommand{\term}[3][]{%
  \hypertarget{t:#2}{\textbf{#3}}%
  \if\relax\detokenize{#1}\relax\else%
    \safemargin{#1}%
  \fi%
}
\renewcommand{\footnote}[1]{\safemargin{#1}}
\newcommand{\ix}[3]{\hyperlink{t:#1}{#2}\dotfill #3\par}

% Socratic Callout Boxes
\newtcolorbox{corebox}{colback=blue!4,colframe=blue!45!black,boxrule=0.6pt,arc=2pt,
  left=6pt,right=6pt,top=3pt,bottom=3pt,fontupper=\small}
\newtcolorbox{ponder}{colback=green!3,colframe=green!35!black,boxrule=0.5pt,arc=2pt,
  left=6pt,right=6pt,top=2pt,bottom=2pt,fontupper=\small,
  title={\small\bfseries Think it through},coltitle=white,colbacktitle=green!35!black,
  toptitle=1pt,bottomtitle=1pt}

\newcommand{\chap}[2]{\section*{#1}\addcontentsline{toc}{section}{#1}\vspace{-2pt}\begin{corebox}\textbf{Core idea.} #2\end{corebox}}
\newcommand{\partq}[2]{\part*{#1}{\small\itshape #2}}

% ---------------------------------------------------------------------------
% Document-specific additions
\usepackage{lmodern}     % scalable fonts: avoids "Font ecsx0800 not found" on some TeX installs
\usepackage{needspace}   % \Needspace before short listings
\newcommand{\ch}[1]{{\small\color{black!55}[#1]}}   % source tag, e.g. \ch{Ch 13, p269}
\newcommand{\mantra}[1]{\par\smallskip{\color{blue!45!black}\textbf{Rule:}~\emph{#1}}\par}
\newcommand{\partline}[2]{\phantomsection\partq{#1}{#2}}  % titlesec adds the TOC entry
\usepackage{amssymb}
\setlength{\emergencystretch}{3em}  % let lines with unbreakable inline code stretch instead of overflowing
% JavaScript and TypeScript for listings, plus the non-ASCII characters the code uses.
\lstdefinelanguage{JavaScript}{
  morekeywords={async,await,break,case,catch,class,const,continue,default,delete,do,else,export,extends,
    false,finally,for,from,function,get,if,import,in,instanceof,let,new,null,of,return,set,static,super,
    switch,this,throw,true,try,typeof,undefined,var,void,while,yield},
  sensitive=true,morecomment=[l]{//},morecomment=[s]{/*}{*/},
  morestring=[b]",morestring=[b]',morestring=[b]`}
\lstdefinelanguage{TypeScript}[]{JavaScript}{
  morekeywords={type,interface,enum,namespace,implements,private,protected,public,readonly,keyof,
    unknown,never,any,number,string,boolean,satisfies,as,is,declare,abstract}}
\lstset{keywordstyle=\color{blue!55!black}\bfseries,commentstyle=\color{green!35!black}\itshape,
  stringstyle=\color{red!45!black},extendedchars=true,
  literate={°}{{\textdegree}}1 {π}{{$\pi$}}1 {²}{{$^2$}}1 {×}{{$\times$}}1 {–}{{--}}1 {—}{{---}}1
    {…}{{\ldots}}1 {→}{{$\rightarrow$}}1 {≥}{{$\geq$}}1}
% Card-by-card solution guide
\newcommand{\trap}{\par\raggedright\textbf{\color{red!70!black}Trap:}~}
\newcommand{\core}[1]{\textit{#1}\par}
\newcommand{\practice}{\textbf{\color{blue!50!black}Practice:}~}
\newcommand{\deckhead}[2]{\par\ifdim\dimexpr\pagegoal-\pagetotal\relax<0.3\textheight\newpage\fi\bigskip\phantomsection\addcontentsline{toc}{section}{#1}{\Large\bfseries\color{blue!50!black}#1}\hfill{\small #2 cards}\par\nopagebreak\vspace{3pt}}
\newcommand{\cardhead}[2]{\par\ifdim\dimexpr\pagegoal-\pagetotal\relax<0.28\textheight\newpage\fi\bigskip\noindent{\color{blue!50!black}\rule{\linewidth}{0.8pt}}\par\nopagebreak\noindent{\large\bfseries\color{blue!50!black}#1:}\enspace{\ttfamily #2}\par\nopagebreak}
\newcommand{\cardsub}[1]{\par\medskip\noindent{\small\bfseries\scshape\color{black!70}#1}\par\nopagebreak\vspace{1pt}}
\newcommand{\readchap}[2]{\par\medskip\noindent{\bfseries\color{blue!40!black}#1}\par\nopagebreak\noindent{\small\textit{Core idea:} #2}\par\nopagebreak}
\newcommand{\readq}[1]{\par\smallskip\noindent{\bfseries\itshape\small #1}\par\nopagebreak}
\newcommand{\cardend}{\par\medskip}
\newtcolorbox{sumbox}[1]{breakable,colback=orange!3,colframe=orange!45!black,boxrule=0.5pt,arc=2pt,
  left=6pt,right=6pt,top=3pt,bottom=3pt,fontupper=\small,title={\small\bfseries #1},coltitle=white,
  colbacktitle=orange!45!black,toptitle=1pt,bottomtitle=1pt}
\newlist{cardrows}{itemize}{1}
\setlist[cardrows]{label={},leftmargin=5.3em,labelwidth=4.9em,labelsep=0.4em,labelindent=0pt,align=left,font=\small\bfseries,before={\small\raggedright},itemsep=0.5pt,topsep=2pt}
% ---------------------------------------------------------------------------
\begin{document}
{\noindent\raggedright\LARGE\bfseries\color{blue!40!black} Part 1 · JavaScript Fundamentals\par}
\vspace{2pt}{\large A Socratic study guide to the 64 cards of Part 1 of “JavaScript \& TypeScript: Zero to Mastery”}\par\vspace{6pt}
\begin{corebox}\textbf{The question this guide answers.} What do you need to know to write JavaScript that does what you mean?\par\smallskip
Part A asks what a program is made of (values, names, comparisons); B, how code repeats and handles text; C, how logic is packaged into functions; D, how many values are held (arrays, objects, Map and Set); E, how nested problems are solved and failures survived (recursion, errors, JSON); F, how bigger programs fit together (classes, asynchronous code, modules). Every section heading is a question, answered in its first sentence. The Card-by-Card Solution Guide then works every card in the deck.\end{corebox}
\vspace{6pt}\tableofcontents\vspace{6pt}
\newpage
\partline{A. What is a program made of?}{Values, the names that point at them, and how they are compared.}
\phantomsection\label{ch:c01}
\chap{1. What is a value, and how do you name one?}{Every piece of data is a value of some type; a variable is a name that points at a value, and \code{const} pins the name, not the value.}

\subsection{What kinds of values are there?}% sec:c01-types
JavaScript has seven primitive types and objects. A \term[\textbf{Value:} one piece of data the program works with, such as \code{42} or \code{"hi"}. Like a coin: you can hand it around and look at it.]{value}{value} is either a \term[\textbf{Primitive:} a simple value that cannot be changed: a string, number, bigint, boolean, undefined, symbol or null. Like a number stamped on a coin.]{primitive}{primitive}---string, number, bigint, boolean, \code{undefined}, symbol, \code{null}---or an \term[\textbf{Object:} a value that holds other values under names and can be changed. Like a filing cabinet whose drawers you can open, fill and empty.]{object}{object}. Arrays and functions are objects too.

The \term[\textbf{typeof:} an operator that returns a string naming a value's type. Like a label maker that reads a value and prints what kind it is.]{typeof}{typeof} operator names a value's type. It has two traps: \code{typeof null} is \code{"object"}, and an array is reported as \code{"object"} as well. Test for \code{null} with \code{=== null} and for arrays with \code{Array.isArray}.
\Needspace{6\baselineskip}
\begin{lstlisting}[language=JavaScript]
console.log(typeof 3.5, typeof "hi", typeof true);
console.log(typeof {}, typeof [], Array.isArray([]));
console.log(typeof null, null === null);
\end{lstlisting}
{\small\textit{Prints:}~\code{number string boolean}\enspace\textcolor{black!45}{|}\enspace \code{object object true}\enspace\textcolor{black!45}{|}\enspace \code{object true}}\par

\subsection{How do you turn text into a number safely?}% sec:c01-numbers
Use \code{Number(text)}, then check the result. \code{Number} reads the whole string: \code{Number("12abc")} is \code{NaN}, the special ``not a number'' value. \code{parseInt} stops at the first character it cannot read, so \code{parseInt("12abc", 10)} is \code{12}. Watch empty text: \code{Number("")} and \code{Number("   ")} are \code{0}. \code{Number.isInteger} then tells you whether you got a whole number.
\Needspace{6\baselineskip}
\begin{lstlisting}[language=JavaScript]
console.log(Number("42"), Number("12abc"), Number(""));
console.log(parseInt("12abc", 10), Number.isInteger(3.5));
\end{lstlisting}
{\small\textit{Prints:}~\code{42 NaN 0}\enspace\textcolor{black!45}{|}\enspace \code{12 false}}\par

\subsection{What do let, const and var actually fix?}% sec:c01-names
A \term[\textbf{Variable:} a name that points at a value. Like a name tag tied to a parcel: you can move the tag to another parcel (reassign) or open the parcel (change an object).]{variable}{variable} is a name pointing at a value. \code{let} and \code{const} names exist only inside the block \code{\{ ... \}} that declares them; \code{const} also forbids pointing the name at a different value. It does not freeze the value: an object behind a \code{const} can still be changed.
\Needspace{6\baselineskip}
\begin{lstlisting}[language=JavaScript]
const team = ["Ada"];
team.push("Grace");
console.log(team.length);
\end{lstlisting}
{\small\textit{Prints:}~\code{2}}\par

\code{var} is older: its name exists, as \code{undefined}, from the top of its function. A \code{let} or \code{const} name cannot be read before its declaration line runs; that stretch is the \term[\textbf{Temporal dead zone (TDZ):} the part of a block before a \code{let}/\code{const} line runs, where reading the name throws a ReferenceError. Like a roped-off area before the shop opens.]{tdz}{temporal dead zone}, and reading it there throws a \code{ReferenceError}.

\subsection{What happens when you pass a value to a function?}% sec:c01-sharing
The function gets a copy of the value. For an object, that value is a \term[\textbf{Reference:} the value that stands for an object: where it lives, not the object itself. Like a street address on a card, not a duplicate house.]{reference}{reference}: the function and the caller now hold the same address. Changing a property through it is seen by the caller; pointing the parameter at a new object is not.
\Needspace{6\baselineskip}
\begin{lstlisting}[language=JavaScript]
function addTag(list) {
  list.push("new");
  list = ["other"];
}
const tags = ["old"];
addTag(tags);
console.log(tags.join(","));
\end{lstlisting}
{\small\textit{Prints:}~\code{old,new}}\par

\mantra{A name points; const pins the name, not the thing.}
\begin{ponder}A function receives an array and calls \code{arr.push(1)}. Another does \code{arr = [...arr, 1]}. Which change does the caller see, and why?\end{ponder}

\phantomsection\label{ch:c02}
\chap{2. How does JavaScript compare values and choose a branch?}{\code{===} compares without converting; truthiness decides an \code{if}; and \code{||}, \code{\&\&} and \code{??} hand back one of their operands.}

\subsection{Why prefer === over ==?}% sec:c02-equality
Because \code{===} never converts. \code{==} first applies \term[\textbf{Coercion:} JavaScript converting a value to another type on its own, e.g. \code{""} to \code{0}. Like a currency exchange that happens before you can compare prices.]{coercion}{coercion}, so \code{0 == ""} is \code{true}. The \term[\textbf{Strict equality (===):} comparison that treats values of different types as different. Like comparing two prices only if they are in the same currency.]{strict}{strict equality} operator treats operands of different types as different. Two other surprises: \code{null == undefined} is \code{true} but \code{null == 0} is not, and \code{NaN} equals nothing, not even itself---use \code{Number.isNaN}.
\Needspace{6\baselineskip}
\begin{lstlisting}[language=JavaScript]
console.log(1 == "1", 1 === "1");
console.log(null == undefined, null == 0);
console.log(NaN === NaN, Number.isNaN(NaN));
\end{lstlisting}
{\small\textit{Prints:}~\code{true false}\enspace\textcolor{black!45}{|}\enspace \code{true false}\enspace\textcolor{black!45}{|}\enspace \code{false true}}\par

\subsection{What counts as true in an if?}% sec:c02-truthy
An \code{if} runs its block when the condition is \term[\textbf{Truthy / falsy:} whether a value counts as true or false in a condition. Like a bouncer's list: eight values are turned away, everyone else gets in.]{truthy}{truthy}. Exactly eight values are falsy: \code{false}, \code{0}, \code{-0}, \code{0n}, \code{""}, \code{null}, \code{undefined} and \code{NaN}. Everything else is truthy, including \code{[]}, \code{\{\}} and the strings \code{"0"} and \code{"false"}.

\subsection{How do you write a chain of decisions?}% sec:c02-branches
With \code{if} statements that return early, checked from the most demanding case down. Each later \code{if} then only sees what the earlier ones let through, so every boundary is decided once. Test each boundary value itself.
\Needspace{6\baselineskip}
\begin{lstlisting}[language=JavaScript]
function label(temp) {
  if (temp >= 30) return "hot";
  if (temp >= 15) return "mild";
  return "cold";
}
console.log(label(30), label(29), label(14));
\end{lstlisting}
{\small\textit{Prints:}~\code{hot mild cold}}\par

\subsection{What do ||, \&\& and ?? really return?}% sec:c02-logical
One of their operands, not \code{true} or \code{false}. \code{a || b} returns \code{a} if it is truthy, else \code{b}; \code{a \&\& b} returns \code{a} if it is falsy, else \code{b}. The \term[\textbf{Nullish coalescing (??):} returns its right side only when its left side is \code{null} or \code{undefined}. Like a substitute who steps in only if the player is absent, not if the player scores zero.]{nullish}{nullish coalescing} operator \code{??} returns its right side only for \code{null} or \code{undefined}, so a real \code{0} or \code{""} survives. That makes it the right tool for defaults.
\Needspace{6\baselineskip}
\begin{lstlisting}[language=JavaScript]
const saved = 0;
console.log(saved || 10, saved ?? 10, null ?? 10);
\end{lstlisting}
{\small\textit{Prints:}~\code{10 0 10}}\par

\mantra{Compare with ===; default with ??.}
\begin{ponder}An untouched form field holds \code{null}; a cleared one holds \code{""}. Should a blank (cleared) name become \code{"Anonymous"}? Which of \code{||} and \code{??} gives the behaviour you chose?\end{ponder}

\partline{B. How does code repeat, and how is text handled?}{Loops, and the string methods that loops often drive.}
\phantomsection\label{ch:c03}
\chap{3. How do you make code repeat?}{A loop repeats a block while a condition holds: \code{for} to count, \code{for...of} to walk values, and \code{return} or \code{break} to stop early.}

\subsection{How does a counting for loop work?}% sec:c03-for
A \term[\textbf{Loop:} code that runs a block again and again. Like laps on a track with a lap counter.]{loop}{loop} repeats a block. A \code{for} loop has three parts: start (\code{let i = 1}), keep-going test (\code{i <= n}) and step (\code{i++}). Each run of the block is an \term[\textbf{Iteration:} one run of a loop's body. One lap.]{iteration}{iteration}. Write \code{<=} to include the last value; \code{<} stops one short---the classic \term[\textbf{Off-by-one error:} a loop that runs one time too many or too few. Like counting fence posts instead of gaps.]{offbyone}{off-by-one error}.

The remainder operator \code{\%} gives what is left after division, so \code{i \% 3 === 0} means ``divisible by 3''.
\Needspace{6\baselineskip}
\begin{lstlisting}[language=JavaScript]
let total = 0;
for (let i = 1; i <= 4; i++) total += i;
console.log(total, 9 % 3 === 0, 10 % 3);
\end{lstlisting}
{\small\textit{Prints:}~\code{10 true 1}}\par

\subsection{How do you walk through the values of an array or string?}% sec:c03-forof
With \code{for...of}, which hands you each value in turn. On a string it gives one character at a time.
\Needspace{6\baselineskip}
\begin{lstlisting}[language=JavaScript]
for (const ch of "hey") console.log(ch.toUpperCase());
\end{lstlisting}
{\small\textit{Prints:}~\code{H}\enspace\textcolor{black!45}{|}\enspace \code{E}\enspace\textcolor{black!45}{|}\enspace \code{Y}}\par

\subsection{How do you stop a loop early?}% sec:c03-exit
\code{return} ends the whole function, loop included; \code{break} ends only the loop. To find the first match, return as soon as you see it; if you only record matches, you end up with the last one.
\Needspace{6\baselineskip}
\begin{lstlisting}[language=JavaScript]
for (const n of [3, 8, 1, 9]) {
  if (n > 5) {
    console.log("first big:", n);
    break;
  }
}
\end{lstlisting}
{\small\textit{Prints:}~\code{first big: 8}}\par

\subsection{Why does let matter inside a loop?}% sec:c03-let
A function made inside a loop remembers the variable it can see. That pairing of a function with the variables around it is a \term[\textbf{Closure:} a function together with the variables it could see when it was created. Like a snow globe: the scene inside is sealed in with it.]{closure}{closure}. \code{let} creates a fresh variable for each iteration, so each function keeps its own; \code{var} makes one variable for the whole loop, so every function sees its final value.
\Needspace{6\baselineskip}
\begin{lstlisting}[language=JavaScript]
const fns = [];
for (let i = 0; i < 3; i++) fns.push(() => i);
console.log(fns.map((f) => f()).join(" "));
\end{lstlisting}
{\small\textit{Prints:}~\code{0 1 2}}\par

\subsection{Why is for...in the wrong loop for arrays?}% sec:c03-forin
\code{for...in} walks an object's keys, as strings, including extra and inherited ones. On an array that means \code{"0"}, \code{"1"}, \ldots{} plus any property someone added. Use \code{for...of} for an array's values.

\mantra{Count with for, walk with for...of, leave with return.}
\begin{ponder}A loop runs \code{for (let i = 0; i <= arr.length; i++)}. What is \code{arr[i]} on the last lap, and why does JavaScript not stop you?\end{ponder}

\phantomsection\label{ch:c04}
\chap{4. How do you work with text?}{Strings never change: every string method returns a new string, and \code{+} joins text as soon as one side is a string.}

\subsection{Can you change a string in place?}% sec:c04-immutable
No. Strings are \term[\textbf{Immutable:} cannot be changed after it is made; you make a new one instead. Like words printed in ink: to fix a typo you print a new copy.]{immutable}{immutable}. A \term[\textbf{Method:} a function attached to a value and called with a dot, as in \code{word.toUpperCase()}. Like a button on an appliance.]{method}{method} such as \code{toUpperCase} or \code{slice} returns a new string; the original stays as it was. \code{slice(1)} gives everything after the first character. Reading \code{word[0]} on an empty string gives \code{undefined}, and calling a method on \code{undefined} throws.
\Needspace{6\baselineskip}
\begin{lstlisting}[language=JavaScript]
const word = "hello";
const shout = word.toUpperCase();
console.log(word, shout, word.slice(1, 3));
\end{lstlisting}
{\small\textit{Prints:}~\code{hello HELLO el}}\par

\subsection{How do you break text apart and put it back?}% sec:c04-split
\code{split} cuts a string into an array of pieces at a pattern; \code{join} glues an array back into a string. \code{trim} removes spaces from both ends. Splitting on a single space leaves empty pieces between repeated spaces; the pattern \code{/\textbackslash s+/} (one or more whitespace characters) does not. Strings have no \code{reverse}: spread the string into an array, reverse that, and join.
\Needspace{6\baselineskip}
\begin{lstlisting}[language=JavaScript]
const parts = "  a   b ".trim().split(/\s+/);
console.log(parts.length, parts.join("-"));
console.log([..."abc"].reverse().join(""));
\end{lstlisting}
{\small\textit{Prints:}~\code{2 a-b}\enspace\textcolor{black!45}{|}\enspace \code{cba}}\par

\subsection{Why does "5" + 3 give "53"?}% sec:c04-plus
Because \code{+} both adds numbers and joins strings: as soon as either side is a string, it joins. The other arithmetic operators (\code{-}, \code{*}, \code{/}) convert their operands to numbers. Evaluation runs left to right, so \code{1 + 2 + "3"} adds first and joins second.
\Needspace{6\baselineskip}
\begin{lstlisting}[language=JavaScript]
console.log("5" + 3, "5" - 3, 1 + 2 + "3");
\end{lstlisting}
{\small\textit{Prints:}~\code{53 2 33}}\par

\mantra{Strings are read-only: build a new one.}
\begin{ponder}\code{capitalize} must return \code{""} for \code{""}. Would \code{word.charAt(0)} instead of \code{word[0]} remove the need for the empty-string check? What does \code{"".charAt(0)} give?\end{ponder}

\partline{C. How do you package logic?}{Functions, hoisting and closures.}
\phantomsection\label{ch:c05}
\chap{5. How do you package logic so it can be reused?}{A function bundles steps behind a name; declarations are usable before their line, and a closure keeps the variables it was created with.}

\subsection{What is a function, and what does hoisting change?}% sec:c05-define
A \term[\textbf{Function:} a named, reusable block of steps that takes inputs and can return a result. Like a recipe card: write it once, cook from it many times.]{function}{function} takes \term[\textbf{Parameter:} a name in a function's definition that receives an input. A blank on the recipe card, e.g. ``\_\_ cups of flour''.]{parameter}{parameters}, runs its body and hands back a value with \code{return}. A function declaration is \term[\textbf{Hoisting:} JavaScript making a declaration usable before the line it is written on. Like a speaker whose name is on the programme before they walk on stage.]{hoisting}{hoisted}: you can call it above the line where it is written. A function stored in a \code{const} is not usable until that line runs.
\Needspace{6\baselineskip}
\begin{lstlisting}[language=JavaScript]
console.log(double(4));
function double(n) {
  return n * 2;
}
\end{lstlisting}
{\small\textit{Prints:}~\code{8}}\par

\subsection{How can a function keep private state?}% sec:c05-closures
By closing over a variable (\hyperref[ch:c03]{Ch.~3}). Each call to an outer function creates new variables; the inner functions it returns keep using those, and nothing outside can reach them. An \term[\textbf{Arrow function:} the short form \code{(x) => x * 2} of a function expression. Like a sticky note instead of a recipe card.]{arrow}{arrow function} is a compact way to write such an inner function.
\Needspace{6\baselineskip}
\begin{lstlisting}[language=JavaScript]
function makeGreeter(greeting) {
  return (name) => `${greeting}, ${name}`;
}
const hi = makeGreeter("Hi");
const yo = makeGreeter("Yo");
console.log(hi("Ada"), yo("Ada"));
\end{lstlisting}
{\small\textit{Prints:}~\code{Hi, Ada Yo, Ada}}\par

To remember ``already done'', keep a separate flag. Using the stored result as the flag fails when the result is itself \code{undefined}.

\subsection{How does a function accept any number of arguments?}% sec:c05-rest
With a \term[\textbf{Rest parameter:} \code{...name} as the last parameter, gathering all remaining arguments into an array. Like a sack that holds however many apples you hand over.]{rest}{rest parameter}: \code{function sum(...numbers)} receives every argument in a real array. Adding them up with \code{reduce} (\hyperref[ch:c06]{Ch.~6}) needs a starting value, or zero arguments would throw.
\Needspace{6\baselineskip}
\begin{lstlisting}[language=JavaScript]
function count(...items) {
  return items.length;
}
console.log(count(), count("a", "b"));
\end{lstlisting}
{\small\textit{Prints:}~\code{0 2}}\par

\mantra{Name it once, call it anywhere; let closures carry state.}
\begin{ponder}You call \code{makeCounter(0)} twice. How many count variables now exist, and which one does each \code{increment} change? What if \code{count} were declared outside \code{makeCounter}?\end{ponder}

\partline{D. How do you hold many values?}{Arrays, objects, Maps and Sets.}
\phantomsection\label{ch:c06}
\chap{6. How do you work with a list of values?}{\code{map} transforms, \code{filter} keeps, \code{reduce} folds, \code{forEach} only acts; \code{sort} changes the array in place and compares as text unless told otherwise.}

\subsection{How do you transform or select items?}% sec:c06-map
An \term[\textbf{Array:} an ordered list of values, numbered from 0. Like a row of numbered lockers.]{array}{array} has methods that take a \term[\textbf{Callback:} a function you pass to another function, which calls it for you. Like instructions you hand a helper: ``do this with each item''.]{callback}{callback}. \code{map} returns a new array of the same length with each item transformed; \code{filter} returns a new array of the items whose callback returned a truthy value.
\Needspace{6\baselineskip}
\begin{lstlisting}[language=JavaScript]
const prices = [5, 20, 8];
const big = prices.filter((p) => p > 6).map((p) => p * 2);
console.log(big.join(" "), prices.join(" "));
\end{lstlisting}
{\small\textit{Prints:}~\code{40 16 5 20 8}}\par

\subsection{How do you boil a list down to one value?}% sec:c06-reduce
With \code{reduce}. It carries an \term[\textbf{Accumulator:} the running result \code{reduce} passes from one item to the next. Like a running tally on a notepad.]{accumulator}{accumulator} through the array: each callback gets the tally so far and the next item and returns the new tally. Give it a starting value (\code{0}, \code{\{\}}, \ldots): without one, an empty array throws a \code{TypeError}. The callback must return the accumulator every time.
\Needspace{6\baselineskip}
\begin{lstlisting}[language=JavaScript]
const add = (sum, n) => sum + n;
console.log([4, 6].reduce(add, 0), [].reduce(add, 0));
try {
  [].reduce(add);
} catch (error) {
  console.log(error.name);
}
\end{lstlisting}
{\small\textit{Prints:}~\code{10 0}\enspace\textcolor{black!45}{|}\enspace \code{TypeError}}\par

\subsection{When is forEach the wrong tool?}% sec:c06-foreach
When you want a result. \code{forEach} runs the callback for each item and always returns \code{undefined}; it is for side effects such as adding to a total. To build a new array, use \code{map}.
\Needspace{6\baselineskip}
\begin{lstlisting}[language=JavaScript]
const result = [1, 2].forEach((n) => n * 2);
console.log(result, [1, 2].map((n) => n * 2).join(","));
\end{lstlisting}
{\small\textit{Prints:}~\code{undefined 2,4}}\par

\subsection{Why does sort put 100 before 9?}% sec:c06-sort
Without a compare function, \code{sort} compares items as strings, and \code{"100"} comes before \code{"9"}. Pass \code{(a, b) => a - b} to sort numbers ascending. \code{sort} also changes the array \term[\textbf{In place:} changing the original rather than making a new copy. Like rearranging the books on your shelf instead of buying a sorted set.]{inplace}{in place}; copy first (\code{[...nums]}) if the caller's array must stay as it was.
\Needspace{6\baselineskip}
\begin{lstlisting}[language=JavaScript]
console.log([10, 9, 100].sort().join(" "));
console.log([10, 9, 100].sort((a, b) => a - b).join(" "));
\end{lstlisting}
{\small\textit{Prints:}~\code{10 100 9}\enspace\textcolor{black!45}{|}\enspace \code{9 10 100}}\par

\subsection{Is a copy of an array a new array?}% sec:c06-copy
Only if you make one. \code{const b = a} copies the reference (\hyperref[ch:c01]{Ch.~1}), so both names point at one array. \code{[...a]} builds a new array with the same items; the items themselves are shared, not copied.
\Needspace{6\baselineskip}
\begin{lstlisting}[language=JavaScript]
const a = [1];
const b = a;
const copy = [...a];
b.push(2);
console.log(a.length, copy.length, a === b);
\end{lstlisting}
{\small\textit{Prints:}~\code{2 1 true}}\par

\mantra{Return from map, act in forEach, start reduce with a value.}
\begin{ponder}A \code{reduce} callback returns the accumulator on one branch but forgets to on another. What is the accumulator for the next item after the forgetful branch runs?\end{ponder}

\phantomsection\label{ch:c07}
\chap{7. How do you hold named data?}{An object maps property names to values; destructuring unpacks it, spread copies one level of it, \code{in} asks whether a property exists, and \code{?.} stops at a missing link.}

\subsection{How do you pull properties out into variables?}% sec:c07-destructure
An object stores values under names; each is a \term[\textbf{Property:} a named value inside an object, like \code{user.name}. A labelled drawer in the filing cabinet.]{property}{property}. \term[\textbf{Destructuring:} syntax that unpacks properties into variables, as in \code{const \{ a, b \} = obj}. Like emptying labelled drawers onto matching labelled trays.]{destructuring}{Destructuring} copies properties into variables in one line, with renames (\code{\{ a: b \}}) and defaults (\code{\{ role = "guest" \}}). A default is used when the property is missing or \code{undefined}---not when it is \code{null}.
\Needspace{6\baselineskip}
\begin{lstlisting}[language=JavaScript]
const { name, role = "guest" } = { name: "Ada", role: null };
const { title = "none" } = {};
console.log(name, role, title);
\end{lstlisting}
{\small\textit{Prints:}~\code{Ada null none}}\par

\subsection{How do you copy or merge objects?}% sec:c07-spread
With spread: \code{\{ ...a, ...b \}} builds a new object from both; when a name appears twice, the later value wins. So spread the defaults first and the overrides second. Spread makes a \term[\textbf{Shallow copy:} a copy of the top level only; anything nested is shared, not copied. Like photocopying a folder's index page: the pages it lists are still the originals.]{shallow}{shallow copy}: a nested array or object is shared between the copy and the original.
\Needspace{6\baselineskip}
\begin{lstlisting}[language=JavaScript]
const base = { theme: "light", tags: ["a"] };
const mine = { ...base, theme: "dark" };
mine.tags.push("b");
console.log(mine.theme, base.theme, base.tags.join(""));
\end{lstlisting}
{\small\textit{Prints:}~\code{dark light ab}}\par

\subsection{How do you tell a missing property from an empty one?}% sec:c07-in
Ask \code{key in obj}: it is \code{true} when the property exists, whatever its value. Testing the value instead (\code{if (obj[key])}) wrongly treats \code{null}, \code{0} and \code{""} as missing.
\Needspace{6\baselineskip}
\begin{lstlisting}[language=JavaScript]
const box = { count: 0 };
console.log("count" in box, Boolean(box.count), "size" in box);
\end{lstlisting}
{\small\textit{Prints:}~\code{true false false}}\par

\subsection{How do you read a property that might not be there?}% sec:c07-optional
With \term[\textbf{Optional chaining (?.):} reads a property but returns \code{undefined} instead of throwing when the value before it is \code{null} or \code{undefined}. Like a hallway with ``if the door is locked, turn back'' signs.]{optional}{optional chaining}. \code{a?.b} is \code{undefined} when \code{a} is \code{null} or \code{undefined}; plain \code{a.b} would throw a \code{TypeError}. Combine it with \code{??} to supply a fallback.
\Needspace{6\baselineskip}
\begin{lstlisting}[language=JavaScript]
const user = { profile: null };
console.log(user.profile?.email ?? "no email");
\end{lstlisting}
{\small\textit{Prints:}~\code{no email}}\par

\mantra{Ask in whether it exists; ask ?? what to use instead.}
\begin{ponder}\code{withDefaults} spreads the defaults, then the options. If \code{options} is \code{\{ fontSize: undefined \}}, which \code{fontSize} ends up in the result? Would a destructuring default behave the same way?\end{ponder}

\phantomsection\label{ch:c08}
\chap{8. Which collection fits: object, Map or Set?}{A \code{Set} keeps one of each value; a \code{Map} pairs any key with a value in insertion order; a plain object turns keys into strings and inherits keys you never added.}

\subsection{How do you remove duplicates?}% sec:c08-set
Put the values in a \term[\textbf{Set:} a collection in which each value appears at most once, kept in the order first added. Like a guest list where writing a name twice changes nothing.]{set}{Set}. A value can occur in a Set only once, compared with \code{===}, and the Set remembers the order values were first added. Spread it back into an array with \code{[...set]}.
\Needspace{6\baselineskip}
\begin{lstlisting}[language=JavaScript]
const tags = new Set(["js", "ts", "js"]);
console.log(tags.size, [...tags].join(","));
\end{lstlisting}
{\small\textit{Prints:}~\code{2 js,ts}}\par

\subsection{When is a Map better than an object?}% sec:c08-map
When you count or group by key. A \term[\textbf{Map:} a collection of key--value pairs where the key can be any value and pairs stay in insertion order. Like a coat check: any ticket can claim exactly one coat.]{map}{Map} stores key--value pairs: \code{set(key, value)}, \code{get(key)}, \code{has(key)}, \code{size}. Keys keep their type, so the number \code{1} and the string \code{"1"} are different keys; an object would turn both into \code{"1"}. A Map also iterates in insertion order, which settles ties ``first come, first served''. \code{counts.get(word) ?? 0} starts a new key at zero.
\Needspace{6\baselineskip}
\begin{lstlisting}[language=JavaScript]
const counts = new Map();
for (const w of ["a", "b", "a"]) {
  counts.set(w, (counts.get(w) ?? 0) + 1);
}
console.log(counts.get("a"), counts.size);
\end{lstlisting}
{\small\textit{Prints:}~\code{2 2}}\par

\subsection{Why can a plain object already contain "constructor"?}% sec:c08-inherited
Because every plain object inherits properties through its \term[\textbf{Prototype chain:} the series of objects an object borrows properties from when it lacks its own. Like a family recipe book you can read from even though it isn't yours.]{prototype}{prototype chain}. \code{\{\}["constructor"]} is a function, not \code{undefined}, so ``create the group if it doesn't exist yet'' skips it and the code then calls \code{push} on a function. A Map starts truly empty; \code{Object.fromEntries(map)} turns it into a plain object at the end.
\Needspace{6\baselineskip}
\begin{lstlisting}[language=JavaScript]
const groups = {};
console.log(typeof groups["constructor"]);
console.log(new Map().has("constructor"));
\end{lstlisting}
{\small\textit{Prints:}~\code{function}\enspace\textcolor{black!45}{|}\enspace \code{false}}\par

\mantra{Counting or grouping by key? Reach for a Map.}
\begin{ponder}\code{mostFrequentWord} keeps the first word on a tie because it uses \code{>}. What changes with \code{>=}? What would change if the counts were kept in a plain object whose keys were numbers?\end{ponder}

\partline{E. How do you solve nested problems and survive failure?}{Recursion, errors and JSON.}
\phantomsection\label{ch:c09}
\chap{9. How can a function solve a problem by calling itself?}{A recursive function answers the smallest case directly and hands a smaller version of the problem to itself.}

\subsection{What makes recursion stop?}% sec:c09-base
A base case. \term[\textbf{Recursion:} a function calling itself on a smaller part of the problem. Like nesting dolls: open one, find a smaller one, until the smallest.]{recursion}{Recursion} has two kinds of input: a \term[\textbf{Base case:} the input a recursive function answers without calling itself. The smallest doll, which doesn't open.]{basecase}{base case}, answered directly, and a recursive case, which calls the function again on something smaller. Without a base case that every input eventually reaches, the calls never stop and JavaScript throws a \code{RangeError} when its call stack is full.
\Needspace{6\baselineskip}
\begin{lstlisting}[language=JavaScript]
function forever(n) {
  return forever(n + 1);
}
try {
  forever(0);
} catch (error) {
  console.log(error.name);
}
\end{lstlisting}
{\small\textit{Prints:}~\code{RangeError}}\par

\subsection{Why does recursion fit nested data?}% sec:c09-nested
Because each nested piece is the same problem, smaller: an array inside an array, a folder inside a folder. The function handles one level and calls itself for each nested piece, so it works at any depth. \code{flat()} removes only one level of nesting.
\Needspace{6\baselineskip}
\begin{lstlisting}[language=JavaScript]
function depth(item) {
  if (!Array.isArray(item)) return 0;
  return 1 + Math.max(0, ...item.map(depth));
}
console.log(depth([1, [2, [3]]]), depth(7));
\end{lstlisting}
{\small\textit{Prints:}~\code{3 0}}\par

\subsection{In what order do recursive calls finish?}% sec:c09-order
Last started, first finished. Each call waits on the \term[\textbf{Call stack:} the list of function calls that have started but not yet returned. Like a stack of plates: each call adds one, each return removes the top one.]{callstack}{call stack} until the smaller call returns, then uses that result. So results are put together on the way back out.
\Needspace{6\baselineskip}
\begin{lstlisting}[language=JavaScript]
function nest(n) {
  if (n === 0) return "core";
  const inner = nest(n - 1);
  console.log("back in", n);
  return `(${inner})`;
}
console.log(nest(2));
\end{lstlisting}
{\small\textit{Prints:}~\code{back in 1}\enspace\textcolor{black!45}{|}\enspace \code{back in 2}\enspace\textcolor{black!45}{|}\enspace \code{((core))}}\par

\mantra{Write the smallest case first; then shrink the problem toward it.}
\begin{ponder}\code{factorial} returns 1 for any \code{n <= 1}. What does \code{factorial(-3)} return? Is that a correct answer, or would throwing an error be better?\end{ponder}

\phantomsection\label{ch:c10}
\chap{10. What should a program do when something goes wrong?}{\code{throw} stops the function and travels up the calls until a \code{catch} handles it; catch only where you can do something useful, and \code{finally} runs whichever way the block exits.}

\subsection{What happens when you throw?}% sec:c10-throw
The function stops and an \term[\textbf{Exception:} an error object thrown to signal that something went wrong. Like a fire alarm: work stops and everyone heads for the nearest exit.]{exception}{exception} travels up the call stack to the first \code{catch} that waits for it; with no \code{catch}, the program ends. Throw the most specific built-in error: \code{RangeError} for a value outside the allowed range. An error has a \code{name} and a \code{message}.
\Needspace{6\baselineskip}
\begin{lstlisting}[language=JavaScript]
function checkAge(age) {
  if (age < 0) throw new RangeError("negative age");
  return age;
}
try {
  checkAge(-1);
} catch (error) {
  console.log(`${error.name}: ${error.message}`);
}
\end{lstlisting}
{\small\textit{Prints:}~\code{RangeError: negative age}}\par

\subsection{Where should you catch an error?}% sec:c10-catch
Only where you can do something useful. \code{try \{ \ldots \} catch (error) \{ \ldots \}} runs the \code{catch} block if the \code{try} block throws. Catching and returning a default \emph{swallows} the error: the caller can no longer tell ``broken input'' from ``no data''. Handle the cases you expect (such as blank text) and let the rest propagate. \code{JSON.parse} throws a \code{SyntaxError} on text that isn't valid JSON.
\Needspace{6\baselineskip}
\begin{lstlisting}[language=JavaScript]
try {
  JSON.parse("{oops");
} catch (error) {
  console.log(error.name);
}
\end{lstlisting}
{\small\textit{Prints:}~\code{SyntaxError}}\par

\subsection{What does finally guarantee?}% sec:c10-finally
A \code{finally} block always runs before control leaves the \code{try} construct---after a normal finish, an error, or a \code{return}.
\Needspace{6\baselineskip}
\begin{lstlisting}[language=JavaScript]
function f() {
  try {
    return "value";
  } finally {
    console.log("cleanup");
  }
}
console.log(f());
\end{lstlisting}
{\small\textit{Prints:}~\code{cleanup}\enspace\textcolor{black!45}{|}\enspace \code{value}}\par

\subsection{How do you make your own kind of error?}% sec:c10-custom
Extend \code{Error} (the class syntax is in \hyperref[ch:c12]{Ch.~12}) and set \code{name}. The subclass inherits \code{message}, but its name stays \code{"Error"} unless you set it. Callers can then tell your errors apart.
\Needspace{6\baselineskip}
\begin{lstlisting}[language=JavaScript]
class NotFound extends Error {
  name = "NotFound";
}
console.log(String(new NotFound("no user")));
\end{lstlisting}
{\small\textit{Prints:}~\code{NotFound: no user}}\par

\mantra{Throw early, catch where you can act, never catch and forget.}
\begin{ponder}\code{readConfig} returns \code{\{\}} for blank text but throws on bad JSON. Why is blank text ``expected'' and bad JSON not? Who should decide---the function or its caller?\end{ponder}

\phantomsection\label{ch:c11}
\chap{11. How do you turn data into text and back?}{\code{JSON.stringify} writes values as text and quietly drops what JSON cannot hold; \code{JSON.parse} reads the text back and throws on bad input.}

\subsection{How do you read JSON?}% sec:c11-parse
\term[\textbf{JSON:} a text format for data (objects, arrays, strings, numbers, booleans, null) that any program can read. Like a shipping label written in a standard layout.]{json}{JSON} is text. \code{JSON.parse(text)} turns it into JavaScript values. Missing fields come back as \code{undefined}, so supply defaults with \code{??} (\hyperref[ch:c02]{Ch.~2}), which keeps a saved \code{0} or \code{false}.
\Needspace{6\baselineskip}
\begin{lstlisting}[language=JavaScript]
const data = JSON.parse('{"volume": 0}');
console.log(data.volume ?? 50, data.muted ?? false);
\end{lstlisting}
{\small\textit{Prints:}~\code{0 false}}\par

\subsection{What does JSON.stringify leave out?}% sec:c11-stringify
Everything JSON cannot represent. \code{undefined} and functions are left out of objects and become \code{null} in arrays. A \code{Date} is written as a string, and a \code{Map} or \code{Set} becomes \code{"\{\}"} because only an object's own enumerable properties are written. The third argument sets the indentation: \code{JSON.stringify(value, null, 2)}.
\Needspace{6\baselineskip}
\begin{lstlisting}[language=JavaScript]
const value = { a: 1, b: undefined, c: [undefined], m: new Map() };
console.log(JSON.stringify(value));
console.log(JSON.stringify([1], null, 2).split("\n").length);
\end{lstlisting}
{\small\textit{Prints:}~\code{\{\textquotedbl{}a\textquotedbl{}:1,\textquotedbl{}c\textquotedbl{}:[null],\textquotedbl{}m\textquotedbl{}:\{\}\}}\enspace\textcolor{black!45}{|}\enspace \code{3}}\par

\mantra{stringify drops what JSON can't spell; parse throws on what it can't read.}
\begin{ponder}You save \code{\{ when: new Date() \}} with \code{stringify} and read it back with \code{parse}. What type is \code{when} now? How would you get a \code{Date} back?\end{ponder}

\partline{F. How do bigger programs fit together?}{Classes, asynchronous code and modules.}
\phantomsection\label{ch:c12}
\chap{12. How do classes bundle data with behaviour, and what is this?}{A class is a template for objects that share methods; \code{\#fields} are private; \code{this} is set by how a method is called, so a detached method or a \code{function} callback loses it.}

\subsection{What does a class give you?}% sec:c12-class
A \term[\textbf{Class:} a template for making objects that share the same methods. Like a blueprint.]{class}{class} describes objects; \code{new} makes an \term[\textbf{Instance:} one object made from a class. A house built from the blueprint.]{instance}{instance}. In the \code{constructor}, \code{this} is the new instance, so the constructor stores its starting data on \code{this}. Methods are written once in the class and shared by every instance.
\Needspace{6\baselineskip}
\begin{lstlisting}[language=JavaScript]
class Point {
  constructor(x, y) { this.x = x; this.y = y; }
  sum() { return this.x + this.y; }
}
console.log(new Point(2, 3).sum());
\end{lstlisting}
{\small\textit{Prints:}~\code{5}}\par

\subsection{How do you keep state private?}% sec:c12-private
Prefix the field name with \code{\#}. A \term[\textbf{Private field:} a class field whose name starts with \code{\#}; only code inside the class can read or write it. Like a safe that opens only from the inside.]{privatefield}{private field} must be declared in the class body and can only be used inside the class. From outside, \code{obj.balance} is just a different, missing property (\code{undefined}).
\Needspace{6\baselineskip}
\begin{lstlisting}[language=JavaScript]
class Safe {
  #code = 42;
  check(n) { return n === this.#code; }
}
const safe = new Safe();
console.log(safe.check(42), safe.code);
\end{lstlisting}
{\small\textit{Prints:}~\code{true undefined}}\par

\subsection{Why does this go missing?}% sec:c12-this
Because \code{this} is decided at the moment of the call: \code{greeter.greet()} sets it to \code{greeter}, but calling a method pulled off its object sets it to \code{undefined} (class and module code is strict), and reading a property of \code{undefined} throws. \code{fn.call(obj)} passes \code{this} for one call; \code{fn.bind(obj)} makes a copy with \code{this} fixed. An arrow function has no \code{this} of its own and uses the one around it, which is why arrow callbacks inside methods work and \code{function} callbacks do not.
\Needspace{6\baselineskip}
\begin{lstlisting}[language=JavaScript]
const counter = { n: 1, get() { return this.n; } };
const loose = counter.get;
try {
  loose();
} catch (error) {
  console.log(error.name);
}
console.log(loose.call(counter), loose.bind(counter)());
\end{lstlisting}
{\small\textit{Prints:}~\code{TypeError}\enspace\textcolor{black!45}{|}\enspace \code{1 1}}\par

\subsection{How does one class build on another?}% sec:c12-extends
With \code{extends}: the child class gets the parent's methods and can replace (override) some. A child constructor must call \code{super(...)}---the parent's constructor---before it uses \code{this}. A method inherited from the parent that calls \code{this.area()} gets the child's version.
\Needspace{6\baselineskip}
\begin{lstlisting}[language=JavaScript]
class Animal {
  constructor(name) { this.name = name; }
  describe() { return `${this.name}: ${this.sound()}`; }
  sound() { return "..."; }
}
class Dog extends Animal {
  sound() { return "woof"; }
}
console.log(new Dog("Rex").describe());
\end{lstlisting}
{\small\textit{Prints:}~\code{Rex: woof}}\par

\mantra{this is decided at the call, not where the method was written.}
\begin{ponder}\code{setTimeout(account.deposit, 0, 10)} fails, but \code{setTimeout(() => account.deposit(10), 0)} works. What is \code{this} inside \code{deposit} in each case?\end{ponder}

\phantomsection\label{ch:c13}
\chap{13. How does JavaScript wait without stopping?}{A promise stands for a result that isn't ready yet; \code{await} pauses only its own async function; promise callbacks run before timers; \code{Promise.all} waits for every promise and fails on the first rejection.}

\subsection{What is a promise, and what does await do?}% sec:c13-promise
A \term[\textbf{Promise:} an object standing for a result that will arrive later, or a failure. Like a claim ticket for a coat that isn't ready yet.]{promise}{promise} is either still pending, fulfilled with a value, or rejected with an error. Inside an \term[\textbf{Async function:} a function declared with \code{async}; it always returns a promise and may use \code{await}. Like a cook who can step away from one dish while it simmers.]{async}{async function}, \code{await p} pauses that function until \code{p} settles, then gives its value (or throws its error). The rest of the program keeps running meanwhile.
\Needspace{6\baselineskip}
\begin{lstlisting}[language=JavaScript]
async function slowDouble(n) {
  await null;
  return n * 2;
}
slowDouble(4).then((v) => console.log("got", v));
console.log("first");
\end{lstlisting}
{\small\textit{Prints:}~\code{first}\enspace\textcolor{black!45}{|}\enspace \code{got 8}}\par

\subsection{In what order do callbacks run?}% sec:c13-order
Synchronous code first. Promise callbacks are queued as \term[\textbf{Microtask:} a small job (such as a promise callback) run as soon as the current code finishes, before timers. Like the express lane at a checkout.]{microtask}{microtasks} and run as soon as the current code finishes; timer callbacks (\code{setTimeout}) wait in a separate queue and run after that, even with a 0 ms delay.
\Needspace{6\baselineskip}
\begin{lstlisting}[language=JavaScript]
setTimeout(() => console.log("timer"), 0);
Promise.resolve().then(() => console.log("promise"));
console.log("sync");
\end{lstlisting}
{\small\textit{Prints:}~\code{sync}\enspace\textcolor{black!45}{|}\enspace \code{promise}\enspace\textcolor{black!45}{|}\enspace \code{timer}}\par

\subsection{How do you run several requests at once?}% sec:c13-all
Start them all, then wait for all. A promise's work starts when the promise is created; \code{Promise.all} only waits. It fulfills with an array of the values in input order, or rejects with the first rejection. \code{await} inside a loop waits for each request before starting the next, so the requests run one at a time. \code{Promise.allSettled} always fulfills, with a status for each promise.
\Needspace{6\baselineskip}
\begin{lstlisting}[language=JavaScript]
const wait = (ms, v) => new Promise((r) => setTimeout(() => r(v), ms));
const start = Date.now();
const both = await Promise.all([wait(100, "a"), wait(100, "b")]);
console.log(both.join(""), Date.now() - start < 190);
\end{lstlisting}
{\small\textit{Prints:}~\code{ab true}}\par

\subsection{Why must you check response.ok?}% sec:c13-fetch
Because \code{fetch} rejects only when the request itself fails (a bad URL, no network). A 404 or 500 answer still \emph{fulfills} with a response. \code{response.ok} is true only for status codes 200--299, so check it and throw yourself.
\Needspace{6\baselineskip}
\begin{lstlisting}[language=JavaScript]
const response = await fetch(url);
if (!response.ok) throw new Error(`HTTP ${response.status}`);
const data = await response.json();
\end{lstlisting}

\mantra{Start everything, then await all of it.}
\begin{ponder}With \code{Promise.all}, \code{loadNames} returns names in the order of \code{ids}, even if user~3's answer arrives first. What guarantees that order?\end{ponder}

\phantomsection\label{ch:c14}
\chap{14. How do separate files share code?}{Each file is a module with its own scope: \code{export} shares a name, \code{import} borrows it---named imports in braces, the default import without---and a module runs once, so its importers share its state.}

\subsection{How do you share a function between files?}% sec:c14-import
Export it from one \term[\textbf{Module:} a file with its own scope that shares only what it exports. Like a room with labelled hatches: only what is passed through a hatch leaves the room.]{module}{module} and import it in another. A \term[\textbf{Named export:} a value exported under its own name and imported in braces, \code{import \{ PI \}}. A hatch with a nameplate.]{namedexport}{named export} is imported in braces under its exact name (rename with \code{as}). A module can have one \term[\textbf{Default export:} the one unnamed export of a module, imported without braces under any name you choose. The main hatch.]{defaultexport}{default export}, imported without braces under any name.
\Needspace{6\baselineskip}
\begin{lstlisting}[language=JavaScript]
// shapes.js
export const sides = { triangle: 3 };
export default function area(w, h) {
  return w * h;
}
// main.js
import area, { sides } from "./shapes.js";
console.log(area(2, 3), sides.triangle);
\end{lstlisting}
{\small\textit{Prints:}~\code{6 3}}\par

\subsection{What happens if you import a name that isn't exported?}% sec:c14-missing
Nothing runs. Imports are checked before any code in the importing file runs, so asking for a named export that doesn't exist---for example asking for \code{\{ add \}} when \code{add} is the default export---fails at once with a \code{SyntaxError}.

\subsection{If two files import the same module, do they get two copies?}% sec:c14-once
No. A module is evaluated once, and every import is a live, read-only view of its exports. Two imports of the same function share the same variables behind it.
\Needspace{6\baselineskip}
\begin{lstlisting}[language=JavaScript]
// log.js
export const lines = [];
export function log(x) {
  lines.push(x);
}
// main.js
import { log, lines } from "./log.js";
import { log as record } from "./log.js";
log("a");
record("b");
console.log(lines.join(","));
\end{lstlisting}
{\small\textit{Prints:}~\code{a,b}}\par

\mantra{Export on purpose; import by name.}
\begin{ponder}A module keeps a counter in a \code{let} and exports it. Can a file that imports it write \code{count = 0} to reset it? What does ``read-only live binding'' predict?\end{ponder}

\phantomsection\label{synthesis}
\section*{How the Parts Answer the Question}\addcontentsline{toc}{section}{How the Parts Answer the Question}

Writing JavaScript that does what you mean takes three kinds of knowledge, and the Parts supply them in order.

\textbf{What the values really are (A).} Most surprises in this Part come from a value being something other than you assumed: \code{typeof null} is \code{"object"}, \code{0 == ""} is true, \code{config.port || 3000} throws away a real \code{0}, and a \code{const} object can still change. Knowing the seven primitive types, references, truthiness and the three logical operators removes those surprises at the source.

\textbf{How control moves (B, C, E).} Loops decide how often code runs and when it stops (\code{<=} versus \code{<}, \code{return} at the first match). Functions decide what code can see: closures keep private state, and recursion turns nested data into the same small problem. Errors decide where control goes when something fails: throw early, catch only where you can act, and let \code{finally} clean up.

\textbf{How data and programs are shaped (D, F).} Arrays, objects, Maps and Sets each fit a different job, and each has one trap worth knowing: \code{sort} compares text, spread copies one level, plain objects inherit keys. Classes bundle data with behaviour but tie \code{this} to the call; promises let one task wait without stopping the rest; modules share code by name and run once.

\textbf{The strongest objection} is that you could learn all of this by running code and fixing the errors. The answer: the traps in this Part mostly raise \emph{no} error. \code{forEach} quietly returns \code{undefined}, \code{||} quietly replaces \code{0}, \code{fetch} quietly resolves on a 404, and \code{sort} quietly misorders numbers. Code that runs and gives a wrong answer is only caught by knowing what it should do---which is what the cards drill.

\section*{Test Yourself}\addcontentsline{toc}{section}{Test Yourself}
Answer without looking; the link after each question points to the chapter that holds the answer.
\begin{enumerate}[leftmargin=1.6em,itemsep=3pt]
\item What does \code{typeof null} return, and how should you test for \code{null}? \hfill{\small\hyperref[ch:c01]{$\rightarrow$ Ch.~ 1}}
\item Why does \code{config.port || 3000} break for port 0, and what fixes it? \hfill{\small\hyperref[ch:c02]{$\rightarrow$ Ch.~ 2}}
\item A loop must add the numbers 1 to n. Which comparison belongs in its test, and what goes wrong with the other one? \hfill{\small\hyperref[ch:c03]{$\rightarrow$ Ch.~ 3}}
\item Why is \code{\textquotedbl{}5\textquotedbl{} + 3} equal to \code{\textquotedbl{}53\textquotedbl{}} while \code{\textquotedbl{}5\textquotedbl{} - 3} is \code{2}? \hfill{\small\hyperref[ch:c04]{$\rightarrow$ Ch.~ 4}}
\item How do two counters made by the same function keep separate counts? \hfill{\small\hyperref[ch:c05]{$\rightarrow$ Ch.~ 5}}
\item Why does \code{[10, 9, 100].sort()} give 10, 100, 9, and how do you sort numbers? \hfill{\small\hyperref[ch:c06]{$\rightarrow$ Ch.~ 6}}
\item When is a destructuring default not used? \hfill{\small\hyperref[ch:c07]{$\rightarrow$ Ch.~ 7}}
\item Why can \code{\{\}[\textquotedbl{}constructor\textquotedbl{}]} break a grouping function, and what avoids it? \hfill{\small\hyperref[ch:c08]{$\rightarrow$ Ch.~ 8}}
\item Which two kinds of input must every recursive function handle? \hfill{\small\hyperref[ch:c09]{$\rightarrow$ Ch.~ 9}}
\item Why is catching an error and returning \code{\{\}} usually wrong? \hfill{\small\hyperref[ch:c10]{$\rightarrow$ Ch.~ 10}}
\item What happens to \code{undefined} when \code{JSON.stringify} meets it in an object? In an array? \hfill{\small\hyperref[ch:c11]{$\rightarrow$ Ch.~ 11}}
\item Why does a method pulled off its object lose \code{this}, and what restores it? \hfill{\small\hyperref[ch:c12]{$\rightarrow$ Ch.~ 12}}
\item Does \code{Promise.all} start the work it waits for? What does it return? \hfill{\small\hyperref[ch:c13]{$\rightarrow$ Ch.~ 13}}
\item What happens when a file imports a named export that does not exist? \hfill{\small\hyperref[ch:c14]{$\rightarrow$ Ch.~ 14}}
\end{enumerate}
\newpage\phantomsection\partline{Card-by-Card Solution Guide}{One entry for every card in this Part's decks.}
\begin{corebox}\textbf{How this part works.} Each entry is self-contained: it repeats the card, the terms it needs and the chapter sections it drills, so you can study a card without turning back. That repetition is deliberate (the one exception to saying each thing once). For each card: the card as it appears in your deck; how to solve it (approach, what must stay true, what lives in memory, the types involved, how it fails, how to test it) and the trap that most often costs it; the solution, checked by actually running it; and the readings from the chapters.\end{corebox}
\subsection*{Which card drills which chapters?}
{\footnotesize\begin{tabularx}{\linewidth}{@{}rXl@{}}\toprule
\textbf{\#} & \textbf{Card} & \textbf{Chapters it drills}\\\midrule
1 & What does this print? (js01-typeof) & Ch.~\hyperref[ch:c01]{1}\\
2 & \code{describeType(value)} & Ch.~\hyperref[ch:c01]{1}\\
3 & \code{parseCount(text)} & Ch.~\hyperref[ch:c01]{1}\\
4 & What does this print? (js01-const) & Ch.~\hyperref[ch:c01]{1}\\
5 & What does this print? (js01-pass-by-sharing) & Ch.~\hyperref[ch:c01]{1}\\
6 & What does this print? (js01-tdz) & Ch.~\hyperref[ch:c01]{1}\\
7 & What does this print? (js02-equality) & Ch.~\hyperref[ch:c02]{2}\\
8 & \code{Fix getPort(config)} & Ch.~\hyperref[ch:c02]{2}\\
9 & \code{letterGrade(score)} & Ch.~\hyperref[ch:c02]{2}\\
10 & What does this print? (js02-logical-operands) & Ch.~\hyperref[ch:c02]{2}\\
11 & \code{countTruthy(values)} & Ch.~\hyperref[ch:c02]{2}\\
12 & \code{sumTo(n)} & Ch.~\hyperref[ch:c03]{3}\\
13 & \code{countVowels(text)} & Ch.~\hyperref[ch:c03]{3}\\
14 & \code{fizzBuzz(n)} & Ch.~\hyperref[ch:c03]{3}\\
15 & \code{firstIndexOver(nums, limit)} & Ch.~\hyperref[ch:c03]{3}\\
16 & What does this print? (js03-closure-in-loop) & Ch.~\hyperref[ch:c03]{3}\\
17 & What does this print? (js03-for-in-vs-of) & Ch.~\hyperref[ch:c03]{3}\\
18 & \code{capitalize(word)} & Ch.~\hyperref[ch:c04]{4}\\
19 & \code{isPalindrome(text)} & Ch.~\hyperref[ch:c04]{4}\\
20 & \code{initials(fullName)} & Ch.~\hyperref[ch:c04]{4}\\
21 & What does this print? (js04-plus-coercion) & Ch.~\hyperref[ch:c04]{4}\\
22 & What does this print? (js04-immutable) & Ch.~\hyperref[ch:c04]{4}\\
23 & What does this print? (js05-hoisting) & Ch.~\hyperref[ch:c05]{5}\\
24 & \code{makeCounter(start)} & Ch.~\hyperref[ch:c05]{5}\\
25 & \code{once(fn)} & Ch.~\hyperref[ch:c05]{5}\\
26 & \code{sum(...numbers)} & Ch.~\hyperref[ch:c05]{5}, \hyperref[ch:c06]{6}\\
27 & What does this print? (js06-sort-default) & Ch.~\hyperref[ch:c06]{6}\\
28 & \code{squaresOfEvens(nums)} & Ch.~\hyperref[ch:c06]{6}\\
29 & \code{average(nums)} & Ch.~\hyperref[ch:c06]{6}\\
30 & \code{Fix doubled(nums)} & Ch.~\hyperref[ch:c06]{6}\\
\bottomrule\end{tabularx}}\par
{\footnotesize\begin{tabularx}{\linewidth}{@{}rXl@{}}\toprule
\textbf{\#} & \textbf{Card} & \textbf{Chapters it drills}\\\midrule
31 & \code{groupByCategory(items)} & Ch.~\hyperref[ch:c06]{6}\\
32 & What does this print? (js06-shared-reference) & Ch.~\hyperref[ch:c06]{6}\\
33 & \code{formatUser(user)} & Ch.~\hyperref[ch:c07]{7}\\
34 & What does this print? (js07-shallow-spread) & Ch.~\hyperref[ch:c07]{7}\\
35 & \code{withDefaults(options)} & Ch.~\hyperref[ch:c07]{7}\\
36 & \code{pick(obj, keys)} & Ch.~\hyperref[ch:c07]{7}\\
37 & What does this print? (js07-optional-chaining) & Ch.~\hyperref[ch:c07]{7}\\
38 & \code{unique(values)} & Ch.~\hyperref[ch:c08]{8}\\
39 & \code{mostFrequentWord(text)} & Ch.~\hyperref[ch:c08]{8}\\
40 & What does this print? (js08-map-keys) & Ch.~\hyperref[ch:c08]{8}\\
41 & \code{Fix groupByCategory for any category} & Ch.~\hyperref[ch:c08]{8}\\
42 & \code{factorial(n)} & Ch.~\hyperref[ch:c09]{9}\\
43 & \code{sumNested(items)} & Ch.~\hyperref[ch:c09]{9}\\
44 & \code{countFiles(folder)} & Ch.~\hyperref[ch:c09]{9}\\
45 & What does this print? (js09-trace) & Ch.~\hyperref[ch:c09]{9}\\
46 & \code{parseAge(text)} & Ch.~\hyperref[ch:c10]{10}\\
47 & What does this print? (js10-finally) & Ch.~\hyperref[ch:c10]{10}\\
48 & \code{Fix readConfig(json)} & Ch.~\hyperref[ch:c10]{10}, \hyperref[ch:c11]{11}\\
49 & \code{ValidationError} & Ch.~\hyperref[ch:c10]{10}, \hyperref[ch:c12]{12}\\
50 & What does this print? (js11-stringify-drops) & Ch.~\hyperref[ch:c11]{11}\\
51 & \code{loadSettings(json)} & Ch.~\hyperref[ch:c11]{11}\\
52 & \code{prettyJson(text)} & Ch.~\hyperref[ch:c11]{11}\\
53 & What does this print? (js11-map-to-json) & Ch.~\hyperref[ch:c11]{11}, \hyperref[ch:c08]{8}\\
54 & \code{class BankAccount} & Ch.~\hyperref[ch:c12]{12}\\
55 & What does this print? (js12-lost-this) & Ch.~\hyperref[ch:c12]{12}\\
56 & \code{Fix Cart.total()} & Ch.~\hyperref[ch:c12]{12}\\
57 & \code{Circle and Rectangle extend Shape} & Ch.~\hyperref[ch:c12]{12}\\
58 & What does this print? (js13-event-loop) & Ch.~\hyperref[ch:c13]{13}\\
59 & \code{Make loadNames(ids) concurrent} & Ch.~\hyperref[ch:c13]{13}\\
60 & \code{Fix getJson(url)} & Ch.~\hyperref[ch:c13]{13}\\
\bottomrule\end{tabularx}}\par
{\footnotesize\begin{tabularx}{\linewidth}{@{}rXl@{}}\toprule
\textbf{\#} & \textbf{Card} & \textbf{Chapters it drills}\\\midrule
61 & What does this print? (js13-all-vs-allsettled) & Ch.~\hyperref[ch:c13]{13}\\
62 & What does this print? (js14-default-and-named) & Ch.~\hyperref[ch:c14]{14}\\
63 & What does this print? (js14-missing-named-export) & Ch.~\hyperref[ch:c14]{14}\\
64 & What does this print? (js14-shared-state) & Ch.~\hyperref[ch:c14]{14}\\
\bottomrule\end{tabularx}}\par
\deckhead{01 · Values, Variables \& Types}{6}
\begin{corebox}\textbf{Core idea.} Every piece of data is a value of some type; a variable is a name that points at a value, and \code{const} pins the name, not the value.\end{corebox}
\begin{sumbox}{Deck summary sheet --- 01 · Values, Variables \& Types}
JavaScript has seven primitive types (string, number, bigint, boolean, undefined, symbol, null) and objects. \code{typeof} names them, with two traps: \code{typeof null} is "object", and arrays are objects too. \code{const} stops a variable being reassigned, not the object it points to from changing.\par
\smallskip{\footnotesize Read more: \href{https://developer.mozilla.org/en-US/docs/Web/JavaScript/Guide/Data_structures}{MDN — JavaScript data types and data structures}; \href{https://developer.mozilla.org/en-US/docs/Web/JavaScript/Reference/Statements/let}{MDN — let}}\end{sumbox}
\cardhead{Card 1 --- What does this print?}{predict · js01-typeof}
\safemargin{\textbf{Terms for this card}\par\smallskip \textbf{Value:} one piece of data the program works with, such as \code{42} or \code{"hi"}. Like a coin: you can hand it around and look at it.\par\smallskip \textbf{Primitive:} a simple value that cannot be changed: a string, number, bigint, boolean, undefined, symbol or null. Like a number stamped on a coin.\par\smallskip \textbf{Object:} a value that holds other values under names and can be changed. Like a filing cabinet whose drawers you can open, fill and empty.\par\smallskip \textbf{typeof:} an operator that returns a string naming a value's type. Like a label maker that reads a value and prints what kind it is.}
\noindent{\small\textbf{Covered in:} \hyperref[ch:c01]{Ch.~1} of this guide; documentation: \href{https://developer.mozilla.org/en-US/docs/Web/JavaScript/Guide/Data_structures}{MDN — JavaScript data types}.}\par
\cardsub{The card}
What does this print?\par
\begin{lstlisting}[language=JavaScript]
console.log(typeof 42);
console.log(typeof "42");
console.log(typeof null);
console.log(typeof []);
console.log(typeof undefined);
console.log(typeof 10n);
\end{lstlisting}
\cardsub{How to solve it}
\begin{cardrows}
\item[Solve] Evaluate each \code{typeof} on its own line, remembering the two exceptions.
\item[Invariants] \code{typeof} always returns one of a few fixed type names, as a string.
\item[Memory] Nothing is stored; each value is made and thrown away.
\item[Types] number, string, null, array (an object), undefined, bigint.
\item[Failure] Answering "null" for \code{typeof null} or "array" for \code{typeof []}.
\item[Tests] Run it, then change one line and predict again.
\end{cardrows}
\trap \code{typeof null} and \code{typeof []} are both "object".\par
\cardsub{Solution}
{\small\textbf{Prints:}}
\begin{lstlisting}
number
string
object
object
undefined
bigint
\end{lstlisting}
{\small\textit{Verified: Node v22.23.4 prints exactly this.}}\par
\smallskip{\small\textbf{Why:} typeof null is "object" (a bug kept for compatibility), so test for null with === null. Arrays are objects, so use Array.isArray to spot them. 10n is a BigInt.}\par
\cardsub{From the readings}
\readchap{1. What is a value, and how do you name one?}{Every piece of data is a value of some type; a variable is a name that points at a value, and \code{const} pins the name, not the value.}
\readq{What kinds of values are there?}
JavaScript has seven primitive types and objects. A \textbf{value} is either a \textbf{primitive}---string, number, bigint, boolean, \code{undefined}, symbol, \code{null}---or an \textbf{object}. Arrays and functions are objects too.

The \textbf{typeof} operator names a value's type. It has two traps: \code{typeof null} is \code{"object"}, and an array is reported as \code{"object"} as well. Test for \code{null} with \code{=== null} and for arrays with \code{Array.isArray}.
\Needspace{6\baselineskip}
\begin{lstlisting}[language=JavaScript]
console.log(typeof 3.5, typeof "hi", typeof true);
console.log(typeof {}, typeof [], Array.isArray([]));
console.log(typeof null, null === null);
\end{lstlisting}
{\small\textit{Prints:}~\code{number string boolean}\enspace\textcolor{black!45}{|}\enspace \code{object object true}\enspace\textcolor{black!45}{|}\enspace \code{object true}}\par

\cardend
\cardhead{Card 2 --- describeType(value)}{write · javascript · Easy}
\safemargin{\textbf{Terms for this card}\par\smallskip \textbf{Value:} one piece of data the program works with, such as \code{42} or \code{"hi"}. Like a coin: you can hand it around and look at it.\par\smallskip \textbf{Primitive:} a simple value that cannot be changed: a string, number, bigint, boolean, undefined, symbol or null. Like a number stamped on a coin.\par\smallskip \textbf{Object:} a value that holds other values under names and can be changed. Like a filing cabinet whose drawers you can open, fill and empty.\par\smallskip \textbf{typeof:} an operator that returns a string naming a value's type. Like a label maker that reads a value and prints what kind it is.}
\noindent{\small\textbf{Covered in:} \hyperref[ch:c01]{Ch.~1} of this guide; documentation: \href{https://developer.mozilla.org/en-US/docs/Web/JavaScript/Reference/Global_Objects/Array/isArray}{MDN — Array.isArray()}.}\par
\cardsub{The card}
\textbf{describeType(value)}\par
Write \code{describeType(value)}: return \code{\textquotedbl{}null\textquotedbl{}} for \code{null}, \code{\textquotedbl{}array\textquotedbl{}} for arrays, and what \code{typeof} says for everything else.\par
\begin{lstlisting}[language=JavaScript]
function describeType(value) {
  // TODO: handle null and arrays before falling back to typeof
  return typeof value;
}
\end{lstlisting}
\cardsub{How to solve it}
\begin{cardrows}
\item[Solve] Test \code{=== null} first, then \code{Array.isArray}, then fall back to \code{typeof}.
\item[Invariants] Each special case returns before the next check sees the value.
\item[Memory] No state: one value in, one string out.
\item[Types] Any value in; a string out.
\item[Failure] Using \code{!value} or \code{== null} also catches \code{0}, \code{\textquotedbl{}\textquotedbl{}} or \code{undefined}.
\item[Tests] \code{null}, \code{[1, 2]}, \code{\{\}}, \code{\textquotedbl{}\textquotedbl{}}, \code{0} and \code{undefined} must each keep their own answer.
\end{cardrows}
\trap Testing for null with \code{!value} mislabels 0 and the empty string.\par
\cardsub{Solution}
\begin{lstlisting}[language=JavaScript]
function describeType(value) {
  if (value === null) return "null";
  if (Array.isArray(value)) return "array";
  return typeof value;
}
\end{lstlisting}
{\small \textit{Verified: passes all 7 test cases on the card in the app's code runner; the starter code fails them.}}\par
\smallskip{\small\textbf{Why:} \code{=== null} is the only reliable null check, and \code{Array.isArray} the only reliable array check. \code{!value} would also catch 0, "" and undefined; \code{== null} would also catch undefined.}\par
\cardsub{From the readings}
\readchap{1. What is a value, and how do you name one?}{Every piece of data is a value of some type; a variable is a name that points at a value, and \code{const} pins the name, not the value.}
\readq{What kinds of values are there?}
JavaScript has seven primitive types and objects. A \textbf{value} is either a \textbf{primitive}---string, number, bigint, boolean, \code{undefined}, symbol, \code{null}---or an \textbf{object}. Arrays and functions are objects too.

The \textbf{typeof} operator names a value's type. It has two traps: \code{typeof null} is \code{"object"}, and an array is reported as \code{"object"} as well. Test for \code{null} with \code{=== null} and for arrays with \code{Array.isArray}.
\Needspace{6\baselineskip}
\begin{lstlisting}[language=JavaScript]
console.log(typeof 3.5, typeof "hi", typeof true);
console.log(typeof {}, typeof [], Array.isArray([]));
console.log(typeof null, null === null);
\end{lstlisting}
{\small\textit{Prints:}~\code{number string boolean}\enspace\textcolor{black!45}{|}\enspace \code{object object true}\enspace\textcolor{black!45}{|}\enspace \code{object true}}\par

\cardend
\cardhead{Card 3 --- parseCount(text)}{write · javascript · Easy}
\noindent{\small\textbf{Covered in:} \hyperref[ch:c01]{Ch.~1} of this guide; documentation: \href{https://developer.mozilla.org/en-US/docs/Web/JavaScript/Reference/Global_Objects/Number/isInteger}{MDN — Number.isInteger()}.}\par
\cardsub{The card}
\textbf{parseCount(text)}\par
Write \code{parseCount(text)}: return the whole number written in \code{text}, or \code{null} if it isn't one. \code{\textquotedbl{}007\textquotedbl{}} is 7, \code{\textquotedbl{}-4\textquotedbl{}} is -4; \code{\textquotedbl{}\textquotedbl{}}, \code{\textquotedbl{}   \textquotedbl{}}, \code{\textquotedbl{}3.5\textquotedbl{}} and \code{\textquotedbl{}12abc\textquotedbl{}} are all \code{null}.\par
\begin{lstlisting}[language=JavaScript]
function parseCount(text) {
  return parseInt(text, 10);
}
\end{lstlisting}
\cardsub{How to solve it}
\begin{cardrows}
\item[Solve] Reject blank text, convert with \code{Number}, and accept the result only if \code{Number.isInteger} says so.
\item[Invariants] Only text that is entirely a whole number is accepted.
\item[Memory] One local number \code{n}.
\item[Types] string in; number or \code{null} out.
\item[Failure] \code{parseInt} accepts "12abc"; plain \code{Number} turns blank text into 0.
\item[Tests] "007", "-4", "", "   ", "3.5", "12abc".
\end{cardrows}
\trap \code{Number(\textquotedbl{}\textquotedbl{})} is 0, so blank input needs its own check.\par
\cardsub{Solution}
\begin{lstlisting}[language=JavaScript]
function parseCount(text) {
  if (text.trim() === "") return null;
  const n = Number(text);
  return Number.isInteger(n) ? n : null;
}
\end{lstlisting}
{\small \textit{Verified: passes all 7 test cases on the card in the app's code runner; the starter code fails them.}}\par
\smallskip{\small\textbf{Why:} \code{parseInt(\textquotedbl{}12abc\textquotedbl{})} stops at the first bad character and returns 12; \code{Number(\textquotedbl{}12abc\textquotedbl{})} is NaN. But \code{Number(\textquotedbl{}\textquotedbl{})} and \code{Number(\textquotedbl{}   \textquotedbl{})} are 0, so empty input needs its own check.}\par
\cardsub{From the readings}
\readchap{1. What is a value, and how do you name one?}{Every piece of data is a value of some type; a variable is a name that points at a value, and \code{const} pins the name, not the value.}
\readq{How do you turn text into a number safely?}
Use \code{Number(text)}, then check the result. \code{Number} reads the whole string: \code{Number("12abc")} is \code{NaN}, the special ``not a number'' value. \code{parseInt} stops at the first character it cannot read, so \code{parseInt("12abc", 10)} is \code{12}. Watch empty text: \code{Number("")} and \code{Number("   ")} are \code{0}. \code{Number.isInteger} then tells you whether you got a whole number.
\Needspace{6\baselineskip}
\begin{lstlisting}[language=JavaScript]
console.log(Number("42"), Number("12abc"), Number(""));
console.log(parseInt("12abc", 10), Number.isInteger(3.5));
\end{lstlisting}
{\small\textit{Prints:}~\code{42 NaN 0}\enspace\textcolor{black!45}{|}\enspace \code{12 false}}\par

\cardend
\cardhead{Card 4 --- What does this print?}{predict · js01-const}
\safemargin{\textbf{Terms for this card}\par\smallskip \textbf{Variable:} a name that points at a value. Like a name tag tied to a parcel: you can move the tag to another parcel (reassign) or open the parcel (change an object).\par\smallskip \textbf{Temporal dead zone (TDZ):} the part of a block before a \code{let}/\code{const} line runs, where reading the name throws a ReferenceError. Like a roped-off area before the shop opens.}
\noindent{\small\textbf{Covered in:} \hyperref[ch:c01]{Ch.~1} of this guide; documentation: \href{https://developer.mozilla.org/en-US/docs/Web/JavaScript/Reference/Statements/const}{MDN — const}.}\par
\cardsub{The card}
What does this print?\par
\begin{lstlisting}[language=JavaScript]
const settings = { theme: "dark" };
settings.theme = "light";
console.log(settings.theme);
try {
  settings = {};
} catch (error) {
  console.log(error.name);
}
\end{lstlisting}
\cardsub{How to solve it}
\begin{cardrows}
\item[Solve] Line 2 changes a property (allowed); the reassignment inside \code{try} throws, so \code{catch} prints the error's name.
\item[Invariants] A \code{const} name always points at the same object.
\item[Memory] One object, whose \code{theme} changes from dark to light.
\item[Types] object, string, TypeError.
\item[Failure] Believing \code{const} freezes the object, or that the reassignment is a SyntaxError.
\item[Tests] Run it, then change one line and predict again.
\end{cardrows}
\trap \code{const} blocks reassignment, not changes inside the object.\par
\cardsub{Solution}
{\small\textbf{Prints:}}
\begin{lstlisting}
light
TypeError
\end{lstlisting}
{\small\textit{Verified: Node v22.23.4 prints exactly this.}}\par
\smallskip{\small\textbf{Why:} const stops the variable from being reassigned (a TypeError at run time), but the object it points to can still change.}\par
\cardsub{From the readings}
\readchap{1. What is a value, and how do you name one?}{Every piece of data is a value of some type; a variable is a name that points at a value, and \code{const} pins the name, not the value.}
\readq{What do let, const and var actually fix?}
A \textbf{variable} is a name pointing at a value. \code{let} and \code{const} names exist only inside the block \code{\{ ... \}} that declares them; \code{const} also forbids pointing the name at a different value. It does not freeze the value: an object behind a \code{const} can still be changed.
\Needspace{6\baselineskip}
\begin{lstlisting}[language=JavaScript]
const team = ["Ada"];
team.push("Grace");
console.log(team.length);
\end{lstlisting}
{\small\textit{Prints:}~\code{2}}\par

\code{var} is older: its name exists, as \code{undefined}, from the top of its function. A \code{let} or \code{const} name cannot be read before its declaration line runs; that stretch is the \textbf{temporal dead zone}, and reading it there throws a \code{ReferenceError}.
\cardend
\cardhead{Card 5 --- What does this print?}{predict · js01-pass-by-sharing}
\safemargin{\textbf{Terms for this card}\par\smallskip \textbf{Reference:} the value that stands for an object: where it lives, not the object itself. Like a street address on a card, not a duplicate house.}
\noindent{\small\textbf{Covered in:} \hyperref[ch:c01]{Ch.~1} of this guide; documentation: \href{https://developer.mozilla.org/en-US/docs/Web/JavaScript/Guide/Functions}{MDN — Functions}.}\par
\cardsub{The card}
What does this print?\par
\begin{lstlisting}[language=JavaScript]
function rename(user) {
  user.name = "Grace";
  user = { name: "Linus" };
}
const person = { name: "Ada" };
rename(person);
console.log(person.name);
\end{lstlisting}
\cardsub{How to solve it}
\begin{cardrows}
\item[Solve] Follow the reference: the property change reaches the caller's object; the reassignment only moves the local name.
\item[Invariants] The caller's \code{person} keeps pointing at the same object throughout.
\item[Memory] Two objects: the original (renamed Grace) and a new one (Linus) that nothing outside sees.
\item[Types] objects holding strings.
\item[Failure] Answering Linus (the caller sees the reassignment) or Ada (nothing is shared).
\item[Tests] Run it, and log \code{user.name} inside the function after each line.
\end{cardrows}
\trap Reassigning a parameter never reaches the caller; changing a property through it does.\par
\cardsub{Solution}
{\small\textbf{Prints:}}
\begin{lstlisting}
Grace
\end{lstlisting}
{\small\textit{Verified: Node v22.23.4 prints exactly this.}}\par
\smallskip{\small\textbf{Why:} JavaScript passes every argument by value. For an object that value is a reference, so changing a property through it is seen by the caller, but pointing the parameter at a new object is not.}\par
\cardsub{From the readings}
\readchap{1. What is a value, and how do you name one?}{Every piece of data is a value of some type; a variable is a name that points at a value, and \code{const} pins the name, not the value.}
\readq{What happens when you pass a value to a function?}
The function gets a copy of the value. For an object, that value is a \textbf{reference}: the function and the caller now hold the same address. Changing a property through it is seen by the caller; pointing the parameter at a new object is not.
\Needspace{6\baselineskip}
\begin{lstlisting}[language=JavaScript]
function addTag(list) {
  list.push("new");
  list = ["other"];
}
const tags = ["old"];
addTag(tags);
console.log(tags.join(","));
\end{lstlisting}
{\small\textit{Prints:}~\code{old,new}}\par

\cardend
\cardhead{Card 6 --- What does this print?}{predict · js01-tdz}
\safemargin{\textbf{Terms for this card}\par\smallskip \textbf{Variable:} a name that points at a value. Like a name tag tied to a parcel: you can move the tag to another parcel (reassign) or open the parcel (change an object).\par\smallskip \textbf{Temporal dead zone (TDZ):} the part of a block before a \code{let}/\code{const} line runs, where reading the name throws a ReferenceError. Like a roped-off area before the shop opens.}
\noindent{\small\textbf{Covered in:} \hyperref[ch:c01]{Ch.~1} of this guide; documentation: \href{https://developer.mozilla.org/en-US/docs/Web/JavaScript/Reference/Statements/let}{MDN — let}.}\par
\cardsub{The card}
What does this print?\par
\begin{lstlisting}[language=JavaScript]
console.log(typeof hoisted);
try {
  console.log(early);
} catch (error) {
  console.log(error.name);
}
let early = 1;
var hoisted = 2;
\end{lstlisting}
\cardsub{How to solve it}
\begin{cardrows}
\item[Solve] \code{var hoisted} exists as \code{undefined} from the top; \code{early} is in its temporal dead zone until its line runs.
\item[Invariants] A \code{let} name cannot be read before its declaration runs.
\item[Memory] Both names exist from the start; only the \code{var} has a value (\code{undefined}).
\item[Types] string ("undefined"), ReferenceError.
\item[Failure] Expecting \code{typeof hoisted} to be "number", or \code{early} to print undefined.
\item[Tests] Run it, then move \code{let early = 1;} to the top and run again.
\end{cardrows}
\trap Before its line, a \code{var} reads as undefined but a \code{let} throws.\par
\cardsub{Solution}
{\small\textbf{Prints:}}
\begin{lstlisting}
undefined
ReferenceError
\end{lstlisting}
{\small\textit{Verified: Node v22.23.4 prints exactly this.}}\par
\smallskip{\small\textbf{Why:} A var exists (as undefined) from the top of its function. A let or const exists too, but reading it before its declaration line throws: it is in the temporal dead zone.}\par
\cardsub{From the readings}
\readchap{1. What is a value, and how do you name one?}{Every piece of data is a value of some type; a variable is a name that points at a value, and \code{const} pins the name, not the value.}
\readq{What do let, const and var actually fix?}
A \textbf{variable} is a name pointing at a value. \code{let} and \code{const} names exist only inside the block \code{\{ ... \}} that declares them; \code{const} also forbids pointing the name at a different value. It does not freeze the value: an object behind a \code{const} can still be changed.
\Needspace{6\baselineskip}
\begin{lstlisting}[language=JavaScript]
const team = ["Ada"];
team.push("Grace");
console.log(team.length);
\end{lstlisting}
{\small\textit{Prints:}~\code{2}}\par

\code{var} is older: its name exists, as \code{undefined}, from the top of its function. A \code{let} or \code{const} name cannot be read before its declaration line runs; that stretch is the \textbf{temporal dead zone}, and reading it there throws a \code{ReferenceError}.
\cardend
\deckhead{02 · Operators, Equality \& Conditionals}{5}
\begin{corebox}\textbf{Core idea.} \code{===} compares without converting; truthiness decides an \code{if}; and \code{||}, \code{\&\&} and \code{??} hand back one of their operands.\end{corebox}
\begin{sumbox}{Deck summary sheet --- 02 · Operators, Equality \& Conditionals}
\code{===} compares without converting types; \code{==} converts first, with surprising results. \code{||} picks the right side for any falsy value, \code{??} only for null and undefined. \code{\&\&}, \code{||} and \code{??} return one of their operands, not true/false.\par
\smallskip{\footnotesize Read more: \href{https://developer.mozilla.org/en-US/docs/Web/JavaScript/Equality_comparisons_and_sameness}{MDN — Equality comparisons and sameness}; \href{https://developer.mozilla.org/en-US/docs/Web/JavaScript/Reference/Operators/Nullish_coalescing}{MDN — Nullish coalescing operator (??)}}\end{sumbox}
\cardhead{Card 7 --- What does this print?}{predict · js02-equality}
\safemargin{\textbf{Terms for this card}\par\smallskip \textbf{Coercion:} JavaScript converting a value to another type on its own, e.g. \code{""} to \code{0}. Like a currency exchange that happens before you can compare prices.\par\smallskip \textbf{Strict equality (===):} comparison that treats values of different types as different. Like comparing two prices only if they are in the same currency.}
\noindent{\small\textbf{Covered in:} \hyperref[ch:c02]{Ch.~2} of this guide; documentation: \href{https://developer.mozilla.org/en-US/docs/Web/JavaScript/Reference/Operators/Strict_equality}{MDN — Strict equality (===)}.}\par
\cardsub{The card}
What does this print?\par
\begin{lstlisting}[language=JavaScript]
console.log(0 == "");
console.log(0 === "");
console.log(null == undefined);
console.log(null == 0);
console.log(NaN === NaN);
\end{lstlisting}
\cardsub{How to solve it}
\begin{cardrows}
\item[Solve] For each \code{==}, convert both sides first; for \code{===}, different types are never equal.
\item[Invariants] \code{===} never converts; \code{NaN} is never equal to anything.
\item[Memory] Nothing stored.
\item[Types] number, string ("" becomes 0), null, undefined, NaN.
\item[Failure] Assuming \code{null == 0} (it is false) or \code{NaN === NaN} (also false).
\item[Tests] Run it, then change one line and predict again.
\end{cardrows}
\trap \code{null == undefined} is true, but \code{null == 0} is false.\par
\cardsub{Solution}
{\small\textbf{Prints:}}
\begin{lstlisting}
true
false
true
false
false
\end{lstlisting}
{\small\textit{Verified: Node v22.23.4 prints exactly this.}}\par
\smallskip{\small\textbf{Why:} == converts "" to 0 before comparing; === does not convert. null == undefined is true, but null is not == 0. NaN is not equal to anything, itself included: use Number.isNaN.}\par
\cardsub{From the readings}
\readchap{2. How does JavaScript compare values and choose a branch?}{\code{===} compares without converting; truthiness decides an \code{if}; and \code{||}, \code{\&\&} and \code{??} hand back one of their operands.}
\readq{Why prefer === over ==?}
Because \code{===} never converts. \code{==} first applies \textbf{coercion}, so \code{0 == ""} is \code{true}. The \textbf{strict equality} operator treats operands of different types as different. Two other surprises: \code{null == undefined} is \code{true} but \code{null == 0} is not, and \code{NaN} equals nothing, not even itself---use \code{Number.isNaN}.
\Needspace{6\baselineskip}
\begin{lstlisting}[language=JavaScript]
console.log(1 == "1", 1 === "1");
console.log(null == undefined, null == 0);
console.log(NaN === NaN, Number.isNaN(NaN));
\end{lstlisting}
{\small\textit{Prints:}~\code{true false}\enspace\textcolor{black!45}{|}\enspace \code{true false}\enspace\textcolor{black!45}{|}\enspace \code{false true}}\par

\cardend
\cardhead{Card 8 --- Fix getPort(config)}{fix · javascript · Easy}
\safemargin{\textbf{Terms for this card}\par\smallskip \textbf{Nullish coalescing (??):} returns its right side only when its left side is \code{null} or \code{undefined}. Like a substitute who steps in only if the player is absent, not if the player scores zero.\par\smallskip \textbf{Truthy / falsy:} whether a value counts as true or false in a condition. Like a bouncer's list: eight values are turned away, everyone else gets in.}
\noindent{\small\textbf{Covered in:} \hyperref[ch:c02]{Ch.~2} of this guide; documentation: \href{https://developer.mozilla.org/en-US/docs/Web/JavaScript/Reference/Operators/Nullish_coalescing}{MDN — Nullish coalescing operator (??)}.}\par
\cardsub{The card}
\textbf{Fix getPort(config)}\par
\code{getPort(config)} should return \code{config.port}, or \code{3000} when the port is missing (\code{undefined} or \code{null}). Port \code{0} is valid and must be kept. Fix the bug.\par
\begin{lstlisting}[language=JavaScript]
function getPort(config) {
  return config.port || 3000;
}
\end{lstlisting}
\cardsub{How to solve it}
\begin{cardrows}
\item[Solve] Replace \code{||} with \code{??}, which falls back only for \code{null} or \code{undefined}.
\item[Invariants] A port that is present, including 0, is returned unchanged.
\item[Memory] No state.
\item[Types] object in; number out.
\item[Failure] \code{||} (and a truthiness ternary) replace 0 as if it were missing.
\item[Tests] 8080, 0, a missing port, and \code{null}.
\end{cardrows}
\trap \code{||} treats a real 0 as missing.\par
\cardsub{Solution}
\begin{lstlisting}[language=JavaScript]
function getPort(config) {
  return config.port ?? 3000;
}
\end{lstlisting}
{\small \textit{Verified: passes all 4 test cases on the card in the app's code runner; the starter code fails them.}}\par
\smallskip{\small\textbf{Why:} \code{||} replaces every falsy value, 0 included; \code{??} replaces only null and undefined.}\par
\cardsub{From the readings}
\readchap{2. How does JavaScript compare values and choose a branch?}{\code{===} compares without converting; truthiness decides an \code{if}; and \code{||}, \code{\&\&} and \code{??} hand back one of their operands.}
\readq{What do ||, \&\& and ?? really return?}
One of their operands, not \code{true} or \code{false}. \code{a || b} returns \code{a} if it is truthy, else \code{b}; \code{a \&\& b} returns \code{a} if it is falsy, else \code{b}. The \textbf{nullish coalescing} operator \code{??} returns its right side only for \code{null} or \code{undefined}, so a real \code{0} or \code{""} survives. That makes it the right tool for defaults.
\Needspace{6\baselineskip}
\begin{lstlisting}[language=JavaScript]
const saved = 0;
console.log(saved || 10, saved ?? 10, null ?? 10);
\end{lstlisting}
{\small\textit{Prints:}~\code{10 0 10}}\par

\readq{What counts as true in an if?}
An \code{if} runs its block when the condition is \textbf{truthy}. Exactly eight values are falsy: \code{false}, \code{0}, \code{-0}, \code{0n}, \code{""}, \code{null}, \code{undefined} and \code{NaN}. Everything else is truthy, including \code{[]}, \code{\{\}} and the strings \code{"0"} and \code{"false"}.
\cardend
\cardhead{Card 9 --- letterGrade(score)}{write · javascript · Easy}
\noindent{\small\textbf{Covered in:} \hyperref[ch:c02]{Ch.~2} of this guide; documentation: \href{https://developer.mozilla.org/en-US/docs/Web/JavaScript/Reference/Statements/if...else}{MDN — if...else}.}\par
\cardsub{The card}
\textbf{letterGrade(score)}\par
Write \code{letterGrade(score)}: \code{\textquotedbl{}A\textquotedbl{}} for 90 and up, \code{\textquotedbl{}B\textquotedbl{}} for 80–89, \code{\textquotedbl{}C\textquotedbl{}} for 70–79, \code{\textquotedbl{}F\textquotedbl{}} below 70.\par
\begin{lstlisting}[language=JavaScript]
function letterGrade(score) {
  // TODO
  return "";
}
\end{lstlisting}
\cardsub{How to solve it}
\begin{cardrows}
\item[Solve] Return early from the highest band down: 90, then 80, then 70, else F.
\item[Invariants] Each later \code{if} only sees scores below the earlier thresholds.
\item[Memory] No state.
\item[Types] number in; one-letter string out.
\item[Failure] \code{\textgreater{}} instead of \code{\textgreater{}=} moves 90, 80 and 70 down a band; checking low bands first catches everything.
\item[Tests] Every boundary: 90, 89, 80, 79, 70, 69, plus 0 and 95.
\end{cardrows}
\trap Using \code{\textgreater{}} makes a score of exactly 90 a B.\par
\cardsub{Solution}
\begin{lstlisting}[language=JavaScript]
function letterGrade(score) {
  if (score >= 90) return "A";
  if (score >= 80) return "B";
  if (score >= 70) return "C";
  return "F";
}
\end{lstlisting}
{\small \textit{Verified: passes all 8 test cases on the card in the app's code runner; the starter code fails them.}}\par
\smallskip{\small\textbf{Why:} Check the highest band first and return early: each later \code{if} then only sees scores below the previous threshold. Test every boundary (90, 89, 80, \ldots{}).}\par
\cardsub{From the readings}
\readchap{2. How does JavaScript compare values and choose a branch?}{\code{===} compares without converting; truthiness decides an \code{if}; and \code{||}, \code{\&\&} and \code{??} hand back one of their operands.}
\readq{How do you write a chain of decisions?}
With \code{if} statements that return early, checked from the most demanding case down. Each later \code{if} then only sees what the earlier ones let through, so every boundary is decided once. Test each boundary value itself.
\Needspace{6\baselineskip}
\begin{lstlisting}[language=JavaScript]
function label(temp) {
  if (temp >= 30) return "hot";
  if (temp >= 15) return "mild";
  return "cold";
}
console.log(label(30), label(29), label(14));
\end{lstlisting}
{\small\textit{Prints:}~\code{hot mild cold}}\par

\cardend
\cardhead{Card 10 --- What does this print?}{predict · js02-logical-operands}
\safemargin{\textbf{Terms for this card}\par\smallskip \textbf{Nullish coalescing (??):} returns its right side only when its left side is \code{null} or \code{undefined}. Like a substitute who steps in only if the player is absent, not if the player scores zero.}
\noindent{\small\textbf{Covered in:} \hyperref[ch:c02]{Ch.~2} of this guide; documentation: \href{https://developer.mozilla.org/en-US/docs/Web/JavaScript/Reference/Operators/Logical_OR}{MDN — Logical OR (||)}.}\par
\cardsub{The card}
What does this print?\par
\begin{lstlisting}[language=JavaScript]
console.log("" || "fallback");
console.log("value" && 42);
console.log(0 ?? "zero");
console.log(null ?? "none");
console.log(0 || null);
\end{lstlisting}
\cardsub{How to solve it}
\begin{cardrows}
\item[Solve] For each line apply the rule: \code{||} gives the first truthy value (or the last), \code{\&\&} the first falsy (or the last), \code{??} the right side only for null/undefined.
\item[Invariants] These operators return an operand, never a fresh boolean.
\item[Memory] Nothing stored.
\item[Types] strings, numbers, null.
\item[Failure] Answering true/false, or thinking \code{0 ?? \textquotedbl{}zero\textquotedbl{}} gives "zero".
\item[Tests] Run it, then change one line and predict again.
\end{cardrows}
\trap \code{0 ?? \textquotedbl{}zero\textquotedbl{}} is 0, because 0 is not null or undefined.\par
\cardsub{Solution}
{\small\textbf{Prints:}}
\begin{lstlisting}
fallback
42
0
none
null
\end{lstlisting}
{\small\textit{Verified: Node v22.23.4 prints exactly this.}}\par
\smallskip{\small\textbf{Why:} a || b returns a if it is truthy, else b. a \&\& b returns a if it is falsy, else b. a ?? b returns b only when a is null or undefined, so 0 ?? "zero" is 0.}\par
\cardsub{From the readings}
\readchap{2. How does JavaScript compare values and choose a branch?}{\code{===} compares without converting; truthiness decides an \code{if}; and \code{||}, \code{\&\&} and \code{??} hand back one of their operands.}
\readq{What do ||, \&\& and ?? really return?}
One of their operands, not \code{true} or \code{false}. \code{a || b} returns \code{a} if it is truthy, else \code{b}; \code{a \&\& b} returns \code{a} if it is falsy, else \code{b}. The \textbf{nullish coalescing} operator \code{??} returns its right side only for \code{null} or \code{undefined}, so a real \code{0} or \code{""} survives. That makes it the right tool for defaults.
\Needspace{6\baselineskip}
\begin{lstlisting}[language=JavaScript]
const saved = 0;
console.log(saved || 10, saved ?? 10, null ?? 10);
\end{lstlisting}
{\small\textit{Prints:}~\code{10 0 10}}\par

\cardend
\cardhead{Card 11 --- countTruthy(values)}{write · javascript · Easy}
\safemargin{\textbf{Terms for this card}\par\smallskip \textbf{Truthy / falsy:} whether a value counts as true or false in a condition. Like a bouncer's list: eight values are turned away, everyone else gets in.}
\noindent{\small\textbf{Covered in:} \hyperref[ch:c02]{Ch.~2} of this guide; documentation: \href{https://developer.mozilla.org/en-US/docs/Web/JavaScript/Reference/Global_Objects/Array/filter}{MDN — Array.prototype.filter()}.}\par
\cardsub{The card}
\textbf{countTruthy(values)}\par
Write \code{countTruthy(values)}: how many items in the array are truthy?\par
\begin{lstlisting}[language=JavaScript]
function countTruthy(values) {
  let count = 0;
  // TODO
  return count;
}
\end{lstlisting}
\cardsub{How to solve it}
\begin{cardrows}
\item[Solve] \code{values.filter(Boolean).length} keeps exactly the truthy values.
\item[Invariants] Only the eight falsy values are dropped.
\item[Memory] A filtered copy of the array.
\item[Types] array of anything in; number out.
\item[Failure] Counting non-null values instead of truthy ones, or comparing with \code{!= false} (which drops \code{[]} and "0").
\item[Tests] \code{[]}, \code{\{\}}, "0" and "false" (truthy) next to 0, "", null and false.
\end{cardrows}
\trap "0", "false", [] and \{\} are truthy.\par
\cardsub{Solution}
\begin{lstlisting}[language=JavaScript]
function countTruthy(values) {
  return values.filter(Boolean).length;
}
\end{lstlisting}
{\small \textit{Verified: passes all 4 test cases on the card in the app's code runner; the starter code fails them.}}\par
\smallskip{\small\textbf{Why:} The falsy values are false, 0, -0, 0n, "", null, undefined and NaN. Everything else is truthy, including [], \{\} and the strings "0" and "false".}\par
\cardsub{From the readings}
\readchap{2. How does JavaScript compare values and choose a branch?}{\code{===} compares without converting; truthiness decides an \code{if}; and \code{||}, \code{\&\&} and \code{??} hand back one of their operands.}
\readq{What counts as true in an if?}
An \code{if} runs its block when the condition is \textbf{truthy}. Exactly eight values are falsy: \code{false}, \code{0}, \code{-0}, \code{0n}, \code{""}, \code{null}, \code{undefined} and \code{NaN}. Everything else is truthy, including \code{[]}, \code{\{\}} and the strings \code{"0"} and \code{"false"}.
\cardend
\deckhead{03 · Loops}{6}
\begin{corebox}\textbf{Core idea.} A loop repeats a block while a condition holds: \code{for} to count, \code{for...of} to walk values, and \code{return} or \code{break} to stop early.\end{corebox}
\begin{sumbox}{Deck summary sheet --- 03 · Loops}
\code{for} counts, \code{while} repeats until a condition fails, \code{for...of} walks the values of an array or string, and \code{for...in} walks an object's keys (as strings, including inherited ones). A loop can \code{return} or \code{break} early.\par
\smallskip{\footnotesize Read more: \href{https://developer.mozilla.org/en-US/docs/Web/JavaScript/Guide/Loops_and_iteration}{MDN — Loops and iteration}; \href{https://developer.mozilla.org/en-US/docs/Web/JavaScript/Reference/Statements/for...in}{MDN — for...in}; \href{https://developer.mozilla.org/en-US/docs/Web/JavaScript/Guide/Closures}{MDN — Closures}}\end{sumbox}
\cardhead{Card 12 --- sumTo(n)}{write · javascript · Easy}
\safemargin{\textbf{Terms for this card}\par\smallskip \textbf{Loop:} code that runs a block again and again. Like laps on a track with a lap counter.\par\smallskip \textbf{Iteration:} one run of a loop's body. One lap.\par\smallskip \textbf{Off-by-one error:} a loop that runs one time too many or too few. Like counting fence posts instead of gaps.}
\noindent{\small\textbf{Covered in:} \hyperref[ch:c03]{Ch.~3} of this guide; documentation: \href{https://developer.mozilla.org/en-US/docs/Web/JavaScript/Reference/Statements/for}{MDN — for}.}\par
\cardsub{The card}
\textbf{sumTo(n)}\par
Write \code{sumTo(n)} with a \code{for} loop: return 1 + 2 + \ldots{} + n. \code{sumTo(0)} is 0.\par
\begin{lstlisting}[language=JavaScript]
function sumTo(n) {
  let total = 0;
  // TODO: loop from 1 to n
  return total;
}
\end{lstlisting}
\cardsub{How to solve it}
\begin{cardrows}
\item[Solve] Loop \code{i} from 1 to n inclusive, adding each \code{i} to a running total.
\item[Invariants] After each lap, \code{total} is the sum of 1..i.
\item[Memory] Two numbers: \code{i} and \code{total}.
\item[Types] number in; number out.
\item[Failure] \code{i \textless{} n} stops one short (off by one).
\item[Tests] 0, 1, 5 and 100.
\end{cardrows}
\trap \code{\textless{}} instead of \code{\textless{}=} drops the last number.\par
\cardsub{Solution}
\begin{lstlisting}[language=JavaScript]
function sumTo(n) {
  let total = 0;
  for (let i = 1; i <= n; i++) {
    total += i;
  }
  return total;
}
\end{lstlisting}
{\small \textit{Verified: passes all 4 test cases on the card in the app's code runner; the starter code fails them.}}\par
\smallskip{\small\textbf{Why:} \code{i \textless{}= n} includes n itself; \code{i \textless{} n} is the classic off-by-one that stops one short.}\par
\cardsub{From the readings}
\readchap{3. How do you make code repeat?}{A loop repeats a block while a condition holds: \code{for} to count, \code{for...of} to walk values, and \code{return} or \code{break} to stop early.}
\readq{How does a counting for loop work?}
A \textbf{loop} repeats a block. A \code{for} loop has three parts: start (\code{let i = 1}), keep-going test (\code{i <= n}) and step (\code{i++}). Each run of the block is an \textbf{iteration}. Write \code{<=} to include the last value; \code{<} stops one short---the classic \textbf{off-by-one error}.

The remainder operator \code{\%} gives what is left after division, so \code{i \% 3 === 0} means ``divisible by 3''.
\Needspace{6\baselineskip}
\begin{lstlisting}[language=JavaScript]
let total = 0;
for (let i = 1; i <= 4; i++) total += i;
console.log(total, 9 % 3 === 0, 10 % 3);
\end{lstlisting}
{\small\textit{Prints:}~\code{10 true 1}}\par

\cardend
\cardhead{Card 13 --- countVowels(text)}{write · javascript · Easy}
\safemargin{\textbf{Terms for this card}\par\smallskip \textbf{Immutable:} cannot be changed after it is made; you make a new one instead. Like words printed in ink: to fix a typo you print a new copy.\par\smallskip \textbf{Method:} a function attached to a value and called with a dot, as in \code{word.toUpperCase()}. Like a button on an appliance.}
\noindent{\small\textbf{Covered in:} \hyperref[ch:c03]{Ch.~3} of this guide; documentation: \href{https://developer.mozilla.org/en-US/docs/Web/JavaScript/Reference/Statements/for...of}{MDN — for...of}.}\par
\cardsub{The card}
\textbf{countVowels(text)}\par
Write \code{countVowels(text)} using \code{for...of}: count a, e, i, o and u in either case.\par
\begin{lstlisting}[language=JavaScript]
function countVowels(text) {
  let count = 0;
  // TODO: for (const ch of ...)
  return count;
}
\end{lstlisting}
\cardsub{How to solve it}
\begin{cardrows}
\item[Solve] Lower-case the text once, walk it with \code{for...of}, and count characters found in "aeiou".
\item[Invariants] \code{count} equals the vowels seen so far.
\item[Memory] One counter; the lower-cased copy of the text.
\item[Types] string in; number out.
\item[Failure] Forgetting to lower-case, so capitals are missed.
\item[Tests] "AEIOU", "rhythm", "" and a sentence with capitals.
\end{cardrows}
\trap Upper-case vowels are skipped unless you lower-case first.\par
\cardsub{Solution}
\begin{lstlisting}[language=JavaScript]
function countVowels(text) {
  let count = 0;
  for (const ch of text.toLowerCase()) {
    if ("aeiou".includes(ch)) count++;
  }
  return count;
}
\end{lstlisting}
{\small \textit{Verified: passes all 5 test cases on the card in the app's code runner; the starter code fails them.}}\par
\smallskip{\small\textbf{Why:} \code{for...of} gives each character of a string. Lower-case once so one check covers both cases.}\par
\cardsub{From the readings}
\readchap{3. How do you make code repeat?}{A loop repeats a block while a condition holds: \code{for} to count, \code{for...of} to walk values, and \code{return} or \code{break} to stop early.}
\readq{How do you walk through the values of an array or string?}
With \code{for...of}, which hands you each value in turn. On a string it gives one character at a time.
\Needspace{6\baselineskip}
\begin{lstlisting}[language=JavaScript]
for (const ch of "hey") console.log(ch.toUpperCase());
\end{lstlisting}
{\small\textit{Prints:}~\code{H}\enspace\textcolor{black!45}{|}\enspace \code{E}\enspace\textcolor{black!45}{|}\enspace \code{Y}}\par

\readchap{4. How do you work with text?}{Strings never change: every string method returns a new string, and \code{+} joins text as soon as one side is a string.}
\readq{Can you change a string in place?}
No. Strings are \textbf{immutable}. A \textbf{method} such as \code{toUpperCase} or \code{slice} returns a new string; the original stays as it was. \code{slice(1)} gives everything after the first character. Reading \code{word[0]} on an empty string gives \code{undefined}, and calling a method on \code{undefined} throws.
\Needspace{6\baselineskip}
\begin{lstlisting}[language=JavaScript]
const word = "hello";
const shout = word.toUpperCase();
console.log(word, shout, word.slice(1, 3));
\end{lstlisting}
{\small\textit{Prints:}~\code{hello HELLO el}}\par

\cardend
\cardhead{Card 14 --- fizzBuzz(n)}{write · javascript · Easy}
\safemargin{\textbf{Terms for this card}\par\smallskip \textbf{Loop:} code that runs a block again and again. Like laps on a track with a lap counter.\par\smallskip \textbf{Iteration:} one run of a loop's body. One lap.\par\smallskip \textbf{Off-by-one error:} a loop that runs one time too many or too few. Like counting fence posts instead of gaps.}
\noindent{\small\textbf{Covered in:} \hyperref[ch:c03]{Ch.~3} of this guide; documentation: \href{https://developer.mozilla.org/en-US/docs/Web/JavaScript/Reference/Operators/Remainder}{MDN — Remainder (\%)}.}\par
\cardsub{The card}
\textbf{fizzBuzz(n)}\par
Write \code{fizzBuzz(n)}: return an array of strings for 1..n — \code{\textquotedbl{}Fizz\textquotedbl{}} for multiples of 3, \code{\textquotedbl{}Buzz\textquotedbl{}} for multiples of 5, \code{\textquotedbl{}FizzBuzz\textquotedbl{}} for both, otherwise the number as a string.\par
\begin{lstlisting}[language=JavaScript]
function fizzBuzz(n) {
  const out = [];
  // TODO
  return out;
}
\end{lstlisting}
\cardsub{How to solve it}
\begin{cardrows}
\item[Solve] Loop 1..n; test divisibility by 15 first, then 3, then 5, else push the number as a string.
\item[Invariants] The output array has exactly n entries, in order.
\item[Memory] One array growing by one string per lap.
\item[Types] number in; array of strings out.
\item[Failure] Testing 3 before 15 makes 15 "Fizz"; pushing the number itself gives the wrong type.
\item[Tests] 5, 15 and 0.
\end{cardrows}
\trap Check the both-divisible case before the single ones.\par
\cardsub{Solution}
\begin{lstlisting}[language=JavaScript]
function fizzBuzz(n) {
  const out = [];
  for (let i = 1; i <= n; i++) {
    if (i % 15 === 0) out.push("FizzBuzz");
    else if (i % 3 === 0) out.push("Fizz");
    else if (i % 5 === 0) out.push("Buzz");
    else out.push(String(i));
  }
  return out;
}
\end{lstlisting}
{\small \textit{Verified: passes all 3 test cases on the card in the app's code runner; the starter code fails them.}}\par
\smallskip{\small\textbf{Why:} Test the most specific case (divisible by both) first, or 15 is caught by the 3 check. The remainder operator \% tells you divisibility.}\par
\cardsub{From the readings}
\readchap{3. How do you make code repeat?}{A loop repeats a block while a condition holds: \code{for} to count, \code{for...of} to walk values, and \code{return} or \code{break} to stop early.}
\readq{How does a counting for loop work?}
A \textbf{loop} repeats a block. A \code{for} loop has three parts: start (\code{let i = 1}), keep-going test (\code{i <= n}) and step (\code{i++}). Each run of the block is an \textbf{iteration}. Write \code{<=} to include the last value; \code{<} stops one short---the classic \textbf{off-by-one error}.

The remainder operator \code{\%} gives what is left after division, so \code{i \% 3 === 0} means ``divisible by 3''.
\Needspace{6\baselineskip}
\begin{lstlisting}[language=JavaScript]
let total = 0;
for (let i = 1; i <= 4; i++) total += i;
console.log(total, 9 % 3 === 0, 10 % 3);
\end{lstlisting}
{\small\textit{Prints:}~\code{10 true 1}}\par

\cardend
\cardhead{Card 15 --- firstIndexOver(nums, limit)}{write · javascript · Easy}
\noindent{\small\textbf{Covered in:} \hyperref[ch:c03]{Ch.~3} of this guide; documentation: \href{https://developer.mozilla.org/en-US/docs/Web/JavaScript/Reference/Statements/return}{MDN — return}.}\par
\cardsub{The card}
\textbf{firstIndexOver(nums, limit)}\par
Write \code{firstIndexOver(nums, limit)}: return the index of the first number greater than \code{limit}, or -1 if there is none. Stop as soon as you find it.\par
\begin{lstlisting}[language=JavaScript]
function firstIndexOver(nums, limit) {
  // TODO
  return -1;
}
\end{lstlisting}
\cardsub{How to solve it}
\begin{cardrows}
\item[Solve] Walk the indexes and \code{return i} at the first number greater than the limit; return -1 after the loop.
\item[Invariants] Every index before \code{i} held a number that was not greater than the limit.
\item[Memory] Only the loop index.
\item[Types] array of numbers and a number in; number out.
\item[Failure] Recording matches without returning gives the last match; \code{\textgreater{}=} accepts equal numbers.
\item[Tests] A match in the middle, none at all, an empty array, and a number equal to the limit.
\end{cardrows}
\trap Without an early \code{return}, the loop finds the last match, not the first.\par
\cardsub{Solution}
\begin{lstlisting}[language=JavaScript]
function firstIndexOver(nums, limit) {
  for (let i = 0; i < nums.length; i++) {
    if (nums[i] > limit) return i;
  }
  return -1;
}
\end{lstlisting}
{\small \textit{Verified: passes all 5 test cases on the card in the app's code runner; the starter code fails them.}}\par
\smallskip{\small\textbf{Why:} \code{return} inside the loop ends it at the first match. Without it you would keep overwriting the answer and end up with the last match.}\par
\cardsub{From the readings}
\readchap{3. How do you make code repeat?}{A loop repeats a block while a condition holds: \code{for} to count, \code{for...of} to walk values, and \code{return} or \code{break} to stop early.}
\readq{How do you stop a loop early?}
\code{return} ends the whole function, loop included; \code{break} ends only the loop. To find the first match, return as soon as you see it; if you only record matches, you end up with the last one.
\Needspace{6\baselineskip}
\begin{lstlisting}[language=JavaScript]
for (const n of [3, 8, 1, 9]) {
  if (n > 5) {
    console.log("first big:", n);
    break;
  }
}
\end{lstlisting}
{\small\textit{Prints:}~\code{first big: 8}}\par

\cardend
\cardhead{Card 16 --- What does this print?}{predict · js03-closure-in-loop}
\safemargin{\textbf{Terms for this card}\par\smallskip \textbf{Closure:} a function together with the variables it could see when it was created. Like a snow globe: the scene inside is sealed in with it.}
\noindent{\small\textbf{Covered in:} \hyperref[ch:c03]{Ch.~3} of this guide; documentation: \href{https://developer.mozilla.org/en-US/docs/Web/JavaScript/Guide/Closures}{MDN — Closures}.}\par
\cardsub{The card}
What does this print?\par
\begin{lstlisting}[language=JavaScript]
const withVar = [];
for (var i = 0; i < 3; i++) {
  withVar.push(() => i);
}
const withLet = [];
for (let j = 0; j < 3; j++) {
  withLet.push(() => j);
}
console.log(withVar.map((f) => f()).join(" "));
console.log(withLet.map((f) => f()).join(" "));
\end{lstlisting}
\cardsub{How to solve it}
\begin{cardrows}
\item[Solve] With \code{var} there is one \code{i}, which ends at 3; with \code{let} each lap gets its own \code{j}.
\item[Invariants] Each arrow function reads the variable it closes over when it is finally called.
\item[Memory] One shared \code{i} versus three separate \code{j} variables.
\item[Types] functions and numbers.
\item[Failure] Answering "0 1 2" for both loops.
\item[Tests] Run it, then change one line and predict again.
\end{cardrows}
\trap \code{var} in a loop gives every function the same, final value.\par
\cardsub{Solution}
{\small\textbf{Prints:}}
\begin{lstlisting}
3 3 3
0 1 2
\end{lstlisting}
{\small\textit{Verified: Node v22.23.4 prints exactly this.}}\par
\smallskip{\small\textbf{Why:} var makes one i for the whole loop, so every function sees its final value, 3. let makes a fresh j for each iteration, so each function keeps its own.}\par
\cardsub{From the readings}
\readchap{3. How do you make code repeat?}{A loop repeats a block while a condition holds: \code{for} to count, \code{for...of} to walk values, and \code{return} or \code{break} to stop early.}
\readq{Why does let matter inside a loop?}
A function made inside a loop remembers the variable it can see. That pairing of a function with the variables around it is a \textbf{closure}. \code{let} creates a fresh variable for each iteration, so each function keeps its own; \code{var} makes one variable for the whole loop, so every function sees its final value.
\Needspace{6\baselineskip}
\begin{lstlisting}[language=JavaScript]
const fns = [];
for (let i = 0; i < 3; i++) fns.push(() => i);
console.log(fns.map((f) => f()).join(" "));
\end{lstlisting}
{\small\textit{Prints:}~\code{0 1 2}}\par

\cardend
\cardhead{Card 17 --- What does this print?}{predict · js03-for-in-vs-of}
\noindent{\small\textbf{Covered in:} \hyperref[ch:c03]{Ch.~3} of this guide; documentation: \href{https://developer.mozilla.org/en-US/docs/Web/JavaScript/Reference/Statements/for...in}{MDN — for...in}.}\par
\cardsub{The card}
What does this print?\par
\begin{lstlisting}[language=JavaScript]
const scores = [10, 20];
scores.label = "round 1";
const keys = [];
for (const key in scores) keys.push(key);
const values = [];
for (const value of scores) values.push(value);
console.log(keys.join(","));
console.log(values.join(","));
\end{lstlisting}
\cardsub{How to solve it}
\begin{cardrows}
\item[Solve] \code{for...in} lists keys (as strings), including the extra \code{label} property; \code{for...of} lists values.
\item[Invariants] \code{for...of} on an array yields only its items.
\item[Memory] One array with an extra named property.
\item[Types] strings (keys) and numbers (values).
\item[Failure] Forgetting that \code{label} appears, or that the keys are the strings "0" and "1".
\item[Tests] Run it, then change one line and predict again.
\end{cardrows}
\trap \code{for...in} on an array also visits added and inherited properties.\par
\cardsub{Solution}
{\small\textbf{Prints:}}
\begin{lstlisting}
0,1,label
10,20
\end{lstlisting}
{\small\textit{Verified: Node v22.23.4 prints exactly this.}}\par
\smallskip{\small\textbf{Why:} for...in visits every enumerable string key, extra and inherited properties included, and gives indexes as strings. Use for...of for the values of an array.}\par
\cardsub{From the readings}
\readchap{3. How do you make code repeat?}{A loop repeats a block while a condition holds: \code{for} to count, \code{for...of} to walk values, and \code{return} or \code{break} to stop early.}
\readq{Why is for...in the wrong loop for arrays?}
\code{for...in} walks an object's keys, as strings, including extra and inherited ones. On an array that means \code{"0"}, \code{"1"}, \ldots{} plus any property someone added. Use \code{for...of} for an array's values.
\readq{How do you walk through the values of an array or string?}
With \code{for...of}, which hands you each value in turn. On a string it gives one character at a time.
\Needspace{6\baselineskip}
\begin{lstlisting}[language=JavaScript]
for (const ch of "hey") console.log(ch.toUpperCase());
\end{lstlisting}
{\small\textit{Prints:}~\code{H}\enspace\textcolor{black!45}{|}\enspace \code{E}\enspace\textcolor{black!45}{|}\enspace \code{Y}}\par

\cardend
\deckhead{04 · Strings}{5}
\begin{corebox}\textbf{Core idea.} Strings never change: every string method returns a new string, and \code{+} joins text as soon as one side is a string.\end{corebox}
\begin{sumbox}{Deck summary sheet --- 04 · Strings}
Strings are immutable: \code{toUpperCase}, \code{slice}, \code{replace} and friends return new strings. \code{split} turns a string into an array and \code{join} turns it back. \code{+} with a string concatenates; the other arithmetic operators convert to numbers.\par
\smallskip{\footnotesize Read more: \href{https://developer.mozilla.org/en-US/docs/Web/JavaScript/Reference/Global_Objects/String}{MDN — String}}\end{sumbox}
\cardhead{Card 18 --- capitalize(word)}{write · javascript · Easy}
\safemargin{\textbf{Terms for this card}\par\smallskip \textbf{Immutable:} cannot be changed after it is made; you make a new one instead. Like words printed in ink: to fix a typo you print a new copy.\par\smallskip \textbf{Method:} a function attached to a value and called with a dot, as in \code{word.toUpperCase()}. Like a button on an appliance.}
\noindent{\small\textbf{Covered in:} \hyperref[ch:c04]{Ch.~4} of this guide; documentation: \href{https://developer.mozilla.org/en-US/docs/Web/JavaScript/Reference/Global_Objects/String/slice}{MDN — String.prototype.slice()}.}\par
\cardsub{The card}
\textbf{capitalize(word)}\par
Write \code{capitalize(word)}: first letter upper-case, the rest lower-case. \code{\textquotedbl{}\textquotedbl{}} stays \code{\textquotedbl{}\textquotedbl{}}.\par
\begin{lstlisting}[language=JavaScript]
function capitalize(word) {
  // TODO
  return word;
}
\end{lstlisting}
\cardsub{How to solve it}
\begin{cardrows}
\item[Solve] Return "" for empty input; otherwise upper-case \code{word[0]} and lower-case \code{word.slice(1)}.
\item[Invariants] The original string is never changed; a new one is returned.
\item[Memory] Two new strings joined together.
\item[Types] string in; string out.
\item[Failure] \code{word[0]} is undefined for "", so calling \code{toUpperCase} on it throws.
\item[Tests] "hello", "wORLD", "a" and "".
\end{cardrows}
\trap The empty string has no first character.\par
\cardsub{Solution}
\begin{lstlisting}[language=JavaScript]
function capitalize(word) {
  if (word === "") return "";
  return word[0].toUpperCase() + word.slice(1).toLowerCase();
}
\end{lstlisting}
{\small \textit{Verified: passes all 4 test cases on the card in the app's code runner; the starter code fails them.}}\par
\smallskip{\small\textbf{Why:} \code{word[0]} is undefined for "", and calling a method on undefined throws. \code{slice(1)} returns everything after the first character.}\par
\cardsub{From the readings}
\readchap{4. How do you work with text?}{Strings never change: every string method returns a new string, and \code{+} joins text as soon as one side is a string.}
\readq{Can you change a string in place?}
No. Strings are \textbf{immutable}. A \textbf{method} such as \code{toUpperCase} or \code{slice} returns a new string; the original stays as it was. \code{slice(1)} gives everything after the first character. Reading \code{word[0]} on an empty string gives \code{undefined}, and calling a method on \code{undefined} throws.
\Needspace{6\baselineskip}
\begin{lstlisting}[language=JavaScript]
const word = "hello";
const shout = word.toUpperCase();
console.log(word, shout, word.slice(1, 3));
\end{lstlisting}
{\small\textit{Prints:}~\code{hello HELLO el}}\par

\cardend
\cardhead{Card 19 --- isPalindrome(text)}{write · javascript · Easy}
\safemargin{\textbf{Terms for this card}\par\smallskip \textbf{Immutable:} cannot be changed after it is made; you make a new one instead. Like words printed in ink: to fix a typo you print a new copy.\par\smallskip \textbf{Method:} a function attached to a value and called with a dot, as in \code{word.toUpperCase()}. Like a button on an appliance.}
\noindent{\small\textbf{Covered in:} \hyperref[ch:c04]{Ch.~4} of this guide; documentation: \href{https://developer.mozilla.org/en-US/docs/Web/JavaScript/Reference/Global_Objects/Array/reverse}{MDN — Array.prototype.reverse()}.}\par
\cardsub{The card}
\textbf{isPalindrome(text)}\par
Write \code{isPalindrome(text)}: true if \code{text} reads the same backwards, ignoring case, spaces and punctuation.\par
\begin{lstlisting}[language=JavaScript]
function isPalindrome(text) {
  // TODO
  return false;
}
\end{lstlisting}
\cardsub{How to solve it}
\begin{cardrows}
\item[Solve] Lower-case, strip everything but letters and digits, then compare with the reversed copy.
\item[Invariants] Only letters and digits are compared, in one case.
\item[Memory] The cleaned string and its reversed copy.
\item[Types] string in; boolean out.
\item[Failure] Skipping the lower-casing or the clean-up makes "Race Car" fail.
\item[Tests] "racecar", "Race Car", "A man, a plan, a canal: Panama", "hello" and "".
\end{cardrows}
\trap Strings have no reverse(): spread to an array, reverse, join.\par
\cardsub{Solution}
\begin{lstlisting}[language=JavaScript]
function isPalindrome(text) {
  const letters = text.toLowerCase().replace(/[^a-z0-9]/g, "");
  return letters === [...letters].reverse().join("");
}
\end{lstlisting}
{\small \textit{Verified: passes all 5 test cases on the card in the app's code runner; the starter code fails them.}}\par
\smallskip{\small\textbf{Why:} Clean the text first, then compare it with its reverse. Strings have no reverse(): spread into an array, reverse, and join back.}\par
\cardsub{From the readings}
\readchap{4. How do you work with text?}{Strings never change: every string method returns a new string, and \code{+} joins text as soon as one side is a string.}
\readq{How do you break text apart and put it back?}
\code{split} cuts a string into an array of pieces at a pattern; \code{join} glues an array back into a string. \code{trim} removes spaces from both ends. Splitting on a single space leaves empty pieces between repeated spaces; the pattern \code{/\textbackslash s+/} (one or more whitespace characters) does not. Strings have no \code{reverse}: spread the string into an array, reverse that, and join.
\Needspace{6\baselineskip}
\begin{lstlisting}[language=JavaScript]
const parts = "  a   b ".trim().split(/\s+/);
console.log(parts.length, parts.join("-"));
console.log([..."abc"].reverse().join(""));
\end{lstlisting}
{\small\textit{Prints:}~\code{2 a-b}\enspace\textcolor{black!45}{|}\enspace \code{cba}}\par

\readq{Can you change a string in place?}
No. Strings are \textbf{immutable}. A \textbf{method} such as \code{toUpperCase} or \code{slice} returns a new string; the original stays as it was. \code{slice(1)} gives everything after the first character. Reading \code{word[0]} on an empty string gives \code{undefined}, and calling a method on \code{undefined} throws.
\Needspace{6\baselineskip}
\begin{lstlisting}[language=JavaScript]
const word = "hello";
const shout = word.toUpperCase();
console.log(word, shout, word.slice(1, 3));
\end{lstlisting}
{\small\textit{Prints:}~\code{hello HELLO el}}\par

\cardend
\cardhead{Card 20 --- initials(fullName)}{write · javascript · Easy}
\noindent{\small\textbf{Covered in:} \hyperref[ch:c04]{Ch.~4} of this guide; documentation: \href{https://developer.mozilla.org/en-US/docs/Web/JavaScript/Reference/Global_Objects/String/split}{MDN — String.prototype.split()}.}\par
\cardsub{The card}
\textbf{initials(fullName)}\par
Write \code{initials(fullName)}: \code{\textquotedbl{}Ada Lovelace\textquotedbl{}} $\rightarrow$ \code{\textquotedbl{}A.L.\textquotedbl{}}. Names may be lower-case and have extra spaces.\par
\begin{lstlisting}[language=JavaScript]
function initials(fullName) {
  // TODO
  return "";
}
\end{lstlisting}
\cardsub{How to solve it}
\begin{cardrows}
\item[Solve] Trim, split on \code{/\textbackslash{}s+/}, take each part's first letter upper-cased plus a dot, and join.
\item[Invariants] Every part from the split is a non-empty word.
\item[Memory] An array of words, then an array of initials.
\item[Types] string in; string out.
\item[Failure] \code{split(\textquotedbl{} \textquotedbl{})} leaves empty strings where spaces repeat, and \code{part[0]} is then undefined.
\item[Tests] Two names, three names, extra spaces, and a single name.
\end{cardrows}
\trap Splitting on one space breaks on repeated spaces.\par
\cardsub{Solution}
\begin{lstlisting}[language=JavaScript]
function initials(fullName) {
  return fullName
    .trim()
    .split(/\s+/)
    .map((part) => part[0].toUpperCase() + ".")
    .join("");
}
\end{lstlisting}
{\small \textit{Verified: passes all 4 test cases on the card in the app's code runner; the starter code fails them.}}\par
\smallskip{\small\textbf{Why:} \code{split(\textquotedbl{} \textquotedbl{})} leaves empty strings between repeated spaces; trim first and split on the pattern /\textbackslash{}s+/ (one or more whitespace characters).}\par
\cardsub{From the readings}
\readchap{4. How do you work with text?}{Strings never change: every string method returns a new string, and \code{+} joins text as soon as one side is a string.}
\readq{How do you break text apart and put it back?}
\code{split} cuts a string into an array of pieces at a pattern; \code{join} glues an array back into a string. \code{trim} removes spaces from both ends. Splitting on a single space leaves empty pieces between repeated spaces; the pattern \code{/\textbackslash s+/} (one or more whitespace characters) does not. Strings have no \code{reverse}: spread the string into an array, reverse that, and join.
\Needspace{6\baselineskip}
\begin{lstlisting}[language=JavaScript]
const parts = "  a   b ".trim().split(/\s+/);
console.log(parts.length, parts.join("-"));
console.log([..."abc"].reverse().join(""));
\end{lstlisting}
{\small\textit{Prints:}~\code{2 a-b}\enspace\textcolor{black!45}{|}\enspace \code{cba}}\par

\cardend
\cardhead{Card 21 --- What does this print?}{predict · js04-plus-coercion}
\noindent{\small\textbf{Covered in:} \hyperref[ch:c04]{Ch.~4} of this guide; documentation: \href{https://developer.mozilla.org/en-US/docs/Web/JavaScript/Reference/Operators/Addition}{MDN — Addition (+)}.}\par
\cardsub{The card}
What does this print?\par
\begin{lstlisting}[language=JavaScript]
console.log("5" + 3);
console.log("5" - 3);
console.log("5" * "2");
console.log(1 + 2 + "3");
console.log("1" + 2 + 3);
\end{lstlisting}
\cardsub{How to solve it}
\begin{cardrows}
\item[Solve] For each \code{+}, ask whether either side is a string; for \code{-} and \code{*}, convert to numbers. Go left to right.
\item[Invariants] \code{+} joins as soon as one side is a string.
\item[Memory] Nothing stored.
\item[Types] strings and numbers.
\item[Failure] Answering 8 for \code{\textquotedbl{}5\textquotedbl{} + 3}, or "123" mistaken for 6 or "33".
\item[Tests] Run it, then change one line and predict again.
\end{cardrows}
\trap \code{1 + 2 + \textquotedbl{}3\textquotedbl{}} adds first (3), then joins: "33".\par
\cardsub{Solution}
{\small\textbf{Prints:}}
\begin{lstlisting}
53
2
10
33
123
\end{lstlisting}
{\small\textit{Verified: Node v22.23.4 prints exactly this.}}\par
\smallskip{\small\textbf{Why:} + concatenates as soon as either side is a string; -, * and / convert to numbers. Evaluation runs left to right: 1 + 2 is 3 before "3" is joined on.}\par
\cardsub{From the readings}
\readchap{4. How do you work with text?}{Strings never change: every string method returns a new string, and \code{+} joins text as soon as one side is a string.}
\readq{Why does "5" + 3 give "53"?}
Because \code{+} both adds numbers and joins strings: as soon as either side is a string, it joins. The other arithmetic operators (\code{-}, \code{*}, \code{/}) convert their operands to numbers. Evaluation runs left to right, so \code{1 + 2 + "3"} adds first and joins second.
\Needspace{6\baselineskip}
\begin{lstlisting}[language=JavaScript]
console.log("5" + 3, "5" - 3, 1 + 2 + "3");
\end{lstlisting}
{\small\textit{Prints:}~\code{53 2 33}}\par

\cardend
\cardhead{Card 22 --- What does this print?}{predict · js04-immutable}
\safemargin{\textbf{Terms for this card}\par\smallskip \textbf{Immutable:} cannot be changed after it is made; you make a new one instead. Like words printed in ink: to fix a typo you print a new copy.\par\smallskip \textbf{Method:} a function attached to a value and called with a dot, as in \code{word.toUpperCase()}. Like a button on an appliance.}
\noindent{\small\textbf{Covered in:} \hyperref[ch:c04]{Ch.~4} of this guide; documentation: \href{https://developer.mozilla.org/en-US/docs/Web/JavaScript/Guide/Data_structures}{MDN — JavaScript data types}.}\par
\cardsub{The card}
What does this print?\par
\begin{lstlisting}[language=JavaScript]
const word = "cat";
const upper = word.toUpperCase();
const replaced = word.replace("c", "b");
console.log(word, upper, replaced);
\end{lstlisting}
\cardsub{How to solve it}
\begin{cardrows}
\item[Solve] Each method returns a new string; \code{word} itself is unchanged.
\item[Invariants] A string never changes after it is made.
\item[Memory] Three separate strings.
\item[Types] strings.
\item[Failure] Thinking \code{toUpperCase} or \code{replace} changed \code{word}.
\item[Tests] Run it, then change one line and predict again.
\end{cardrows}
\trap String methods return new strings; the original stays.\par
\cardsub{Solution}
{\small\textbf{Prints:}}
\begin{lstlisting}
cat CAT bat
\end{lstlisting}
{\small\textit{Verified: Node v22.23.4 prints exactly this.}}\par
\smallskip{\small\textbf{Why:} String methods return new strings; the original never changes.}\par
\cardsub{From the readings}
\readchap{4. How do you work with text?}{Strings never change: every string method returns a new string, and \code{+} joins text as soon as one side is a string.}
\readq{Can you change a string in place?}
No. Strings are \textbf{immutable}. A \textbf{method} such as \code{toUpperCase} or \code{slice} returns a new string; the original stays as it was. \code{slice(1)} gives everything after the first character. Reading \code{word[0]} on an empty string gives \code{undefined}, and calling a method on \code{undefined} throws.
\Needspace{6\baselineskip}
\begin{lstlisting}[language=JavaScript]
const word = "hello";
const shout = word.toUpperCase();
console.log(word, shout, word.slice(1, 3));
\end{lstlisting}
{\small\textit{Prints:}~\code{hello HELLO el}}\par

\cardend
\deckhead{05 · Functions, Scope \& Closures}{4}
\begin{corebox}\textbf{Core idea.} A function bundles steps behind a name; declarations are usable before their line, and a closure keeps the variables it was created with.\end{corebox}
\begin{sumbox}{Deck summary sheet --- 05 · Functions, Scope \& Closures}
Function declarations are hoisted; function expressions assigned to \code{const} are not usable before their line. A closure is a function that keeps using the variables of the scope it was created in, which is how you keep state private. Rest parameters (\code{...args}) collect any number of arguments into a real array.\par
\smallskip{\footnotesize Read more: \href{https://developer.mozilla.org/en-US/docs/Web/JavaScript/Guide/Functions}{MDN — Functions}; \href{https://developer.mozilla.org/en-US/docs/Web/JavaScript/Guide/Closures}{MDN — Closures}}\end{sumbox}
\cardhead{Card 23 --- What does this print?}{predict · js05-hoisting}
\safemargin{\textbf{Terms for this card}\par\smallskip \textbf{Function:} a named, reusable block of steps that takes inputs and can return a result. Like a recipe card: write it once, cook from it many times.\par\smallskip \textbf{Parameter:} a name in a function's definition that receives an input. A blank on the recipe card, e.g. ``\_\_ cups of flour''.\par\smallskip \textbf{Hoisting:} JavaScript making a declaration usable before the line it is written on. Like a speaker whose name is on the programme before they walk on stage.}
\noindent{\small\textbf{Covered in:} \hyperref[ch:c05]{Ch.~5} of this guide; documentation: \href{https://developer.mozilla.org/en-US/docs/Glossary/Hoisting}{MDN — Hoisting}.}\par
\cardsub{The card}
What does this print?\par
\begin{lstlisting}[language=JavaScript]
console.log(declared());
try {
  console.log(expressed());
} catch (error) {
  console.log(error.name);
}
function declared() {
  return "declared works";
}
const expressed = function () {
  return "expressed works";
};
\end{lstlisting}
\cardsub{How to solve it}
\begin{cardrows}
\item[Solve] The declaration is usable before its line; the \code{const} function is not, so calling it throws.
\item[Invariants] A \code{const} name cannot be used before its line runs.
\item[Memory] Both names exist from the start; only the declaration has a function yet.
\item[Types] string and ReferenceError.
\item[Failure] Expecting both calls to work, or the error to be a TypeError.
\item[Tests] Run it, then change one line and predict again.
\end{cardrows}
\trap Function declarations are hoisted; function expressions in \code{const} are not.\par
\cardsub{Solution}
{\small\textbf{Prints:}}
\begin{lstlisting}
declared works
ReferenceError
\end{lstlisting}
{\small\textit{Verified: Node v22.23.4 prints exactly this.}}\par
\smallskip{\small\textbf{Why:} A function declaration is created before any code runs. A const is not usable until its line runs, so calling it earlier throws a ReferenceError.}\par
\cardsub{From the readings}
\readchap{5. How do you package logic so it can be reused?}{A function bundles steps behind a name; declarations are usable before their line, and a closure keeps the variables it was created with.}
\readq{What is a function, and what does hoisting change?}
A \textbf{function} takes \textbf{parameters}, runs its body and hands back a value with \code{return}. A function declaration is \textbf{hoisted}: you can call it above the line where it is written. A function stored in a \code{const} is not usable until that line runs.
\Needspace{6\baselineskip}
\begin{lstlisting}[language=JavaScript]
console.log(double(4));
function double(n) {
  return n * 2;
}
\end{lstlisting}
{\small\textit{Prints:}~\code{8}}\par

\cardend
\cardhead{Card 24 --- makeCounter(start)}{write · javascript · Medium}
\safemargin{\textbf{Terms for this card}\par\smallskip \textbf{Arrow function:} the short form \code{(x) => x * 2} of a function expression. Like a sticky note instead of a recipe card.\par\smallskip \textbf{Closure:} a function together with the variables it could see when it was created. Like a snow globe: the scene inside is sealed in with it.}
\noindent{\small\textbf{Covered in:} \hyperref[ch:c05]{Ch.~5} of this guide; documentation: \href{https://developer.mozilla.org/en-US/docs/Web/JavaScript/Guide/Closures}{MDN — Closures}.}\par
\cardsub{The card}
\textbf{makeCounter(start)}\par
Write \code{makeCounter(start)}: return an object with \code{increment()} (adds 1, returns the new count) and \code{current()}. Each counter must keep its own private count. Don't change \code{check}.\par
\begin{lstlisting}[language=JavaScript]
// Test driver: don't change.
function check(start) {
  const a = makeCounter(start);
  const b = makeCounter(100);
  return [a.increment(), a.increment(), b.increment(), a.current(), b.current()];
}

function makeCounter(start) {
  // TODO: keep the count in a variable the returned functions close over
  return { increment() {}, current() {} };
}
\end{lstlisting}
\cardsub{How to solve it}
\begin{cardrows}
\item[Solve] Declare \code{let count = start} inside \code{makeCounter} and return methods that update and read it.
\item[Invariants] Each counter's \code{count} changes only through its own methods.
\item[Memory] One \code{count} variable per call to \code{makeCounter}.
\item[Types] number in; object with two functions out.
\item[Failure] A \code{count} declared outside the function is shared by every counter.
\item[Tests] Two counters stepped in turn must not affect each other.
\end{cardrows}
\trap Put the state inside the outer function, not beside it.\par
\cardsub{Solution}
\begin{lstlisting}[language=JavaScript]
// Test driver: don't change.
function check(start) {
  const a = makeCounter(start);
  const b = makeCounter(100);
  return [a.increment(), a.increment(), b.increment(), a.current(), b.current()];
}

function makeCounter(start) {
  let count = start;
  return {
    increment() {
      count += 1;
      return count;
    },
    current() {
      return count;
    },
  };
}
\end{lstlisting}
{\small \textit{Verified: passes all 2 test cases on the card in the app's code runner; the starter code fails them.}}\par
\smallskip{\small\textbf{Why:} Each call to makeCounter creates a new \code{count}, and the returned methods close over that one. A variable outside the function would be shared by every counter.}\par
\cardsub{From the readings}
\readchap{5. How do you package logic so it can be reused?}{A function bundles steps behind a name; declarations are usable before their line, and a closure keeps the variables it was created with.}
\readq{How can a function keep private state?}
By closing over a variable (\hyperref[ch:c03]{Ch.~3}). Each call to an outer function creates new variables; the inner functions it returns keep using those, and nothing outside can reach them. An \textbf{arrow function} is a compact way to write such an inner function.
\Needspace{6\baselineskip}
\begin{lstlisting}[language=JavaScript]
function makeGreeter(greeting) {
  return (name) => `${greeting}, ${name}`;
}
const hi = makeGreeter("Hi");
const yo = makeGreeter("Yo");
console.log(hi("Ada"), yo("Ada"));
\end{lstlisting}
{\small\textit{Prints:}~\code{Hi, Ada Yo, Ada}}\par

To remember ``already done'', keep a separate flag. Using the stored result as the flag fails when the result is itself \code{undefined}.
\cardend
\cardhead{Card 25 --- once(fn)}{write · javascript · Medium}
\safemargin{\textbf{Terms for this card}\par\smallskip \textbf{Arrow function:} the short form \code{(x) => x * 2} of a function expression. Like a sticky note instead of a recipe card.}
\noindent{\small\textbf{Covered in:} \hyperref[ch:c05]{Ch.~5} of this guide; documentation: \href{https://developer.mozilla.org/en-US/docs/Web/JavaScript/Guide/Closures}{MDN — Closures}.}\par
\cardsub{The card}
\textbf{once(fn)}\par
Write \code{once(fn)}: return a function that calls \code{fn} the first time only and returns that first result on every later call — even when the result is \code{undefined}. Don't change \code{check}.\par
\begin{lstlisting}[language=JavaScript]
// Test driver: don't change.
function check(values) {
  let calls = 0;
  const double = once((x) => {
    calls++;
    return x * 2;
  });
  const results = values.map((v) => double(v));
  let logs = 0;
  const log = once(() => {
    logs++;
  });
  log();
  log();
  return { results, calls, logs };
}

function once(fn) {
  // TODO
  return fn;
}
\end{lstlisting}
\cardsub{How to solve it}
\begin{cardrows}
\item[Solve] Keep a \code{done} flag and a \code{result} in the closure; call \code{fn} only while \code{done} is false.
\item[Invariants] \code{fn} runs at most once; every call returns the first result.
\item[Memory] Two closed-over variables: \code{done} and \code{result}.
\item[Types] function in; function out.
\item[Failure] Using \code{result === undefined} as the flag calls \code{fn} again when it returns undefined.
\item[Tests] Several calls; a function that returns undefined; no calls at all.
\end{cardrows}
\trap A result of undefined is not "not called yet".\par
\cardsub{Solution}
\begin{lstlisting}[language=JavaScript]
// Test driver: don't change.
function check(values) {
  let calls = 0;
  const double = once((x) => {
    calls++;
    return x * 2;
  });
  const results = values.map((v) => double(v));
  let logs = 0;
  const log = once(() => {
    logs++;
  });
  log();
  log();
  return { results, calls, logs };
}

function once(fn) {
  let done = false;
  let result;
  return (...args) => {
    if (!done) {
      done = true;
      result = fn(...args);
    }
    return result;
  };
}
\end{lstlisting}
{\small \textit{Verified: passes all 2 test cases on the card in the app's code runner; the starter code fails them.}}\par
\smallskip{\small\textbf{Why:} Track "already called" in its own flag. Using the result itself as the flag breaks for a function that returns undefined: it gets called again.}\par
\cardsub{From the readings}
\readchap{5. How do you package logic so it can be reused?}{A function bundles steps behind a name; declarations are usable before their line, and a closure keeps the variables it was created with.}
\readq{How can a function keep private state?}
By closing over a variable (\hyperref[ch:c03]{Ch.~3}). Each call to an outer function creates new variables; the inner functions it returns keep using those, and nothing outside can reach them. An \textbf{arrow function} is a compact way to write such an inner function.
\Needspace{6\baselineskip}
\begin{lstlisting}[language=JavaScript]
function makeGreeter(greeting) {
  return (name) => `${greeting}, ${name}`;
}
const hi = makeGreeter("Hi");
const yo = makeGreeter("Yo");
console.log(hi("Ada"), yo("Ada"));
\end{lstlisting}
{\small\textit{Prints:}~\code{Hi, Ada Yo, Ada}}\par

To remember ``already done'', keep a separate flag. Using the stored result as the flag fails when the result is itself \code{undefined}.
\cardend
\cardhead{Card 26 --- sum(...numbers)}{write · javascript · Easy}
\safemargin{\textbf{Terms for this card}\par\smallskip \textbf{Rest parameter:} \code{...name} as the last parameter, gathering all remaining arguments into an array. Like a sack that holds however many apples you hand over.\par\smallskip \textbf{Accumulator:} the running result \code{reduce} passes from one item to the next. Like a running tally on a notepad.}
\noindent{\small\textbf{Covered in:} \hyperref[ch:c05]{Ch.~5}, \hyperref[ch:c06]{Ch.~6} of this guide; documentation: \href{https://developer.mozilla.org/en-US/docs/Web/JavaScript/Reference/Functions/rest_parameters}{MDN — Rest parameters}.}\par
\cardsub{The card}
\textbf{sum(...numbers)}\par
Write \code{sum(...numbers)}: add up any number of arguments. \code{sum()} is 0.\par
\begin{lstlisting}[language=JavaScript]
function sum(numbers) {
  // TODO
  return 0;
}
\end{lstlisting}
\cardsub{How to solve it}
\begin{cardrows}
\item[Solve] Collect the arguments with \code{...numbers} and \code{reduce} them from 0.
\item[Invariants] The total starts at 0 and adds each argument once.
\item[Memory] One array of arguments.
\item[Types] any number of numbers in; number out.
\item[Failure] Without the starting 0, \code{sum()} throws; a plain parameter receives only the first argument.
\item[Tests] \code{sum()}, one argument, several, and negatives.
\end{cardrows}
\trap \code{reduce} needs a starting value for zero arguments.\par
\cardsub{Solution}
\begin{lstlisting}[language=JavaScript]
function sum(...numbers) {
  return numbers.reduce((total, n) => total + n, 0);
}
\end{lstlisting}
{\small \textit{Verified: passes all 4 test cases on the card in the app's code runner; the starter code fails them.}}\par
\smallskip{\small\textbf{Why:} \code{...numbers} gathers all arguments into a real array. reduce needs the initial value 0, or an empty array throws a TypeError.}\par
\cardsub{From the readings}
\readchap{5. How do you package logic so it can be reused?}{A function bundles steps behind a name; declarations are usable before their line, and a closure keeps the variables it was created with.}
\readq{How does a function accept any number of arguments?}
With a \textbf{rest parameter}: \code{function sum(...numbers)} receives every argument in a real array. Adding them up with \code{reduce} (\hyperref[ch:c06]{Ch.~6}) needs a starting value, or zero arguments would throw.
\Needspace{6\baselineskip}
\begin{lstlisting}[language=JavaScript]
function count(...items) {
  return items.length;
}
console.log(count(), count("a", "b"));
\end{lstlisting}
{\small\textit{Prints:}~\code{0 2}}\par

\readchap{6. How do you work with a list of values?}{\code{map} transforms, \code{filter} keeps, \code{reduce} folds, \code{forEach} only acts; \code{sort} changes the array in place and compares as text unless told otherwise.}
\readq{How do you boil a list down to one value?}
With \code{reduce}. It carries an \textbf{accumulator} through the array: each callback gets the tally so far and the next item and returns the new tally. Give it a starting value (\code{0}, \code{\{\}}, \ldots): without one, an empty array throws a \code{TypeError}. The callback must return the accumulator every time.
\Needspace{6\baselineskip}
\begin{lstlisting}[language=JavaScript]
const add = (sum, n) => sum + n;
console.log([4, 6].reduce(add, 0), [].reduce(add, 0));
try {
  [].reduce(add);
} catch (error) {
  console.log(error.name);
}
\end{lstlisting}
{\small\textit{Prints:}~\code{10 0}\enspace\textcolor{black!45}{|}\enspace \code{TypeError}}\par

\cardend
\deckhead{06 · Arrays \& Iteration}{6}
\begin{corebox}\textbf{Core idea.} \code{map} transforms, \code{filter} keeps, \code{reduce} folds, \code{forEach} only acts; \code{sort} changes the array in place and compares as text unless told otherwise.\end{corebox}
\begin{sumbox}{Deck summary sheet --- 06 · Arrays \& Iteration}
\code{map} transforms every item, \code{filter} keeps some, \code{reduce} folds them into one value (give it a starting value), and \code{forEach} only runs side effects — it always returns undefined. \code{sort} sorts in place and, without a compare function, as strings. Assigning an array copies the reference, not the items.\par
\smallskip{\footnotesize Read more: \href{https://developer.mozilla.org/en-US/docs/Web/JavaScript/Reference/Global_Objects/Array/sort}{MDN — Array.prototype.sort()}; \href{https://developer.mozilla.org/en-US/docs/Web/JavaScript/Reference/Global_Objects/Array/reduce}{MDN — Array.prototype.reduce()}}\end{sumbox}
\cardhead{Card 27 --- What does this print?}{predict · js06-sort-default}
\safemargin{\textbf{Terms for this card}\par\smallskip \textbf{In place:} changing the original rather than making a new copy. Like rearranging the books on your shelf instead of buying a sorted set.}
\noindent{\small\textbf{Covered in:} \hyperref[ch:c06]{Ch.~6} of this guide; documentation: \href{https://developer.mozilla.org/en-US/docs/Web/JavaScript/Reference/Global_Objects/Array/sort}{MDN — Array.prototype.sort()}.}\par
\cardsub{The card}
What does this print?\par
\begin{lstlisting}[language=JavaScript]
const nums = [10, 9, 1, 100];
nums.sort();
console.log(nums.join(" "));
nums.sort((a, b) => a - b);
console.log(nums.join(" "));
\end{lstlisting}
\cardsub{How to solve it}
\begin{cardrows}
\item[Solve] The first sort compares as text ("100" < "9"); the second compares numbers.
\item[Invariants] \code{sort} changes the array in place each time.
\item[Memory] One array, re-ordered twice.
\item[Types] numbers compared as strings, then as numbers.
\item[Failure] Expecting "1 9 10 100" from the first sort.
\item[Tests] Run it, then change one line and predict again.
\end{cardrows}
\trap \code{sort()} without a compare function sorts numbers as strings.\par
\cardsub{Solution}
{\small\textbf{Prints:}}
\begin{lstlisting}
1 10 100 9
1 9 10 100
\end{lstlisting}
{\small\textit{Verified: Node v22.23.4 prints exactly this.}}\par
\smallskip{\small\textbf{Why:} Without a compare function, sort compares the items as strings, and "100" comes before "9". (a, b) => a - b sorts numbers ascending. sort also changes the array in place.}\par
\cardsub{From the readings}
\readchap{6. How do you work with a list of values?}{\code{map} transforms, \code{filter} keeps, \code{reduce} folds, \code{forEach} only acts; \code{sort} changes the array in place and compares as text unless told otherwise.}
\readq{Why does sort put 100 before 9?}
Without a compare function, \code{sort} compares items as strings, and \code{"100"} comes before \code{"9"}. Pass \code{(a, b) => a - b} to sort numbers ascending. \code{sort} also changes the array \textbf{in place}; copy first (\code{[...nums]}) if the caller's array must stay as it was.
\Needspace{6\baselineskip}
\begin{lstlisting}[language=JavaScript]
console.log([10, 9, 100].sort().join(" "));
console.log([10, 9, 100].sort((a, b) => a - b).join(" "));
\end{lstlisting}
{\small\textit{Prints:}~\code{10 100 9}\enspace\textcolor{black!45}{|}\enspace \code{9 10 100}}\par

\cardend
\cardhead{Card 28 --- squaresOfEvens(nums)}{write · javascript · Easy}
\safemargin{\textbf{Terms for this card}\par\smallskip \textbf{Array:} an ordered list of values, numbered from 0. Like a row of numbered lockers.\par\smallskip \textbf{Callback:} a function you pass to another function, which calls it for you. Like instructions you hand a helper: ``do this with each item''.}
\noindent{\small\textbf{Covered in:} \hyperref[ch:c06]{Ch.~6} of this guide; documentation: \href{https://developer.mozilla.org/en-US/docs/Web/JavaScript/Reference/Global_Objects/Array/map}{MDN — Array.prototype.map()}.}\par
\cardsub{The card}
\textbf{squaresOfEvens(nums)}\par
Write \code{squaresOfEvens(nums)} with \code{filter} and \code{map}: the squares of the even numbers, in order.\par
\begin{lstlisting}[language=JavaScript]
function squaresOfEvens(nums) {
  // TODO: filter, then map
  return nums;
}
\end{lstlisting}
\cardsub{How to solve it}
\begin{cardrows}
\item[Solve] \code{filter} the even numbers (\code{n \% 2 === 0}), then \code{map} each to \code{n * n}.
\item[Invariants] The output keeps the input order.
\item[Memory] Two new arrays; the input is untouched.
\item[Types] array of numbers in; array of numbers out.
\item[Failure] Testing \code{\% 2 === 1} drops negatives and keeps odds; mapping without filtering squares everything.
\item[Tests] Mixed, empty, with 0 and a negative even number.
\end{cardrows}
\trap \code{-3 \% 2} is -1, not 1.\par
\cardsub{Solution}
\begin{lstlisting}[language=JavaScript]
function squaresOfEvens(nums) {
  return nums.filter((n) => n % 2 === 0).map((n) => n * n);
}
\end{lstlisting}
{\small \textit{Verified: passes all 4 test cases on the card in the app's code runner; the starter code fails them.}}\par
\smallskip{\small\textbf{Why:} filter keeps the items whose callback returns true; map returns a new array of the same length.}\par
\cardsub{From the readings}
\readchap{6. How do you work with a list of values?}{\code{map} transforms, \code{filter} keeps, \code{reduce} folds, \code{forEach} only acts; \code{sort} changes the array in place and compares as text unless told otherwise.}
\readq{How do you transform or select items?}
An \textbf{array} has methods that take a \textbf{callback}. \code{map} returns a new array of the same length with each item transformed; \code{filter} returns a new array of the items whose callback returned a truthy value.
\Needspace{6\baselineskip}
\begin{lstlisting}[language=JavaScript]
const prices = [5, 20, 8];
const big = prices.filter((p) => p > 6).map((p) => p * 2);
console.log(big.join(" "), prices.join(" "));
\end{lstlisting}
{\small\textit{Prints:}~\code{40 16 5 20 8}}\par

\cardend
\cardhead{Card 29 --- average(nums)}{write · javascript · Easy}
\safemargin{\textbf{Terms for this card}\par\smallskip \textbf{Accumulator:} the running result \code{reduce} passes from one item to the next. Like a running tally on a notepad.}
\noindent{\small\textbf{Covered in:} \hyperref[ch:c06]{Ch.~6} of this guide; documentation: \href{https://developer.mozilla.org/en-US/docs/Web/JavaScript/Reference/Global_Objects/Array/reduce}{MDN — Array.prototype.reduce(), exceptions: TypeError}.}\par
\cardsub{The card}
\textbf{average(nums)}\par
Write \code{average(nums)} using \code{reduce}: the mean of the numbers, or 0 for an empty array.\par
\begin{lstlisting}[language=JavaScript]
function average(nums) {
  // TODO
  return 0;
}
\end{lstlisting}
\cardsub{How to solve it}
\begin{cardrows}
\item[Solve] Return 0 for an empty array; otherwise \code{reduce} from 0 and divide by the length.
\item[Invariants] The sum covers every number exactly once.
\item[Memory] One running sum.
\item[Types] array of numbers in; number out.
\item[Failure] No starting value throws on []; no empty check gives 0 / 0 = NaN.
\item[Tests] [2, 4, 6], [], [5] and [1, 2].
\end{cardrows}
\trap An empty array needs its own answer.\par
\cardsub{Solution}
\begin{lstlisting}[language=JavaScript]
function average(nums) {
  if (nums.length === 0) return 0;
  const total = nums.reduce((sum, n) => sum + n, 0);
  return total / nums.length;
}
\end{lstlisting}
{\small \textit{Verified: passes all 4 test cases on the card in the app's code runner; the starter code fails them.}}\par
\smallskip{\small\textbf{Why:} reduce without a starting value throws on an empty array, and 0 / 0 is NaN, so an empty array needs its own answer.}\par
\cardsub{From the readings}
\readchap{6. How do you work with a list of values?}{\code{map} transforms, \code{filter} keeps, \code{reduce} folds, \code{forEach} only acts; \code{sort} changes the array in place and compares as text unless told otherwise.}
\readq{How do you boil a list down to one value?}
With \code{reduce}. It carries an \textbf{accumulator} through the array: each callback gets the tally so far and the next item and returns the new tally. Give it a starting value (\code{0}, \code{\{\}}, \ldots): without one, an empty array throws a \code{TypeError}. The callback must return the accumulator every time.
\Needspace{6\baselineskip}
\begin{lstlisting}[language=JavaScript]
const add = (sum, n) => sum + n;
console.log([4, 6].reduce(add, 0), [].reduce(add, 0));
try {
  [].reduce(add);
} catch (error) {
  console.log(error.name);
}
\end{lstlisting}
{\small\textit{Prints:}~\code{10 0}\enspace\textcolor{black!45}{|}\enspace \code{TypeError}}\par

\cardend
\cardhead{Card 30 --- Fix doubled(nums)}{fix · javascript · Easy}
\safemargin{\textbf{Terms for this card}\par\smallskip \textbf{Array:} an ordered list of values, numbered from 0. Like a row of numbered lockers.\par\smallskip \textbf{Callback:} a function you pass to another function, which calls it for you. Like instructions you hand a helper: ``do this with each item''.}
\noindent{\small\textbf{Covered in:} \hyperref[ch:c06]{Ch.~6} of this guide; documentation: \href{https://developer.mozilla.org/en-US/docs/Web/JavaScript/Reference/Global_Objects/Array/forEach}{MDN — Array.prototype.forEach()}.}\par
\cardsub{The card}
\textbf{Fix doubled(nums)}\par
\code{doubled(nums)} should return a new array with every number doubled, but it returns \code{undefined}. Fix it.\par
\begin{lstlisting}[language=JavaScript]
function doubled(nums) {
  return nums.forEach((n) => n * 2);
}
\end{lstlisting}
\cardsub{How to solve it}
\begin{cardrows}
\item[Solve] Replace \code{forEach} with \code{map}, which returns the new array.
\item[Invariants] The result has one doubled number per input number.
\item[Memory] One new array.
\item[Types] array in; array out.
\item[Failure] \code{forEach} always returns undefined, whatever its callback returns.
\item[Tests] [1, 2, 3], [] and [-1].
\end{cardrows}
\trap \code{forEach} is for side effects; it never builds an array.\par
\cardsub{Solution}
\begin{lstlisting}[language=JavaScript]
function doubled(nums) {
  return nums.map((n) => n * 2);
}
\end{lstlisting}
{\small \textit{Verified: passes all 3 test cases on the card in the app's code runner; the starter code fails them.}}\par
\smallskip{\small\textbf{Why:} forEach runs the callback for its side effects and always returns undefined. To build a new array from the results, use map.}\par
\cardsub{From the readings}
\readchap{6. How do you work with a list of values?}{\code{map} transforms, \code{filter} keeps, \code{reduce} folds, \code{forEach} only acts; \code{sort} changes the array in place and compares as text unless told otherwise.}
\readq{When is forEach the wrong tool?}
When you want a result. \code{forEach} runs the callback for each item and always returns \code{undefined}; it is for side effects such as adding to a total. To build a new array, use \code{map}.
\Needspace{6\baselineskip}
\begin{lstlisting}[language=JavaScript]
const result = [1, 2].forEach((n) => n * 2);
console.log(result, [1, 2].map((n) => n * 2).join(","));
\end{lstlisting}
{\small\textit{Prints:}~\code{undefined 2,4}}\par

\readq{How do you transform or select items?}
An \textbf{array} has methods that take a \textbf{callback}. \code{map} returns a new array of the same length with each item transformed; \code{filter} returns a new array of the items whose callback returned a truthy value.
\Needspace{6\baselineskip}
\begin{lstlisting}[language=JavaScript]
const prices = [5, 20, 8];
const big = prices.filter((p) => p > 6).map((p) => p * 2);
console.log(big.join(" "), prices.join(" "));
\end{lstlisting}
{\small\textit{Prints:}~\code{40 16 5 20 8}}\par

\cardend
\cardhead{Card 31 --- groupByCategory(items)}{write · javascript · Medium}
\safemargin{\textbf{Terms for this card}\par\smallskip \textbf{Accumulator:} the running result \code{reduce} passes from one item to the next. Like a running tally on a notepad.}
\noindent{\small\textbf{Covered in:} \hyperref[ch:c06]{Ch.~6} of this guide; documentation: \href{https://developer.mozilla.org/en-US/docs/Web/JavaScript/Reference/Global_Objects/Array/reduce}{MDN — Array.prototype.reduce()}.}\par
\cardsub{The card}
\textbf{groupByCategory(items)}\par
Write \code{groupByCategory(items)} with \code{reduce}: given \code{\{ name, category \}} objects, return an object mapping each category to the names in it, in order.\par
\begin{lstlisting}[language=JavaScript]
function groupByCategory(items) {
  // TODO: items.reduce((groups, item) => ..., {})
  return {};
}
\end{lstlisting}
\cardsub{How to solve it}
\begin{cardrows}
\item[Solve] \code{reduce} from \code{\{\}}: create the category's array the first time, push the name, return the accumulator.
\item[Invariants] After each item, the accumulator holds every name seen so far, under its category.
\item[Memory] One object of arrays, built up item by item.
\item[Types] array of objects in; object of string arrays out.
\item[Failure] Assigning \code{[item.name]} each time overwrites earlier names; pushing \code{item} stores the whole object.
\item[Tests] Two items in one category, an empty array, a single item.
\end{cardrows}
\trap Return the accumulator from every call of the callback.\par
\cardsub{Solution}
\begin{lstlisting}[language=JavaScript]
function groupByCategory(items) {
  return items.reduce((groups, item) => {
    if (!groups[item.category]) groups[item.category] = [];
    groups[item.category].push(item.name);
    return groups;
  }, {});
}
\end{lstlisting}
{\small \textit{Verified: passes all 3 test cases on the card in the app's code runner; the starter code fails them.}}\par
\smallskip{\small\textbf{Why:} The accumulator starts as \{\} and must be returned from every callback. Create the category's array the first time you see it, then push.}\par
\cardsub{From the readings}
\readchap{6. How do you work with a list of values?}{\code{map} transforms, \code{filter} keeps, \code{reduce} folds, \code{forEach} only acts; \code{sort} changes the array in place and compares as text unless told otherwise.}
\readq{How do you boil a list down to one value?}
With \code{reduce}. It carries an \textbf{accumulator} through the array: each callback gets the tally so far and the next item and returns the new tally. Give it a starting value (\code{0}, \code{\{\}}, \ldots): without one, an empty array throws a \code{TypeError}. The callback must return the accumulator every time.
\Needspace{6\baselineskip}
\begin{lstlisting}[language=JavaScript]
const add = (sum, n) => sum + n;
console.log([4, 6].reduce(add, 0), [].reduce(add, 0));
try {
  [].reduce(add);
} catch (error) {
  console.log(error.name);
}
\end{lstlisting}
{\small\textit{Prints:}~\code{10 0}\enspace\textcolor{black!45}{|}\enspace \code{TypeError}}\par

\cardend
\cardhead{Card 32 --- What does this print?}{predict · js06-shared-reference}
\noindent{\small\textbf{Covered in:} \hyperref[ch:c06]{Ch.~6} of this guide; documentation: \href{https://developer.mozilla.org/en-US/docs/Web/JavaScript/Reference/Operators/Spread_syntax}{MDN — Spread syntax}.}\par
\cardsub{The card}
What does this print?\par
\begin{lstlisting}[language=JavaScript]
const original = [1, 2, 3];
const alias = original;
const copy = [...original];
alias.push(4);
copy.push(5, 6);
console.log(original.length, alias.length, copy.length);
console.log(original === alias, original === copy);
\end{lstlisting}
\cardsub{How to solve it}
\begin{cardrows}
\item[Solve] \code{alias} is the same array as \code{original}; \code{copy} is a new one. Count each length after the pushes.
\item[Invariants] \code{===} on arrays compares references, not contents.
\item[Memory] Two arrays: original/alias (4 items) and copy (5 items).
\item[Types] arrays, numbers, booleans.
\item[Failure] Answering 3 for original (missing that alias is the same array).
\item[Tests] Run it, then change one line and predict again.
\end{cardrows}
\trap \code{const b = a} copies the reference, not the array.\par
\cardsub{Solution}
{\small\textbf{Prints:}}
\begin{lstlisting}
4 4 5
true false
\end{lstlisting}
{\small\textit{Verified: Node v22.23.4 prints exactly this.}}\par
\smallskip{\small\textbf{Why:} alias = original copies the reference, so both names point at one array. [...original] builds a new array with the same items.}\par
\cardsub{From the readings}
\readchap{6. How do you work with a list of values?}{\code{map} transforms, \code{filter} keeps, \code{reduce} folds, \code{forEach} only acts; \code{sort} changes the array in place and compares as text unless told otherwise.}
\readq{Is a copy of an array a new array?}
Only if you make one. \code{const b = a} copies the reference (\hyperref[ch:c01]{Ch.~1}), so both names point at one array. \code{[...a]} builds a new array with the same items; the items themselves are shared, not copied.
\Needspace{6\baselineskip}
\begin{lstlisting}[language=JavaScript]
const a = [1];
const b = a;
const copy = [...a];
b.push(2);
console.log(a.length, copy.length, a === b);
\end{lstlisting}
{\small\textit{Prints:}~\code{2 1 true}}\par

\cardend
\deckhead{07 · Objects, Destructuring \& Spread}{5}
\begin{corebox}\textbf{Core idea.} An object maps property names to values; destructuring unpacks it, spread copies one level of it, \code{in} asks whether a property exists, and \code{?.} stops at a missing link.\end{corebox}
\begin{sumbox}{Deck summary sheet --- 07 · Objects, Destructuring \& Spread}
Destructuring pulls properties into variables, with renames (\code{\{ a: b \}}) and defaults (\code{\{ a = 1 \}}, used only for undefined). Spread (\code{\{ ...obj \}}) copies one level deep: nested objects are shared. Optional chaining (\code{?.}) stops at null or undefined instead of throwing.\par
\smallskip{\footnotesize Read more: \href{https://developer.mozilla.org/en-US/docs/Web/JavaScript/Reference/Operators/Destructuring}{MDN — Destructuring}; \href{https://developer.mozilla.org/en-US/docs/Web/JavaScript/Reference/Operators/Spread_syntax}{MDN — Spread syntax}; \href{https://developer.mozilla.org/en-US/docs/Web/JavaScript/Reference/Operators/Optional_chaining}{MDN — Optional chaining}}\end{sumbox}
\cardhead{Card 33 --- formatUser(user)}{write · javascript · Easy}
\safemargin{\textbf{Terms for this card}\par\smallskip \textbf{Property:} a named value inside an object, like \code{user.name}. A labelled drawer in the filing cabinet.\par\smallskip \textbf{Destructuring:} syntax that unpacks properties into variables, as in \code{const \{ a, b \} = obj}. Like emptying labelled drawers onto matching labelled trays.}
\noindent{\small\textbf{Covered in:} \hyperref[ch:c07]{Ch.~7} of this guide; documentation: \href{https://developer.mozilla.org/en-US/docs/Web/JavaScript/Reference/Operators/Destructuring}{MDN — Destructuring}.}\par
\cardsub{The card}
\textbf{formatUser(user)}\par
Write \code{formatUser(user)}: given \code{\{ firstName, lastName, role \}}, return \code{\textquotedbl{}LastName, FirstName (ROLE)\textquotedbl{}} with the role in upper case. A missing role means \code{\textquotedbl{}GUEST\textquotedbl{}}. Use destructuring with a default.\par
\begin{lstlisting}[language=JavaScript]
function formatUser(user) {
  // TODO: const { firstName, lastName, role = ... } = user;
  return "";
}
\end{lstlisting}
\cardsub{How to solve it}
\begin{cardrows}
\item[Solve] Destructure with \code{role = \textquotedbl{}GUEST\textquotedbl{}}, then build the string with a template literal and \code{toUpperCase}.
\item[Invariants] A missing or undefined role becomes GUEST.
\item[Memory] Three local variables.
\item[Types] object in; string out.
\item[Failure] No default makes \code{role.toUpperCase()} throw; swapping the names gives the wrong format.
\item[Tests] With a role, without one, with \code{role: undefined}, with mixed case.
\end{cardrows}
\trap A destructuring default applies only to missing or undefined values.\par
\cardsub{Solution}
\begin{lstlisting}[language=JavaScript]
function formatUser(user) {
  const { firstName, lastName, role = "GUEST" } = user;
  return `${lastName}, ${firstName} (${role.toUpperCase()})`;
}
\end{lstlisting}
{\small \textit{Verified: passes all 4 test cases on the card in the app's code runner; the starter code fails them.}}\par
\smallskip{\small\textbf{Why:} A destructuring default applies when the property is missing or undefined (not for null or "").}\par
\cardsub{From the readings}
\readchap{7. How do you hold named data?}{An object maps property names to values; destructuring unpacks it, spread copies one level of it, \code{in} asks whether a property exists, and \code{?.} stops at a missing link.}
\readq{How do you pull properties out into variables?}
An object stores values under names; each is a \textbf{property}. \textbf{Destructuring} copies properties into variables in one line, with renames (\code{\{ a: b \}}) and defaults (\code{\{ role = "guest" \}}). A default is used when the property is missing or \code{undefined}---not when it is \code{null}.
\Needspace{6\baselineskip}
\begin{lstlisting}[language=JavaScript]
const { name, role = "guest" } = { name: "Ada", role: null };
const { title = "none" } = {};
console.log(name, role, title);
\end{lstlisting}
{\small\textit{Prints:}~\code{Ada null none}}\par

\cardend
\cardhead{Card 34 --- What does this print?}{predict · js07-shallow-spread}
\safemargin{\textbf{Terms for this card}\par\smallskip \textbf{Shallow copy:} a copy of the top level only; anything nested is shared, not copied. Like photocopying a folder's index page: the pages it lists are still the originals.}
\noindent{\small\textbf{Covered in:} \hyperref[ch:c07]{Ch.~7} of this guide; documentation: \href{https://developer.mozilla.org/en-US/docs/Web/JavaScript/Reference/Operators/Spread_syntax}{MDN — Spread syntax}.}\par
\cardsub{The card}
What does this print?\par
\begin{lstlisting}[language=JavaScript]
const original = { name: "Ada", skills: ["math"] };
const copy = { ...original };
copy.name = "Grace";
copy.skills.push("code");
console.log(original.name);
console.log(original.skills.join(","));
\end{lstlisting}
\cardsub{How to solve it}
\begin{cardrows}
\item[Solve] Spread copied \code{name} (a string) and the reference to \code{skills} (shared array).
\item[Invariants] Top-level properties are copied; nested objects are shared.
\item[Memory] Two objects, one shared array.
\item[Types] objects, strings, arrays.
\item[Failure] Answering "math" (thinking the array was copied).
\item[Tests] Run it, then change one line and predict again.
\end{cardrows}
\trap Spread copies one level deep.\par
\cardsub{Solution}
{\small\textbf{Prints:}}
\begin{lstlisting}
Ada
math,code
\end{lstlisting}
{\small\textit{Verified: Node v22.23.4 prints exactly this.}}\par
\smallskip{\small\textbf{Why:} Spread copies the top-level properties. copy.name is its own string, but copy.skills is the same array as original.skills. For a deep copy use structuredClone.}\par
\cardsub{From the readings}
\readchap{7. How do you hold named data?}{An object maps property names to values; destructuring unpacks it, spread copies one level of it, \code{in} asks whether a property exists, and \code{?.} stops at a missing link.}
\readq{How do you copy or merge objects?}
With spread: \code{\{ ...a, ...b \}} builds a new object from both; when a name appears twice, the later value wins. So spread the defaults first and the overrides second. Spread makes a \textbf{shallow copy}: a nested array or object is shared between the copy and the original.
\Needspace{6\baselineskip}
\begin{lstlisting}[language=JavaScript]
const base = { theme: "light", tags: ["a"] };
const mine = { ...base, theme: "dark" };
mine.tags.push("b");
console.log(mine.theme, base.theme, base.tags.join(""));
\end{lstlisting}
{\small\textit{Prints:}~\code{dark light ab}}\par

\cardend
\cardhead{Card 35 --- withDefaults(options)}{write · javascript · Easy}
\safemargin{\textbf{Terms for this card}\par\smallskip \textbf{Shallow copy:} a copy of the top level only; anything nested is shared, not copied. Like photocopying a folder's index page: the pages it lists are still the originals.}
\noindent{\small\textbf{Covered in:} \hyperref[ch:c07]{Ch.~7} of this guide; documentation: \href{https://developer.mozilla.org/en-US/docs/Web/JavaScript/Reference/Operators/Spread_syntax}{MDN — Spread syntax}.}\par
\cardsub{The card}
\textbf{withDefaults(options)}\par
Write \code{withDefaults(options)}: return a new object starting from \code{\{ theme: \textquotedbl{}light\textquotedbl{}, fontSize: 14, autosave: true \}}, overridden by whatever \code{options} sets. Use spread.\par
\begin{lstlisting}[language=JavaScript]
function withDefaults(options) {
  const defaults = { theme: "light", fontSize: 14, autosave: true };
  // TODO
  return defaults;
}
\end{lstlisting}
\cardsub{How to solve it}
\begin{cardrows}
\item[Solve] Return \code{\{ ...defaults, ...options \}} so the options override the defaults.
\item[Invariants] Every default key is present unless options sets it.
\item[Memory] One new object.
\item[Types] object in; object out.
\item[Failure] Spreading options first lets the defaults overwrite the user's choices.
\item[Tests] Empty options, one override, several overrides including false.
\end{cardrows}
\trap Later properties win, so spread the overrides last.\par
\cardsub{Solution}
\begin{lstlisting}[language=JavaScript]
function withDefaults(options) {
  const defaults = { theme: "light", fontSize: 14, autosave: true };
  return { ...defaults, ...options };
}
\end{lstlisting}
{\small \textit{Verified: passes all 3 test cases on the card in the app's code runner; the starter code fails them.}}\par
\smallskip{\small\textbf{Why:} Later properties win, so spread the defaults first and the overrides second.}\par
\cardsub{From the readings}
\readchap{7. How do you hold named data?}{An object maps property names to values; destructuring unpacks it, spread copies one level of it, \code{in} asks whether a property exists, and \code{?.} stops at a missing link.}
\readq{How do you copy or merge objects?}
With spread: \code{\{ ...a, ...b \}} builds a new object from both; when a name appears twice, the later value wins. So spread the defaults first and the overrides second. Spread makes a \textbf{shallow copy}: a nested array or object is shared between the copy and the original.
\Needspace{6\baselineskip}
\begin{lstlisting}[language=JavaScript]
const base = { theme: "light", tags: ["a"] };
const mine = { ...base, theme: "dark" };
mine.tags.push("b");
console.log(mine.theme, base.theme, base.tags.join(""));
\end{lstlisting}
{\small\textit{Prints:}~\code{dark light ab}}\par

\cardend
\cardhead{Card 36 --- pick(obj, keys)}{write · javascript · Easy}
\noindent{\small\textbf{Covered in:} \hyperref[ch:c07]{Ch.~7} of this guide; documentation: \href{https://developer.mozilla.org/en-US/docs/Web/JavaScript/Reference/Operators/in}{MDN — in}.}\par
\cardsub{The card}
\textbf{pick(obj, keys)}\par
Write \code{pick(obj, keys)}: a new object with only the listed keys that \code{obj} actually has. Keep keys whose value is \code{null}.\par
\begin{lstlisting}[language=JavaScript]
function pick(obj, keys) {
  const result = {};
  // TODO
  return result;
}
\end{lstlisting}
\cardsub{How to solve it}
\begin{cardrows}
\item[Solve] For each key, copy it only \code{if (key in obj)}.
\item[Invariants] The result has only keys that exist on \code{obj}, with their values, null included.
\item[Memory] One new object.
\item[Types] object and array of strings in; object out.
\item[Failure] Copying blindly adds \code{undefined} keys; testing the value drops null, 0 and "".
\item[Tests] Existing keys, a missing key, no keys, and a key whose value is null.
\end{cardrows}
\trap Ask whether the key exists, not whether its value is truthy.\par
\cardsub{Solution}
\begin{lstlisting}[language=JavaScript]
function pick(obj, keys) {
  const result = {};
  for (const key of keys) {
    if (key in obj) result[key] = obj[key];
  }
  return result;
}
\end{lstlisting}
{\small \textit{Verified: passes all 4 test cases on the card in the app's code runner; the starter code fails them.}}\par
\smallskip{\small\textbf{Why:} \code{key in obj} asks whether the property exists. Testing its value instead drops real properties that hold null, 0 or "", and copying blindly adds undefined keys.}\par
\cardsub{From the readings}
\readchap{7. How do you hold named data?}{An object maps property names to values; destructuring unpacks it, spread copies one level of it, \code{in} asks whether a property exists, and \code{?.} stops at a missing link.}
\readq{How do you tell a missing property from an empty one?}
Ask \code{key in obj}: it is \code{true} when the property exists, whatever its value. Testing the value instead (\code{if (obj[key])}) wrongly treats \code{null}, \code{0} and \code{""} as missing.
\Needspace{6\baselineskip}
\begin{lstlisting}[language=JavaScript]
const box = { count: 0 };
console.log("count" in box, Boolean(box.count), "size" in box);
\end{lstlisting}
{\small\textit{Prints:}~\code{true false false}}\par

\cardend
\cardhead{Card 37 --- What does this print?}{predict · js07-optional-chaining}
\safemargin{\textbf{Terms for this card}\par\smallskip \textbf{Optional chaining (?.):} reads a property but returns \code{undefined} instead of throwing when the value before it is \code{null} or \code{undefined}. Like a hallway with ``if the door is locked, turn back'' signs.}
\noindent{\small\textbf{Covered in:} \hyperref[ch:c07]{Ch.~7} of this guide; documentation: \href{https://developer.mozilla.org/en-US/docs/Web/JavaScript/Reference/Operators/Optional_chaining}{MDN — Optional chaining (?.)}.}\par
\cardsub{The card}
What does this print?\par
\begin{lstlisting}[language=JavaScript]
const user = { profile: null };
console.log(user.profile?.email);
console.log(user.settings?.theme ?? "default");
try {
  console.log(user.profile.email);
} catch (error) {
  console.log(error.name);
}
\end{lstlisting}
\cardsub{How to solve it}
\begin{cardrows}
\item[Solve] \code{?.} stops at \code{null} (undefined); \code{?.} plus \code{??} gives the default; plain \code{.} on null throws.
\item[Invariants] \code{?.} never throws on null or undefined.
\item[Memory] One object whose \code{profile} is null.
\item[Types] undefined, string, TypeError.
\item[Failure] Expecting the third line to print undefined instead of TypeError.
\item[Tests] Run it, then change one line and predict again.
\end{cardrows}
\trap Only \code{?.} is safe; one plain dot after a null throws.\par
\cardsub{Solution}
{\small\textbf{Prints:}}
\begin{lstlisting}
undefined
default
TypeError
\end{lstlisting}
{\small\textit{Verified: Node v22.23.4 prints exactly this.}}\par
\smallskip{\small\textbf{Why:} ?. returns undefined as soon as the value on its left is null or undefined. Plain . on null throws a TypeError.}\par
\cardsub{From the readings}
\readchap{7. How do you hold named data?}{An object maps property names to values; destructuring unpacks it, spread copies one level of it, \code{in} asks whether a property exists, and \code{?.} stops at a missing link.}
\readq{How do you read a property that might not be there?}
With \textbf{optional chaining}. \code{a?.b} is \code{undefined} when \code{a} is \code{null} or \code{undefined}; plain \code{a.b} would throw a \code{TypeError}. Combine it with \code{??} to supply a fallback.
\Needspace{6\baselineskip}
\begin{lstlisting}[language=JavaScript]
const user = { profile: null };
console.log(user.profile?.email ?? "no email");
\end{lstlisting}
{\small\textit{Prints:}~\code{no email}}\par

\cardend
\deckhead{08 · Map \& Set}{4}
\begin{corebox}\textbf{Core idea.} A \code{Set} keeps one of each value; a \code{Map} pairs any key with a value in insertion order; a plain object turns keys into strings and inherits keys you never added.\end{corebox}
\begin{sumbox}{Deck summary sheet --- 08 · Map \& Set}
A \code{Set} holds unique values; a \code{Map} holds key $\rightarrow$ value pairs where keys can be any type and stay in insertion order. Plain objects turn keys into strings and inherit keys such as "constructor", which makes them risky as dictionaries.\par
\smallskip{\footnotesize Read more: \href{https://developer.mozilla.org/en-US/docs/Web/JavaScript/Guide/Keyed_collections}{MDN — Keyed collections}}\end{sumbox}
\cardhead{Card 38 --- unique(values)}{write · javascript · Easy}
\safemargin{\textbf{Terms for this card}\par\smallskip \textbf{Set:} a collection in which each value appears at most once, kept in the order first added. Like a guest list where writing a name twice changes nothing.}
\noindent{\small\textbf{Covered in:} \hyperref[ch:c08]{Ch.~8} of this guide; documentation: \href{https://developer.mozilla.org/en-US/docs/Web/JavaScript/Reference/Global_Objects/Set}{MDN — Set}.}\par
\cardsub{The card}
\textbf{unique(values)}\par
Write \code{unique(values)} with a \code{Set}: the values without duplicates, in first-seen order.\par
\begin{lstlisting}[language=JavaScript]
function unique(values) {
  // TODO
  return values;
}
\end{lstlisting}
\cardsub{How to solve it}
\begin{cardrows}
\item[Solve] Return \code{[...new Set(values)]}.
\item[Invariants] First-seen order is kept; each value appears once.
\item[Memory] One Set, then one array.
\item[Types] array in; array out.
\item[Failure] Sorting the result changes the order.
\item[Tests] Repeats, empty, "a" vs "A", and 1 vs "1".
\end{cardrows}
\trap A Set compares with ===, so 1 and "1" both stay.\par
\cardsub{Solution}
\begin{lstlisting}[language=JavaScript]
function unique(values) {
  return [...new Set(values)];
}
\end{lstlisting}
{\small \textit{Verified: passes all 4 test cases on the card in the app's code runner; the starter code fails them.}}\par
\smallskip{\small\textbf{Why:} A Set keeps one of each value, compared with ===, in insertion order. Spread it back into an array.}\par
\cardsub{From the readings}
\readchap{8. Which collection fits: object, Map or Set?}{A \code{Set} keeps one of each value; a \code{Map} pairs any key with a value in insertion order; a plain object turns keys into strings and inherits keys you never added.}
\readq{How do you remove duplicates?}
Put the values in a \textbf{Set}. A value can occur in a Set only once, compared with \code{===}, and the Set remembers the order values were first added. Spread it back into an array with \code{[...set]}.
\Needspace{6\baselineskip}
\begin{lstlisting}[language=JavaScript]
const tags = new Set(["js", "ts", "js"]);
console.log(tags.size, [...tags].join(","));
\end{lstlisting}
{\small\textit{Prints:}~\code{2 js,ts}}\par

\cardend
\cardhead{Card 39 --- mostFrequentWord(text)}{write · javascript · Medium}
\safemargin{\textbf{Terms for this card}\par\smallskip \textbf{Map:} a collection of key--value pairs where the key can be any value and pairs stay in insertion order. Like a coat check: any ticket can claim exactly one coat.}
\noindent{\small\textbf{Covered in:} \hyperref[ch:c08]{Ch.~8} of this guide; documentation: \href{https://developer.mozilla.org/en-US/docs/Web/JavaScript/Reference/Global_Objects/Map}{MDN — Map}.}\par
\cardsub{The card}
\textbf{mostFrequentWord(text)}\par
Write \code{mostFrequentWord(text)} with a \code{Map}: the most common word, lower-cased. On a tie, return the word that appeared first. Empty text gives \code{\textquotedbl{}\textquotedbl{}}.\par
\begin{lstlisting}[language=JavaScript]
function mostFrequentWord(text) {
  const counts = new Map();
  // TODO
  return "";
}
\end{lstlisting}
\cardsub{How to solve it}
\begin{cardrows}
\item[Solve] Count lower-cased words in a Map, then walk it keeping the first word with a strictly larger count.
\item[Invariants] The Map iterates in insertion order, so ties keep the earliest word.
\item[Memory] A Map of word to count, plus the best word so far.
\item[Types] string in; string out.
\item[Failure] \code{\textgreater{}=} picks the last word on a tie; not lower-casing splits "Go" and "go".
\item[Tests] A clear winner, mixed case, a tie, and empty text.
\end{cardrows}
\trap Use \code{\textgreater{}} so a tie keeps the first word.\par
\cardsub{Solution}
\begin{lstlisting}[language=JavaScript]
function mostFrequentWord(text) {
  const counts = new Map();
  for (const word of text.toLowerCase().split(/\s+/).filter(Boolean)) {
    counts.set(word, (counts.get(word) ?? 0) + 1);
  }
  let best = "";
  let bestCount = 0;
  for (const [word, count] of counts) {
    if (count > bestCount) {
      best = word;
      bestCount = count;
    }
  }
  return best;
}
\end{lstlisting}
{\small \textit{Verified: passes all 5 test cases on the card in the app's code runner; the starter code fails them.}}\par
\smallskip{\small\textbf{Why:} A Map iterates in insertion order, so a strict > keeps the first word on a tie. \code{counts.get(word) ?? 0} starts a new word at zero.}\par
\cardsub{From the readings}
\readchap{8. Which collection fits: object, Map or Set?}{A \code{Set} keeps one of each value; a \code{Map} pairs any key with a value in insertion order; a plain object turns keys into strings and inherits keys you never added.}
\readq{When is a Map better than an object?}
When you count or group by key. A \textbf{Map} stores key--value pairs: \code{set(key, value)}, \code{get(key)}, \code{has(key)}, \code{size}. Keys keep their type, so the number \code{1} and the string \code{"1"} are different keys; an object would turn both into \code{"1"}. A Map also iterates in insertion order, which settles ties ``first come, first served''. \code{counts.get(word) ?? 0} starts a new key at zero.
\Needspace{6\baselineskip}
\begin{lstlisting}[language=JavaScript]
const counts = new Map();
for (const w of ["a", "b", "a"]) {
  counts.set(w, (counts.get(w) ?? 0) + 1);
}
console.log(counts.get("a"), counts.size);
\end{lstlisting}
{\small\textit{Prints:}~\code{2 2}}\par

\cardend
\cardhead{Card 40 --- What does this print?}{predict · js08-map-keys}
\safemargin{\textbf{Terms for this card}\par\smallskip \textbf{Map:} a collection of key--value pairs where the key can be any value and pairs stay in insertion order. Like a coat check: any ticket can claim exactly one coat.}
\noindent{\small\textbf{Covered in:} \hyperref[ch:c08]{Ch.~8} of this guide; documentation: \href{https://developer.mozilla.org/en-US/docs/Web/JavaScript/Guide/Keyed_collections}{MDN — Keyed collections}.}\par
\cardsub{The card}
What does this print?\par
\begin{lstlisting}[language=JavaScript]
const obj = {};
obj[1] = "number key";
obj["1"] = "string key";
const map = new Map();
map.set(1, "number key");
map.set("1", "string key");
console.log(Object.keys(obj).length, map.size);
console.log(map.get(1));
\end{lstlisting}
\cardsub{How to solve it}
\begin{cardrows}
\item[Solve] The object turns 1 and "1" into one key; the Map keeps two.
\item[Invariants] Map keys keep their type.
\item[Memory] An object with one key; a Map with two entries.
\item[Types] numbers and strings as keys.
\item[Failure] Answering "2 2" or "string key".
\item[Tests] Run it, then change one line and predict again.
\end{cardrows}
\trap Object keys are strings; Map keys can be anything.\par
\cardsub{Solution}
{\small\textbf{Prints:}}
\begin{lstlisting}
1 2
number key
\end{lstlisting}
{\small\textit{Verified: Node v22.23.4 prints exactly this.}}\par
\smallskip{\small\textbf{Why:} Object keys are strings, so 1 and "1" are the same key and the second write replaces the first. A Map keeps the number 1 and the string "1" apart.}\par
\cardsub{From the readings}
\readchap{8. Which collection fits: object, Map or Set?}{A \code{Set} keeps one of each value; a \code{Map} pairs any key with a value in insertion order; a plain object turns keys into strings and inherits keys you never added.}
\readq{When is a Map better than an object?}
When you count or group by key. A \textbf{Map} stores key--value pairs: \code{set(key, value)}, \code{get(key)}, \code{has(key)}, \code{size}. Keys keep their type, so the number \code{1} and the string \code{"1"} are different keys; an object would turn both into \code{"1"}. A Map also iterates in insertion order, which settles ties ``first come, first served''. \code{counts.get(word) ?? 0} starts a new key at zero.
\Needspace{6\baselineskip}
\begin{lstlisting}[language=JavaScript]
const counts = new Map();
for (const w of ["a", "b", "a"]) {
  counts.set(w, (counts.get(w) ?? 0) + 1);
}
console.log(counts.get("a"), counts.size);
\end{lstlisting}
{\small\textit{Prints:}~\code{2 2}}\par

\cardend
\cardhead{Card 41 --- Fix groupByCategory for any category}{fix · javascript · Medium}
\safemargin{\textbf{Terms for this card}\par\smallskip \textbf{Prototype chain:} the series of objects an object borrows properties from when it lacks its own. Like a family recipe book you can read from even though it isn't yours.\par\smallskip \textbf{Map:} a collection of key--value pairs where the key can be any value and pairs stay in insertion order. Like a coat check: any ticket can claim exactly one coat.}
\noindent{\small\textbf{Covered in:} \hyperref[ch:c08]{Ch.~8} of this guide; documentation: \href{https://developer.mozilla.org/en-US/docs/Web/JavaScript/Reference/Global_Objects/Object/fromEntries}{MDN — Object.fromEntries()}.}\par
\cardsub{The card}
\textbf{Fix groupByCategory for any category}\par
This \code{groupByCategory(items)} crashes when a category is called \code{\textquotedbl{}constructor\textquotedbl{}} or \code{\textquotedbl{}toString\textquotedbl{}}. Fix it by grouping in a \code{Map}, then return a plain object with \code{Object.fromEntries}.\par
\begin{lstlisting}[language=JavaScript]
function groupByCategory(items) {
  const groups = {};
  for (const item of items) {
    if (!groups[item.category]) groups[item.category] = [];
    groups[item.category].push(item.name);
  }
  return groups;
}
\end{lstlisting}
\cardsub{How to solve it}
\begin{cardrows}
\item[Solve] Group in a \code{Map} (\code{has}, \code{set}, \code{get(...).push}), then \code{Object.fromEntries(groups)}.
\item[Invariants] Every category starts as an empty array the first time it is seen.
\item[Memory] A Map of category to names, then a plain object.
\item[Types] array of objects in; object out.
\item[Failure] \code{\{\}[\textquotedbl{}constructor\textquotedbl{}]} already exists (a function), so the plain-object version pushes onto it and throws.
\item[Tests] "constructor", "toString", a normal category, and an empty list.
\end{cardrows}
\trap Plain objects inherit keys you never added.\par
\cardsub{Solution}
\begin{lstlisting}[language=JavaScript]
function groupByCategory(items) {
  const groups = new Map();
  for (const item of items) {
    if (!groups.has(item.category)) groups.set(item.category, []);
    groups.get(item.category).push(item.name);
  }
  return Object.fromEntries(groups);
}
\end{lstlisting}
{\small \textit{Verified: passes all 4 test cases on the card in the app's code runner; the starter code fails them.}}\par
\smallskip{\small\textbf{Why:} Every plain object inherits properties like constructor and toString, so \code{groups[\textquotedbl{}constructor\textquotedbl{}]} already exists and isn't an array. A Map starts truly empty.}\par
\cardsub{From the readings}
\readchap{8. Which collection fits: object, Map or Set?}{A \code{Set} keeps one of each value; a \code{Map} pairs any key with a value in insertion order; a plain object turns keys into strings and inherits keys you never added.}
\readq{Why can a plain object already contain "constructor"?}
Because every plain object inherits properties through its \textbf{prototype chain}. \code{\{\}["constructor"]} is a function, not \code{undefined}, so ``create the group if it doesn't exist yet'' skips it and the code then calls \code{push} on a function. A Map starts truly empty; \code{Object.fromEntries(map)} turns it into a plain object at the end.
\Needspace{6\baselineskip}
\begin{lstlisting}[language=JavaScript]
const groups = {};
console.log(typeof groups["constructor"]);
console.log(new Map().has("constructor"));
\end{lstlisting}
{\small\textit{Prints:}~\code{function}\enspace\textcolor{black!45}{|}\enspace \code{false}}\par

\readq{When is a Map better than an object?}
When you count or group by key. A \textbf{Map} stores key--value pairs: \code{set(key, value)}, \code{get(key)}, \code{has(key)}, \code{size}. Keys keep their type, so the number \code{1} and the string \code{"1"} are different keys; an object would turn both into \code{"1"}. A Map also iterates in insertion order, which settles ties ``first come, first served''. \code{counts.get(word) ?? 0} starts a new key at zero.
\Needspace{6\baselineskip}
\begin{lstlisting}[language=JavaScript]
const counts = new Map();
for (const w of ["a", "b", "a"]) {
  counts.set(w, (counts.get(w) ?? 0) + 1);
}
console.log(counts.get("a"), counts.size);
\end{lstlisting}
{\small\textit{Prints:}~\code{2 2}}\par

\cardend
\deckhead{09 · Recursion}{4}
\begin{corebox}\textbf{Core idea.} A recursive function answers the smallest case directly and hands a smaller version of the problem to itself.\end{corebox}
\begin{sumbox}{Deck summary sheet --- 09 · Recursion}
A recursive function calls itself on a smaller piece of the problem and stops at a base case. It fits nested data — arrays inside arrays, folders inside folders — whose depth you don't know in advance.\par
\smallskip{\footnotesize Read more: \href{https://developer.mozilla.org/en-US/docs/Web/JavaScript/Guide/Functions\#recursion}{MDN — Functions: recursion}}\end{sumbox}
\cardhead{Card 42 --- factorial(n)}{write · javascript · Easy}
\safemargin{\textbf{Terms for this card}\par\smallskip \textbf{Recursion:} a function calling itself on a smaller part of the problem. Like nesting dolls: open one, find a smaller one, until the smallest.\par\smallskip \textbf{Base case:} the input a recursive function answers without calling itself. The smallest doll, which doesn't open.}
\noindent{\small\textbf{Covered in:} \hyperref[ch:c09]{Ch.~9} of this guide; documentation: \href{https://developer.mozilla.org/en-US/docs/Glossary/Recursion}{MDN — Recursion}.}\par
\cardsub{The card}
\textbf{factorial(n)}\par
Write \code{factorial(n)} recursively: n $\times$ (n $-$ 1) $\times$ \ldots{} $\times$ 1, with \code{factorial(0)} = 1.\par
\begin{lstlisting}[language=JavaScript]
function factorial(n) {
  // TODO: base case, then n * factorial(n - 1)
  return n;
}
\end{lstlisting}
\cardsub{How to solve it}
\begin{cardrows}
\item[Solve] Return 1 when \code{n \textless{}= 1}; otherwise \code{n * factorial(n - 1)}.
\item[Invariants] Every call is on a smaller n until the base case.
\item[Memory] One stack frame per pending call.
\item[Types] number in; number out.
\item[Failure] A base case of \code{n === 1} never stops for 0.
\item[Tests] 0, 1, 5 and 10.
\end{cardrows}
\trap The base case must cover 0.\par
\cardsub{Solution}
\begin{lstlisting}[language=JavaScript]
function factorial(n) {
  if (n <= 1) return 1;
  return n * factorial(n - 1);
}
\end{lstlisting}
{\small \textit{Verified: passes all 4 test cases on the card in the app's code runner; the starter code fails them.}}\par
\smallskip{\small\textbf{Why:} Without a base case that covers 0, factorial(0) recurses forever and overflows the stack.}\par
\cardsub{From the readings}
\readchap{9. How can a function solve a problem by calling itself?}{A recursive function answers the smallest case directly and hands a smaller version of the problem to itself.}
\readq{What makes recursion stop?}
A base case. \textbf{Recursion} has two kinds of input: a \textbf{base case}, answered directly, and a recursive case, which calls the function again on something smaller. Without a base case that every input eventually reaches, the calls never stop and JavaScript throws a \code{RangeError} when its call stack is full.
\Needspace{6\baselineskip}
\begin{lstlisting}[language=JavaScript]
function forever(n) {
  return forever(n + 1);
}
try {
  forever(0);
} catch (error) {
  console.log(error.name);
}
\end{lstlisting}
{\small\textit{Prints:}~\code{RangeError}}\par

\cardend
\cardhead{Card 43 --- sumNested(items)}{write · javascript · Easy}
\noindent{\small\textbf{Covered in:} \hyperref[ch:c09]{Ch.~9} of this guide; documentation: \href{https://developer.mozilla.org/en-US/docs/Glossary/Recursion}{MDN — Recursion}.}\par
\cardsub{The card}
\textbf{sumNested(items)}\par
Write \code{sumNested(items)}: add every number in an array whose items can be arrays, nested to any depth.\par
\begin{lstlisting}[language=JavaScript]
function sumNested(items) {
  let total = 0;
  // TODO
  return total;
}
\end{lstlisting}
\cardsub{How to solve it}
\begin{cardrows}
\item[Solve] For each item, add \code{sumNested(item)} if it is an array, else the item itself.
\item[Invariants] Each call sums one array at one level, delegating deeper levels.
\item[Memory] One running total per pending call.
\item[Types] nested array of numbers in; number out.
\item[Failure] \code{flat()} removes only one level, so deeper arrays end up as strings in the sum.
\item[Tests] Mixed depths, empty, nested empty, a single number.
\end{cardrows}
\trap One level of flattening is not enough.\par
\cardsub{Solution}
\begin{lstlisting}[language=JavaScript]
function sumNested(items) {
  let total = 0;
  for (const item of items) {
    total += Array.isArray(item) ? sumNested(item) : item;
  }
  return total;
}
\end{lstlisting}
{\small \textit{Verified: passes all 4 test cases on the card in the app's code runner; the starter code fails them.}}\par
\smallskip{\small\textbf{Why:} Each array is the same problem, smaller. \code{flat()} only removes one level of nesting.}\par
\cardsub{From the readings}
\readchap{9. How can a function solve a problem by calling itself?}{A recursive function answers the smallest case directly and hands a smaller version of the problem to itself.}
\readq{Why does recursion fit nested data?}
Because each nested piece is the same problem, smaller: an array inside an array, a folder inside a folder. The function handles one level and calls itself for each nested piece, so it works at any depth. \code{flat()} removes only one level of nesting.
\Needspace{6\baselineskip}
\begin{lstlisting}[language=JavaScript]
function depth(item) {
  if (!Array.isArray(item)) return 0;
  return 1 + Math.max(0, ...item.map(depth));
}
console.log(depth([1, [2, [3]]]), depth(7));
\end{lstlisting}
{\small\textit{Prints:}~\code{3 0}}\par

\cardend
\cardhead{Card 44 --- countFiles(folder)}{write · javascript · Easy}
\noindent{\small\textbf{Covered in:} \hyperref[ch:c09]{Ch.~9} of this guide; documentation: \href{https://developer.mozilla.org/en-US/docs/Web/JavaScript/Guide/Functions}{MDN — Functions}.}\par
\cardsub{The card}
\textbf{countFiles(folder)}\par
A folder is \code{\{ name, files: string[], folders: Folder[] \}}. Write \code{countFiles(folder)}: the number of files in it and every folder below it.\par
\begin{lstlisting}[language=JavaScript]
function countFiles(folder) {
  // TODO
  return folder.files.length;
}
\end{lstlisting}
\cardsub{How to solve it}
\begin{cardrows}
\item[Solve] Start with this folder's file count and add \code{countFiles(child)} for every sub-folder.
\item[Invariants] A folder with no sub-folders just returns its own count.
\item[Memory] One running total per pending call.
\item[Types] folder object in; number out.
\item[Failure] Adding only the children's own files misses deeper levels.
\item[Tests] A single folder, and a tree three levels deep with an empty folder.
\end{cardrows}
\trap Recurse into children; don't just look at them.\par
\cardsub{Solution}
\begin{lstlisting}[language=JavaScript]
function countFiles(folder) {
  let total = folder.files.length;
  for (const child of folder.folders) {
    total += countFiles(child);
  }
  return total;
}
\end{lstlisting}
{\small \textit{Verified: passes all 2 test cases on the card in the app's code runner; the starter code fails them.}}\par
\smallskip{\small\textbf{Why:} The base case is a folder with no sub-folders: the loop doesn't run and it returns its own file count.}\par
\cardsub{From the readings}
\readchap{9. How can a function solve a problem by calling itself?}{A recursive function answers the smallest case directly and hands a smaller version of the problem to itself.}
\readq{Why does recursion fit nested data?}
Because each nested piece is the same problem, smaller: an array inside an array, a folder inside a folder. The function handles one level and calls itself for each nested piece, so it works at any depth. \code{flat()} removes only one level of nesting.
\Needspace{6\baselineskip}
\begin{lstlisting}[language=JavaScript]
function depth(item) {
  if (!Array.isArray(item)) return 0;
  return 1 + Math.max(0, ...item.map(depth));
}
console.log(depth([1, [2, [3]]]), depth(7));
\end{lstlisting}
{\small\textit{Prints:}~\code{3 0}}\par

\cardend
\cardhead{Card 45 --- What does this print?}{predict · js09-trace}
\safemargin{\textbf{Terms for this card}\par\smallskip \textbf{Call stack:} the list of function calls that have started but not yet returned. Like a stack of plates: each call adds one, each return removes the top one.}
\noindent{\small\textbf{Covered in:} \hyperref[ch:c09]{Ch.~9} of this guide; documentation: \href{https://developer.mozilla.org/en-US/docs/Web/JavaScript/Guide/Functions}{MDN — Functions}.}\par
\cardsub{The card}
What does this print?\par
\begin{lstlisting}[language=JavaScript]
function countdown(n) {
  if (n === 0) return "liftoff";
  return n + " " + countdown(n - 1);
}
console.log(countdown(3));
\end{lstlisting}
\cardsub{How to solve it}
\begin{cardrows}
\item[Solve] Expand: 3 + countdown(2) $\rightarrow$ 2 + countdown(1) $\rightarrow$ 1 + countdown(0) $\rightarrow$ "liftoff"; join on the way back.
\item[Invariants] Each call returns a string that starts with its own number.
\item[Memory] Up to four calls waiting on the stack.
\item[Types] numbers joined into a string.
\item[Failure] Reversing the order ("1 2 3 liftoff").
\item[Tests] Run it, then change one line and predict again.
\end{cardrows}
\trap Each call puts its number in front of the smaller call's result.\par
\cardsub{Solution}
{\small\textbf{Prints:}}
\begin{lstlisting}
3 2 1 liftoff
\end{lstlisting}
{\small\textit{Verified: Node v22.23.4 prints exactly this.}}\par
\smallskip{\small\textbf{Why:} Each call waits for the smaller call to return, then puts its own number in front.}\par
\cardsub{From the readings}
\readchap{9. How can a function solve a problem by calling itself?}{A recursive function answers the smallest case directly and hands a smaller version of the problem to itself.}
\readq{In what order do recursive calls finish?}
Last started, first finished. Each call waits on the \textbf{call stack} until the smaller call returns, then uses that result. So results are put together on the way back out.
\Needspace{6\baselineskip}
\begin{lstlisting}[language=JavaScript]
function nest(n) {
  if (n === 0) return "core";
  const inner = nest(n - 1);
  console.log("back in", n);
  return `(${inner})`;
}
console.log(nest(2));
\end{lstlisting}
{\small\textit{Prints:}~\code{back in 1}\enspace\textcolor{black!45}{|}\enspace \code{back in 2}\enspace\textcolor{black!45}{|}\enspace \code{((core))}}\par

\cardend
\deckhead{10 · Errors: throw, try, catch, finally}{4}
\begin{corebox}\textbf{Core idea.} \code{throw} stops the function and travels up the calls until a \code{catch} handles it; catch only where you can do something useful, and \code{finally} runs whichever way the block exits.\end{corebox}
\begin{sumbox}{Deck summary sheet --- 10 · Errors: throw, try, catch, finally}
\code{throw} stops the function and travels up the call stack until a \code{try...catch} handles it. Catch only where you can do something useful; catching and returning a default hides the failure. \code{finally} runs whichever way the block exits.\par
\smallskip{\footnotesize Read more: \href{https://developer.mozilla.org/en-US/docs/Web/JavaScript/Guide/Control_flow_and_error_handling}{MDN — Control flow and error handling}}\end{sumbox}
\cardhead{Card 46 --- parseAge(text)}{write · javascript · Easy}
\safemargin{\textbf{Terms for this card}\par\smallskip \textbf{Exception:} an error object thrown to signal that something went wrong. Like a fire alarm: work stops and everyone heads for the nearest exit.}
\noindent{\small\textbf{Covered in:} \hyperref[ch:c10]{Ch.~10} of this guide; documentation: \href{https://developer.mozilla.org/en-US/docs/Web/JavaScript/Guide/Control_flow_and_error_handling}{MDN — Control flow and error handling}.}\par
\cardsub{The card}
\textbf{parseAge(text)}\par
Write \code{parseAge(text)}: return the age as a number. For anything that isn't a whole number from 0 to 150 (including empty text), throw a \code{RangeError} with the message \code{\textquotedbl{}age must be a whole number from 0 to 150\textquotedbl{}}.\par
\begin{lstlisting}[language=JavaScript]
function parseAge(text) {
  // TODO: validate, and throw new RangeError(...) for bad input
  return Number(text);
}
\end{lstlisting}
\cardsub{How to solve it}
\begin{cardrows}
\item[Solve] Convert with \code{Number}; throw a \code{RangeError} for blank text, non-integers, below 0 or above 150.
\item[Invariants] Only whole numbers from 0 to 150 are returned.
\item[Memory] One local number.
\item[Types] string in; number out, or a RangeError.
\item[Failure] Throwing a plain \code{Error}; skipping the blank check ("" becomes 0).
\item[Tests] "42", "0", "150", "151", "-1", "4.5", "abc", "" and "   ".
\end{cardrows}
\trap Blank text converts to 0 and passes unless checked.\par
\cardsub{Solution}
\begin{lstlisting}[language=JavaScript]
function parseAge(text) {
  const age = Number(text);
  if (text.trim() === "" || !Number.isInteger(age) || age < 0 || age > 150) {
    throw new RangeError("age must be a whole number from 0 to 150");
  }
  return age;
}
\end{lstlisting}
{\small \textit{Verified: passes all 9 test cases on the card in the app's code runner; the starter code fails them.}}\par
\smallskip{\small\textbf{Why:} Throw the most specific built-in error (RangeError for out-of-range values) so callers can tell failures apart. Number("") is 0, so empty input needs its own check.}\par
\cardsub{From the readings}
\readchap{10. What should a program do when something goes wrong?}{\code{throw} stops the function and travels up the calls until a \code{catch} handles it; catch only where you can do something useful, and \code{finally} runs whichever way the block exits.}
\readq{What happens when you throw?}
The function stops and an \textbf{exception} travels up the call stack to the first \code{catch} that waits for it; with no \code{catch}, the program ends. Throw the most specific built-in error: \code{RangeError} for a value outside the allowed range. An error has a \code{name} and a \code{message}.
\Needspace{6\baselineskip}
\begin{lstlisting}[language=JavaScript]
function checkAge(age) {
  if (age < 0) throw new RangeError("negative age");
  return age;
}
try {
  checkAge(-1);
} catch (error) {
  console.log(`${error.name}: ${error.message}`);
}
\end{lstlisting}
{\small\textit{Prints:}~\code{RangeError: negative age}}\par

\cardend
\cardhead{Card 47 --- What does this print?}{predict · js10-finally}
\noindent{\small\textbf{Covered in:} \hyperref[ch:c10]{Ch.~10} of this guide; documentation: \href{https://developer.mozilla.org/en-US/docs/Web/JavaScript/Reference/Statements/try...catch}{MDN — try...catch}.}\par
\cardsub{The card}
What does this print?\par
\begin{lstlisting}[language=JavaScript]
function load() {
  try {
    console.log("try");
    return "result";
  } finally {
    console.log("finally");
  }
}
console.log(load());
\end{lstlisting}
\cardsub{How to solve it}
\begin{cardrows}
\item[Solve] "try" prints, the return is prepared, \code{finally} prints, then the caller prints the result.
\item[Invariants] \code{finally} runs before control leaves the try statement.
\item[Memory] The pending return value waits while finally runs.
\item[Types] strings.
\item[Failure] Putting "result" before "finally".
\item[Tests] Run it, then change one line and predict again.
\end{cardrows}
\trap A return inside try still runs finally first.\par
\cardsub{Solution}
{\small\textbf{Prints:}}
\begin{lstlisting}
try
finally
result
\end{lstlisting}
{\small\textit{Verified: Node v22.23.4 prints exactly this.}}\par
\smallskip{\small\textbf{Why:} finally runs on every exit from the try block, return included, before the caller gets the value.}\par
\cardsub{From the readings}
\readchap{10. What should a program do when something goes wrong?}{\code{throw} stops the function and travels up the calls until a \code{catch} handles it; catch only where you can do something useful, and \code{finally} runs whichever way the block exits.}
\readq{What does finally guarantee?}
A \code{finally} block always runs before control leaves the \code{try} construct---after a normal finish, an error, or a \code{return}.
\Needspace{6\baselineskip}
\begin{lstlisting}[language=JavaScript]
function f() {
  try {
    return "value";
  } finally {
    console.log("cleanup");
  }
}
console.log(f());
\end{lstlisting}
{\small\textit{Prints:}~\code{cleanup}\enspace\textcolor{black!45}{|}\enspace \code{value}}\par

\cardend
\cardhead{Card 48 --- Fix readConfig(json)}{fix · javascript · Medium}
\safemargin{\textbf{Terms for this card}\par\smallskip \textbf{JSON:} a text format for data (objects, arrays, strings, numbers, booleans, null) that any program can read. Like a shipping label written in a standard layout.}
\noindent{\small\textbf{Covered in:} \hyperref[ch:c10]{Ch.~10}, \hyperref[ch:c11]{Ch.~11} of this guide; documentation: \href{https://developer.mozilla.org/en-US/docs/Web/JavaScript/Reference/Global_Objects/JSON/parse}{MDN — JSON.parse()}.}\par
\cardsub{The card}
\textbf{Fix readConfig(json)}\par
\code{readConfig(json)} hides broken input by returning \code{\{\}}, so callers never find out the config is invalid. Let a parse error reach the caller, but keep returning \code{\{\}} for empty or blank text.\par
\begin{lstlisting}[language=JavaScript]
function readConfig(json) {
  try {
    return JSON.parse(json);
  } catch {
    return {};
  }
}
\end{lstlisting}
\cardsub{How to solve it}
\begin{cardrows}
\item[Solve] Remove the try/catch; return \code{\{\}} only for blank text; let \code{JSON.parse} throw otherwise.
\item[Invariants] Bad JSON always reaches the caller as a SyntaxError.
\item[Memory] No state.
\item[Types] string in; object out, or a SyntaxError.
\item[Failure] Catching every error turns broken input into empty settings; removing the blank check makes "" throw.
\item[Tests] Valid JSON, "", "   " and "\{not json".
\end{cardrows}
\trap Handle the expected case; don't catch everything.\par
\cardsub{Solution}
\begin{lstlisting}[language=JavaScript]
function readConfig(json) {
  if (json.trim() === "") return {};
  return JSON.parse(json);
}
\end{lstlisting}
{\small \textit{Verified: passes all 4 test cases on the card in the app's code runner; the starter code fails them.}}\par
\smallskip{\small\textbf{Why:} Swallowing an error turns "your file is broken" into "you have no settings". Handle the case you expect (empty text) and let real errors propagate.}\par
\cardsub{From the readings}
\readchap{10. What should a program do when something goes wrong?}{\code{throw} stops the function and travels up the calls until a \code{catch} handles it; catch only where you can do something useful, and \code{finally} runs whichever way the block exits.}
\readq{Where should you catch an error?}
Only where you can do something useful. \code{try \{ \ldots \} catch (error) \{ \ldots \}} runs the \code{catch} block if the \code{try} block throws. Catching and returning a default \emph{swallows} the error: the caller can no longer tell ``broken input'' from ``no data''. Handle the cases you expect (such as blank text) and let the rest propagate. \code{JSON.parse} throws a \code{SyntaxError} on text that isn't valid JSON.
\Needspace{6\baselineskip}
\begin{lstlisting}[language=JavaScript]
try {
  JSON.parse("{oops");
} catch (error) {
  console.log(error.name);
}
\end{lstlisting}
{\small\textit{Prints:}~\code{SyntaxError}}\par

\readchap{11. How do you turn data into text and back?}{\code{JSON.stringify} writes values as text and quietly drops what JSON cannot hold; \code{JSON.parse} reads the text back and throws on bad input.}
\readq{How do you read JSON?}
\textbf{JSON} is text. \code{JSON.parse(text)} turns it into JavaScript values. Missing fields come back as \code{undefined}, so supply defaults with \code{??} (\hyperref[ch:c02]{Ch.~2}), which keeps a saved \code{0} or \code{false}.
\Needspace{6\baselineskip}
\begin{lstlisting}[language=JavaScript]
const data = JSON.parse('{"volume": 0}');
console.log(data.volume ?? 50, data.muted ?? false);
\end{lstlisting}
{\small\textit{Prints:}~\code{0 false}}\par

\cardend
\cardhead{Card 49 --- ValidationError}{write · javascript · Medium}
\safemargin{\textbf{Terms for this card}\par\smallskip \textbf{Exception:} an error object thrown to signal that something went wrong. Like a fire alarm: work stops and everyone heads for the nearest exit.}
\noindent{\small\textbf{Covered in:} \hyperref[ch:c10]{Ch.~10}, \hyperref[ch:c12]{Ch.~12} of this guide; documentation: \href{https://developer.mozilla.org/en-US/docs/Web/JavaScript/Reference/Global_Objects/Error}{MDN — Error}.}\par
\cardsub{The card}
\textbf{ValidationError}\par
Make \code{ValidationError} an error whose \code{name} is \code{\textquotedbl{}ValidationError\textquotedbl{}}. Then make \code{validateUser(user)} throw one with the message \code{\textquotedbl{}name is required\textquotedbl{}} or \code{\textquotedbl{}email is required\textquotedbl{}} when that field is missing or empty, and return \code{true} otherwise.\par
\begin{lstlisting}[language=JavaScript]
function validateUser(user) {
  // TODO: throw a ValidationError for a missing name or email
  return true;
}

class ValidationError extends Error {
  // TODO: set this.name
}
\end{lstlisting}
\cardsub{How to solve it}
\begin{cardrows}
\item[Solve] Give \code{ValidationError} a constructor that calls \code{super(message)} and sets \code{this.name}; throw it for a missing name or email.
\item[Invariants] Every error thrown is a ValidationError with the right message.
\item[Memory] No state.
\item[Types] object in; true out, or a ValidationError.
\item[Failure] Without \code{this.name}, the error still reports itself as "Error".
\item[Tests] A valid user, a missing name, a missing email, an empty name.
\end{cardrows}
\trap A subclass of Error keeps the name "Error" unless you set it.\par
\cardsub{Solution}
\begin{lstlisting}[language=JavaScript]
function validateUser(user) {
  if (!user.name) throw new ValidationError("name is required");
  if (!user.email) throw new ValidationError("email is required");
  return true;
}

class ValidationError extends Error {
  constructor(message) {
    super(message);
    this.name = "ValidationError";
  }
}
\end{lstlisting}
{\small \textit{Verified: passes all 4 test cases on the card in the app's code runner; the starter code fails them.}}\par
\smallskip{\small\textbf{Why:} A subclass of Error inherits message and stack, but its name stays "Error" unless you set it. Callers can then check \code{error instanceof ValidationError}.}\par
\cardsub{From the readings}
\readchap{10. What should a program do when something goes wrong?}{\code{throw} stops the function and travels up the calls until a \code{catch} handles it; catch only where you can do something useful, and \code{finally} runs whichever way the block exits.}
\readq{How do you make your own kind of error?}
Extend \code{Error} (the class syntax is in \hyperref[ch:c12]{Ch.~12}) and set \code{name}. The subclass inherits \code{message}, but its name stays \code{"Error"} unless you set it. Callers can then tell your errors apart.
\Needspace{6\baselineskip}
\begin{lstlisting}[language=JavaScript]
class NotFound extends Error {
  name = "NotFound";
}
console.log(String(new NotFound("no user")));
\end{lstlisting}
{\small\textit{Prints:}~\code{NotFound: no user}}\par

\readq{What happens when you throw?}
The function stops and an \textbf{exception} travels up the call stack to the first \code{catch} that waits for it; with no \code{catch}, the program ends. Throw the most specific built-in error: \code{RangeError} for a value outside the allowed range. An error has a \code{name} and a \code{message}.
\Needspace{6\baselineskip}
\begin{lstlisting}[language=JavaScript]
function checkAge(age) {
  if (age < 0) throw new RangeError("negative age");
  return age;
}
try {
  checkAge(-1);
} catch (error) {
  console.log(`${error.name}: ${error.message}`);
}
\end{lstlisting}
{\small\textit{Prints:}~\code{RangeError: negative age}}\par

\cardend
\deckhead{11 · JSON}{4}
\begin{corebox}\textbf{Core idea.} \code{JSON.stringify} writes values as text and quietly drops what JSON cannot hold; \code{JSON.parse} reads the text back and throws on bad input.\end{corebox}
\begin{sumbox}{Deck summary sheet --- 11 · JSON}
\code{JSON.parse} turns text into values and throws a SyntaxError on bad input. \code{JSON.stringify} turns values into text: undefined and functions disappear from objects (null in arrays), Dates become strings, and Maps become \code{\{\}}.\par
\smallskip{\footnotesize Read more: \href{https://developer.mozilla.org/en-US/docs/Web/JavaScript/Reference/Global_Objects/JSON/stringify}{MDN — JSON.stringify()}}\end{sumbox}
\cardhead{Card 50 --- What does this print?}{predict · js11-stringify-drops}
\noindent{\small\textbf{Covered in:} \hyperref[ch:c11]{Ch.~11} of this guide; documentation: \href{https://developer.mozilla.org/en-US/docs/Web/JavaScript/Reference/Global_Objects/JSON/stringify}{MDN — JSON.stringify()}.}\par
\cardsub{The card}
What does this print?\par
\begin{lstlisting}[language=JavaScript]
const data = {
  name: "Ada",
  nickname: undefined,
  greet() {
    return "hi";
  },
  tags: [undefined, "math"],
  born: new Date(Date.UTC(1815, 11, 10)),
};
console.log(JSON.stringify(data));
\end{lstlisting}
\cardsub{How to solve it}
\begin{cardrows}
\item[Solve] Drop \code{nickname} and \code{greet} from the object, turn the \code{undefined} in the array into null, and write the Date as an ISO string.
\item[Invariants] Only values JSON can hold are written.
\item[Memory] The original object is unchanged; a new string is made.
\item[Types] object in; string out.
\item[Failure] Writing \code{\textquotedbl{}nickname\textquotedbl{}:null} or including the function.
\item[Tests] Run it, then change one line and predict again.
\end{cardrows}
\trap undefined is dropped in objects but becomes null in arrays.\par
\cardsub{Solution}
{\small\textbf{Prints:}}
\begin{lstlisting}
{"name":"Ada","tags":[null,"math"],"born":"1815-12-10T00:00:00.000Z"}
\end{lstlisting}
{\small\textit{Verified: Node v22.23.4 prints exactly this.}}\par
\smallskip{\small\textbf{Why:} undefined and functions are not JSON: they are left out of objects and become null in arrays. A Date is written as its ISO string (and comes back as a string, not a Date).}\par
\cardsub{From the readings}
\readchap{11. How do you turn data into text and back?}{\code{JSON.stringify} writes values as text and quietly drops what JSON cannot hold; \code{JSON.parse} reads the text back and throws on bad input.}
\readq{What does JSON.stringify leave out?}
Everything JSON cannot represent. \code{undefined} and functions are left out of objects and become \code{null} in arrays. A \code{Date} is written as a string, and a \code{Map} or \code{Set} becomes \code{"\{\}"} because only an object's own enumerable properties are written. The third argument sets the indentation: \code{JSON.stringify(value, null, 2)}.
\Needspace{6\baselineskip}
\begin{lstlisting}[language=JavaScript]
const value = { a: 1, b: undefined, c: [undefined], m: new Map() };
console.log(JSON.stringify(value));
console.log(JSON.stringify([1], null, 2).split("\n").length);
\end{lstlisting}
{\small\textit{Prints:}~\code{\{\textquotedbl{}a\textquotedbl{}:1,\textquotedbl{}c\textquotedbl{}:[null],\textquotedbl{}m\textquotedbl{}:\{\}\}}\enspace\textcolor{black!45}{|}\enspace \code{3}}\par

\cardend
\cardhead{Card 51 --- loadSettings(json)}{write · javascript · Easy}
\safemargin{\textbf{Terms for this card}\par\smallskip \textbf{JSON:} a text format for data (objects, arrays, strings, numbers, booleans, null) that any program can read. Like a shipping label written in a standard layout.\par\smallskip \textbf{Nullish coalescing (??):} returns its right side only when its left side is \code{null} or \code{undefined}. Like a substitute who steps in only if the player is absent, not if the player scores zero.}
\noindent{\small\textbf{Covered in:} \hyperref[ch:c11]{Ch.~11} of this guide; documentation: \href{https://developer.mozilla.org/en-US/docs/Web/JavaScript/Reference/Global_Objects/JSON/parse}{MDN — JSON.parse()}.}\par
\cardsub{The card}
\textbf{loadSettings(json)}\par
Write \code{loadSettings(json)}: parse the JSON text and return \code{\{ volume, muted \}}, with \code{volume} 50 and \code{muted} false when missing. A saved 0 or false must be kept. Ignore any other fields.\par
\begin{lstlisting}[language=JavaScript]
function loadSettings(json) {
  const data = JSON.parse(json);
  // TODO
  return data;
}
\end{lstlisting}
\cardsub{How to solve it}
\begin{cardrows}
\item[Solve] Parse, then return \code{\{ volume: data.volume ?? 50, muted: data.muted ?? false \}}.
\item[Invariants] The result has exactly two fields; saved 0 and false survive.
\item[Memory] One parsed object.
\item[Types] string in; object out.
\item[Failure] \code{||} replaces a saved 0 or false; spreading the parsed data keeps unwanted fields.
\item[Tests] All fields, none, volume 0, and an extra field.
\end{cardrows}
\trap \code{||} throws away a saved 0.\par
\cardsub{Solution}
\begin{lstlisting}[language=JavaScript]
function loadSettings(json) {
  const data = JSON.parse(json);
  return { volume: data.volume ?? 50, muted: data.muted ?? false };
}
\end{lstlisting}
{\small \textit{Verified: passes all 4 test cases on the card in the app's code runner; the starter code fails them.}}\par
\smallskip{\small\textbf{Why:} Build the result from just the fields you want, with \code{??} so real 0 and false values survive.}\par
\cardsub{From the readings}
\readchap{11. How do you turn data into text and back?}{\code{JSON.stringify} writes values as text and quietly drops what JSON cannot hold; \code{JSON.parse} reads the text back and throws on bad input.}
\readq{How do you read JSON?}
\textbf{JSON} is text. \code{JSON.parse(text)} turns it into JavaScript values. Missing fields come back as \code{undefined}, so supply defaults with \code{??} (\hyperref[ch:c02]{Ch.~2}), which keeps a saved \code{0} or \code{false}.
\Needspace{6\baselineskip}
\begin{lstlisting}[language=JavaScript]
const data = JSON.parse('{"volume": 0}');
console.log(data.volume ?? 50, data.muted ?? false);
\end{lstlisting}
{\small\textit{Prints:}~\code{0 false}}\par

\readchap{2. How does JavaScript compare values and choose a branch?}{\code{===} compares without converting; truthiness decides an \code{if}; and \code{||}, \code{\&\&} and \code{??} hand back one of their operands.}
\readq{What do ||, \&\& and ?? really return?}
One of their operands, not \code{true} or \code{false}. \code{a || b} returns \code{a} if it is truthy, else \code{b}; \code{a \&\& b} returns \code{a} if it is falsy, else \code{b}. The \textbf{nullish coalescing} operator \code{??} returns its right side only for \code{null} or \code{undefined}, so a real \code{0} or \code{""} survives. That makes it the right tool for defaults.
\Needspace{6\baselineskip}
\begin{lstlisting}[language=JavaScript]
const saved = 0;
console.log(saved || 10, saved ?? 10, null ?? 10);
\end{lstlisting}
{\small\textit{Prints:}~\code{10 0 10}}\par

\cardend
\cardhead{Card 52 --- prettyJson(text)}{write · javascript · Easy}
\safemargin{\textbf{Terms for this card}\par\smallskip \textbf{JSON:} a text format for data (objects, arrays, strings, numbers, booleans, null) that any program can read. Like a shipping label written in a standard layout.}
\noindent{\small\textbf{Covered in:} \hyperref[ch:c11]{Ch.~11} of this guide; documentation: \href{https://developer.mozilla.org/en-US/docs/Web/JavaScript/Reference/Global_Objects/JSON/stringify}{MDN — JSON.stringify()}.}\par
\cardsub{The card}
\textbf{prettyJson(text)}\par
Write \code{prettyJson(text)}: reformat compact JSON text with 2-space indentation.\par
\begin{lstlisting}[language=JavaScript]
function prettyJson(text) {
  // TODO: parse, then stringify with indentation
  return text;
}
\end{lstlisting}
\cardsub{How to solve it}
\begin{cardrows}
\item[Solve] \code{JSON.stringify(JSON.parse(text), null, 2)}.
\item[Invariants] The output holds the same data as the input.
\item[Memory] One parsed value.
\item[Types] string in; string out.
\item[Failure] Stringifying the text itself just adds quotes around it.
\item[Tests] An object with an array, an empty array, a plain string.
\end{cardrows}
\trap Parse first, then stringify with indentation.\par
\cardsub{Solution}
\begin{lstlisting}[language=JavaScript]
function prettyJson(text) {
  return JSON.stringify(JSON.parse(text), null, 2);
}
\end{lstlisting}
{\small \textit{Verified: passes all 3 test cases on the card in the app's code runner; the starter code fails them.}}\par
\smallskip{\small\textbf{Why:} The third argument of JSON.stringify is the indentation. Stringifying the text itself (without parsing) just wraps it in quotes.}\par
\cardsub{From the readings}
\readchap{11. How do you turn data into text and back?}{\code{JSON.stringify} writes values as text and quietly drops what JSON cannot hold; \code{JSON.parse} reads the text back and throws on bad input.}
\readq{What does JSON.stringify leave out?}
Everything JSON cannot represent. \code{undefined} and functions are left out of objects and become \code{null} in arrays. A \code{Date} is written as a string, and a \code{Map} or \code{Set} becomes \code{"\{\}"} because only an object's own enumerable properties are written. The third argument sets the indentation: \code{JSON.stringify(value, null, 2)}.
\Needspace{6\baselineskip}
\begin{lstlisting}[language=JavaScript]
const value = { a: 1, b: undefined, c: [undefined], m: new Map() };
console.log(JSON.stringify(value));
console.log(JSON.stringify([1], null, 2).split("\n").length);
\end{lstlisting}
{\small\textit{Prints:}~\code{\{\textquotedbl{}a\textquotedbl{}:1,\textquotedbl{}c\textquotedbl{}:[null],\textquotedbl{}m\textquotedbl{}:\{\}\}}\enspace\textcolor{black!45}{|}\enspace \code{3}}\par

\readq{How do you read JSON?}
\textbf{JSON} is text. \code{JSON.parse(text)} turns it into JavaScript values. Missing fields come back as \code{undefined}, so supply defaults with \code{??} (\hyperref[ch:c02]{Ch.~2}), which keeps a saved \code{0} or \code{false}.
\Needspace{6\baselineskip}
\begin{lstlisting}[language=JavaScript]
const data = JSON.parse('{"volume": 0}');
console.log(data.volume ?? 50, data.muted ?? false);
\end{lstlisting}
{\small\textit{Prints:}~\code{0 false}}\par

\cardend
\cardhead{Card 53 --- What does this print?}{predict · js11-map-to-json}
\safemargin{\textbf{Terms for this card}\par\smallskip \textbf{Map:} a collection of key--value pairs where the key can be any value and pairs stay in insertion order. Like a coat check: any ticket can claim exactly one coat.}
\noindent{\small\textbf{Covered in:} \hyperref[ch:c11]{Ch.~11}, \hyperref[ch:c08]{Ch.~8} of this guide; documentation: \href{https://developer.mozilla.org/en-US/docs/Web/JavaScript/Reference/Global_Objects/JSON/stringify}{MDN — JSON.stringify()}.}\par
\cardsub{The card}
What does this print?\par
\begin{lstlisting}[language=JavaScript]
const scores = new Map([["ada", 3]]);
console.log(JSON.stringify(scores));
console.log(JSON.stringify(Object.fromEntries(scores)));
\end{lstlisting}
\cardsub{How to solve it}
\begin{cardrows}
\item[Solve] A Map has no own enumerable properties, so it becomes \{\}; converting it first keeps its entries.
\item[Invariants] stringify writes only an object's own enumerable properties.
\item[Memory] One Map, one plain object.
\item[Types] Map, object, strings.
\item[Failure] Expecting \code{\{\textquotedbl{}ada\textquotedbl{}:3\}} on the first line.
\item[Tests] Run it, then change one line and predict again.
\end{cardrows}
\trap A Map stringifies to \{\}.\par
\cardsub{Solution}
{\small\textbf{Prints:}}
\begin{lstlisting}
{}
{"ada":3}
\end{lstlisting}
{\small\textit{Verified: Node v22.23.4 prints exactly this.}}\par
\smallskip{\small\textbf{Why:} JSON.stringify only writes an object's own enumerable properties, and a Map has none. Convert it first.}\par
\cardsub{From the readings}
\readchap{11. How do you turn data into text and back?}{\code{JSON.stringify} writes values as text and quietly drops what JSON cannot hold; \code{JSON.parse} reads the text back and throws on bad input.}
\readq{What does JSON.stringify leave out?}
Everything JSON cannot represent. \code{undefined} and functions are left out of objects and become \code{null} in arrays. A \code{Date} is written as a string, and a \code{Map} or \code{Set} becomes \code{"\{\}"} because only an object's own enumerable properties are written. The third argument sets the indentation: \code{JSON.stringify(value, null, 2)}.
\Needspace{6\baselineskip}
\begin{lstlisting}[language=JavaScript]
const value = { a: 1, b: undefined, c: [undefined], m: new Map() };
console.log(JSON.stringify(value));
console.log(JSON.stringify([1], null, 2).split("\n").length);
\end{lstlisting}
{\small\textit{Prints:}~\code{\{\textquotedbl{}a\textquotedbl{}:1,\textquotedbl{}c\textquotedbl{}:[null],\textquotedbl{}m\textquotedbl{}:\{\}\}}\enspace\textcolor{black!45}{|}\enspace \code{3}}\par

\readchap{8. Which collection fits: object, Map or Set?}{A \code{Set} keeps one of each value; a \code{Map} pairs any key with a value in insertion order; a plain object turns keys into strings and inherits keys you never added.}
\readq{When is a Map better than an object?}
When you count or group by key. A \textbf{Map} stores key--value pairs: \code{set(key, value)}, \code{get(key)}, \code{has(key)}, \code{size}. Keys keep their type, so the number \code{1} and the string \code{"1"} are different keys; an object would turn both into \code{"1"}. A Map also iterates in insertion order, which settles ties ``first come, first served''. \code{counts.get(word) ?? 0} starts a new key at zero.
\Needspace{6\baselineskip}
\begin{lstlisting}[language=JavaScript]
const counts = new Map();
for (const w of ["a", "b", "a"]) {
  counts.set(w, (counts.get(w) ?? 0) + 1);
}
console.log(counts.get("a"), counts.size);
\end{lstlisting}
{\small\textit{Prints:}~\code{2 2}}\par

\cardend
\deckhead{12 · Classes \& this}{4}
\begin{corebox}\textbf{Core idea.} A class is a template for objects that share methods; \code{\#fields} are private; \code{this} is set by how a method is called, so a detached method or a \code{function} callback loses it.\end{corebox}
\begin{sumbox}{Deck summary sheet --- 12 · Classes \& this}
A class bundles state and methods; \code{\#field} makes state truly private. \code{extends} inherits and \code{super(...)} runs the parent constructor. Inside a method \code{this} is whatever the method was called on, so pulling a method off its object, or using a \code{function} callback, loses it. Arrow functions and \code{bind} keep it.\par
\smallskip{\footnotesize Read more: \href{https://developer.mozilla.org/en-US/docs/Web/JavaScript/Guide/Using_classes}{MDN — Using classes}; \href{https://developer.mozilla.org/en-US/docs/Web/JavaScript/Reference/Operators/this}{MDN — this}}\end{sumbox}
\cardhead{Card 54 --- class BankAccount}{write · javascript · Medium}
\safemargin{\textbf{Terms for this card}\par\smallskip \textbf{Class:} a template for making objects that share the same methods. Like a blueprint.\par\smallskip \textbf{Instance:} one object made from a class. A house built from the blueprint.\par\smallskip \textbf{Private field:} a class field whose name starts with \code{\#}; only code inside the class can read or write it. Like a safe that opens only from the inside.}
\noindent{\small\textbf{Covered in:} \hyperref[ch:c12]{Ch.~12} of this guide; documentation: \href{https://developer.mozilla.org/en-US/docs/Web/JavaScript/Guide/Using_classes}{MDN — Using classes}.}\par
\cardsub{The card}
\textbf{class BankAccount}\par
Finish \code{BankAccount}: keep the balance in a private \code{\#balance} field starting at 0. \code{deposit(amount)} and \code{withdraw(amount)} return the new balance; \code{getBalance()} returns it. Throw \code{RangeError(\textquotedbl{}deposit must be positive\textquotedbl{})} for a deposit of 0 or less and \code{RangeError(\textquotedbl{}insufficient funds\textquotedbl{})} when withdrawing more than the balance. Don't change \code{run}.\par
\begin{lstlisting}[language=JavaScript]
// Test driver: don't change. Runs each [method, ...args] and records the result.
function run(actions) {
  const account = new BankAccount("Ada");
  return actions.map(([method, ...args]) => {
    if (method === "publicBalance") return account.balance;
    try {
      return account[method](...args);
    } catch (error) {
      return `${error.name}: ${error.message}`;
    }
  });
}

class BankAccount {
  constructor(owner) {
    this.owner = owner;
  }

  // TODO: #balance, deposit(amount), withdraw(amount), getBalance()
}
\end{lstlisting}
\cardsub{How to solve it}
\begin{cardrows}
\item[Solve] Declare \code{\#balance = 0}; validate before changing it; return the new balance from each method.
\item[Invariants] The balance never goes below 0 and is never reachable from outside.
\item[Memory] One private number per account.
\item[Types] numbers in; numbers out, or RangeErrors.
\item[Failure] A public \code{balance} property is visible from outside; \code{\textgreater{}=} blocks withdrawing the exact balance.
\item[Tests] Deposit and withdraw, overdraw, withdraw everything, non-positive deposits, and the public property.
\end{cardrows}
\trap Validate first, so a failed call leaves the balance unchanged.\par
\cardsub{Solution}
\begin{lstlisting}[language=JavaScript]
// Test driver: don't change. Runs each [method, ...args] and records the result.
function run(actions) {
  const account = new BankAccount("Ada");
  return actions.map(([method, ...args]) => {
    if (method === "publicBalance") return account.balance;
    try {
      return account[method](...args);
    } catch (error) {
      return `${error.name}: ${error.message}`;
    }
  });
}

class BankAccount {
  #balance = 0;

  constructor(owner) {
    this.owner = owner;
  }

  deposit(amount) {
    if (amount <= 0) throw new RangeError("deposit must be positive");
    this.#balance += amount;
    return this.#balance;
  }

  withdraw(amount) {
    if (amount > this.#balance) throw new RangeError("insufficient funds");
    this.#balance -= amount;
    return this.#balance;
  }

  getBalance() {
    return this.#balance;
  }
}
\end{lstlisting}
{\small \textit{Verified: passes all 5 test cases on the card in the app's code runner; the starter code fails them.}}\par
\smallskip{\small\textbf{Why:} A \# field can only be read inside the class: \code{account.balance} from outside is undefined. Validate before changing state so a failed call leaves it untouched.}\par
\cardsub{From the readings}
\readchap{12. How do classes bundle data with behaviour, and what is this?}{A class is a template for objects that share methods; \code{\#fields} are private; \code{this} is set by how a method is called, so a detached method or a \code{function} callback loses it.}
\readq{What does a class give you?}
A \textbf{class} describes objects; \code{new} makes an \textbf{instance}. In the \code{constructor}, \code{this} is the new instance, so the constructor stores its starting data on \code{this}. Methods are written once in the class and shared by every instance.
\Needspace{6\baselineskip}
\begin{lstlisting}[language=JavaScript]
class Point {
  constructor(x, y) { this.x = x; this.y = y; }
  sum() { return this.x + this.y; }
}
console.log(new Point(2, 3).sum());
\end{lstlisting}
{\small\textit{Prints:}~\code{5}}\par

\readq{How do you keep state private?}
Prefix the field name with \code{\#}. A \textbf{private field} must be declared in the class body and can only be used inside the class. From outside, \code{obj.balance} is just a different, missing property (\code{undefined}).
\Needspace{6\baselineskip}
\begin{lstlisting}[language=JavaScript]
class Safe {
  #code = 42;
  check(n) { return n === this.#code; }
}
const safe = new Safe();
console.log(safe.check(42), safe.code);
\end{lstlisting}
{\small\textit{Prints:}~\code{true undefined}}\par

\cardend
\cardhead{Card 55 --- What does this print?}{predict · js12-lost-this}
\noindent{\small\textbf{Covered in:} \hyperref[ch:c12]{Ch.~12} of this guide; documentation: \href{https://developer.mozilla.org/en-US/docs/Web/JavaScript/Reference/Global_Objects/Function/bind}{MDN — Function.prototype.bind()}.}\par
\cardsub{The card}
What does this print?\par
\begin{lstlisting}[language=JavaScript]
class Greeter {
  constructor(name) {
    this.name = name;
  }
  greet() {
    return `Hi, ${this.name}`;
  }
}
const greeter = new Greeter("Ada");
const { greet } = greeter;
try {
  console.log(greet());
} catch (error) {
  console.log(error.name);
}
console.log(greet.call(greeter));
const bound = greeter.greet.bind(greeter);
console.log(bound());
\end{lstlisting}
\cardsub{How to solve it}
\begin{cardrows}
\item[Solve] Called on its own, \code{greet} has \code{this} undefined and throws; \code{call} and \code{bind} supply \code{this}.
\item[Invariants] \code{this} is set by the call, not by where the method was defined.
\item[Memory] One instance; a detached function; a bound copy.
\item[Types] functions, strings, TypeError.
\item[Failure] Expecting "Hi, undefined" instead of a TypeError on the first line.
\item[Tests] Run it, then change one line and predict again.
\end{cardrows}
\trap A detached method loses \code{this}; in class code that throws.\par
\cardsub{Solution}
{\small\textbf{Prints:}}
\begin{lstlisting}
TypeError
Hi, Ada
Hi, Ada
\end{lstlisting}
{\small\textit{Verified: Node v22.23.4 prints exactly this.}}\par
\smallskip{\small\textbf{Why:} this is set by how a function is called, not where it was defined. Called on its own, greet gets this = undefined (class code is strict), so this.name throws. call() passes this for one call; bind() makes a copy with this fixed.}\par
\cardsub{From the readings}
\readchap{12. How do classes bundle data with behaviour, and what is this?}{A class is a template for objects that share methods; \code{\#fields} are private; \code{this} is set by how a method is called, so a detached method or a \code{function} callback loses it.}
\readq{Why does this go missing?}
Because \code{this} is decided at the moment of the call: \code{greeter.greet()} sets it to \code{greeter}, but calling a method pulled off its object sets it to \code{undefined} (class and module code is strict), and reading a property of \code{undefined} throws. \code{fn.call(obj)} passes \code{this} for one call; \code{fn.bind(obj)} makes a copy with \code{this} fixed. An arrow function has no \code{this} of its own and uses the one around it, which is why arrow callbacks inside methods work and \code{function} callbacks do not.
\Needspace{6\baselineskip}
\begin{lstlisting}[language=JavaScript]
const counter = { n: 1, get() { return this.n; } };
const loose = counter.get;
try {
  loose();
} catch (error) {
  console.log(error.name);
}
console.log(loose.call(counter), loose.bind(counter)());
\end{lstlisting}
{\small\textit{Prints:}~\code{TypeError}\enspace\textcolor{black!45}{|}\enspace \code{1 1}}\par

\cardend
\cardhead{Card 56 --- Fix Cart.total()}{fix · javascript · Medium}
\noindent{\small\textbf{Covered in:} \hyperref[ch:c12]{Ch.~12} of this guide; documentation: \href{https://developer.mozilla.org/en-US/docs/Web/JavaScript/Reference/Functions/Arrow_functions}{MDN — Arrow function expressions}.}\par
\cardsub{The card}
\textbf{Fix Cart.total()}\par
\code{total()} throws a TypeError: inside the \code{function} callback, \code{this} is not the cart. Fix it. Don't change \code{check}.\par
\begin{lstlisting}[language=JavaScript]
// Test driver: don't change.
function check(prices) {
  return new Cart(prices).total();
}

class Cart {
  constructor(prices) {
    this.prices = prices;
    this.sum = 0;
  }

  total() {
    this.sum = 0;
    this.prices.forEach(function (price) {
      this.sum += price;
    });
    return this.sum;
  }
}
\end{lstlisting}
\cardsub{How to solve it}
\begin{cardrows}
\item[Solve] Change the \code{function (price)} callback to an arrow function.
\item[Invariants] Inside the arrow function, \code{this} is the cart.
\item[Memory] The cart's \code{sum} field, updated once per price.
\item[Types] array of numbers in; number out.
\item[Failure] A \code{function} callback gets its own \code{this} (undefined), so \code{this.sum} throws.
\item[Tests] Several prices, one price, none.
\end{cardrows}
\trap Arrow functions use the surrounding \code{this}; \code{function} callbacks don't.\par
\cardsub{Solution}
\begin{lstlisting}[language=JavaScript]
// Test driver: don't change.
function check(prices) {
  return new Cart(prices).total();
}

class Cart {
  constructor(prices) {
    this.prices = prices;
    this.sum = 0;
  }

  total() {
    this.sum = 0;
    this.prices.forEach((price) => {
      this.sum += price;
    });
    return this.sum;
  }
}
\end{lstlisting}
{\small \textit{Verified: passes all 3 test cases on the card in the app's code runner; the starter code fails them.}}\par
\smallskip{\small\textbf{Why:} A \code{function} callback gets its own \code{this} (undefined here). An arrow function has no \code{this} of its own, so it uses the method's.}\par
\cardsub{From the readings}
\readchap{12. How do classes bundle data with behaviour, and what is this?}{A class is a template for objects that share methods; \code{\#fields} are private; \code{this} is set by how a method is called, so a detached method or a \code{function} callback loses it.}
\readq{Why does this go missing?}
Because \code{this} is decided at the moment of the call: \code{greeter.greet()} sets it to \code{greeter}, but calling a method pulled off its object sets it to \code{undefined} (class and module code is strict), and reading a property of \code{undefined} throws. \code{fn.call(obj)} passes \code{this} for one call; \code{fn.bind(obj)} makes a copy with \code{this} fixed. An arrow function has no \code{this} of its own and uses the one around it, which is why arrow callbacks inside methods work and \code{function} callbacks do not.
\Needspace{6\baselineskip}
\begin{lstlisting}[language=JavaScript]
const counter = { n: 1, get() { return this.n; } };
const loose = counter.get;
try {
  loose();
} catch (error) {
  console.log(error.name);
}
console.log(loose.call(counter), loose.bind(counter)());
\end{lstlisting}
{\small\textit{Prints:}~\code{TypeError}\enspace\textcolor{black!45}{|}\enspace \code{1 1}}\par

\cardend
\cardhead{Card 57 --- Circle and Rectangle extend Shape}{write · javascript · Medium}
\safemargin{\textbf{Terms for this card}\par\smallskip \textbf{Class:} a template for making objects that share the same methods. Like a blueprint.\par\smallskip \textbf{Instance:} one object made from a class. A house built from the blueprint.}
\noindent{\small\textbf{Covered in:} \hyperref[ch:c12]{Ch.~12} of this guide; documentation: \href{https://developer.mozilla.org/en-US/docs/Web/JavaScript/Reference/Classes/extends}{MDN — extends}.}\par
\cardsub{The card}
\textbf{Circle and Rectangle extend Shape}\par
Write \code{Circle} (constructor \code{radius}, name \code{\textquotedbl{}circle\textquotedbl{}}, area $\pi$r$^2$) and \code{Rectangle} (constructor \code{width, height}, name \code{\textquotedbl{}rectangle\textquotedbl{}}) so both extend \code{Shape}, call \code{super(name)}, and implement \code{area()}. Don't change \code{check} or \code{Shape}.\par
\begin{lstlisting}[language=JavaScript]
// Test driver: don't change.
function check(specs) {
  return specs.map(([kind, ...dims]) => {
    const shape = kind === "circle" ? new Circle(...dims) : new Rectangle(...dims);
    return [shape instanceof Shape, shape.describe()];
  });
}

class Shape {
  constructor(name) {
    this.name = name;
  }
  area() {
    return 0;
  }
  describe() {
    return `${this.name} with area ${this.area().toFixed(2)}`;
  }
}

class Circle extends Shape {
  // TODO
}

class Rectangle extends Shape {
  // TODO
}
\end{lstlisting}
\cardsub{How to solve it}
\begin{cardrows}
\item[Solve] Write both classes with \code{extends Shape}, call \code{super(name)} first, store the dimensions, and implement \code{area()}.
\item[Invariants] Every shape is a Shape and describes itself with its own area.
\item[Memory] One instance per shape with its dimensions.
\item[Types] numbers in; Shape instances out.
\item[Failure] Standalone classes don't extend Shape (instanceof is false); using \code{this} before \code{super} throws.
\item[Tests] A circle of radius 1 and 2, and a 2 $\times$ 3 rectangle.
\end{cardrows}
\trap Call super() before touching this.\par
\cardsub{Solution}
\begin{lstlisting}[language=JavaScript]
// Test driver: don't change.
function check(specs) {
  return specs.map(([kind, ...dims]) => {
    const shape = kind === "circle" ? new Circle(...dims) : new Rectangle(...dims);
    return [shape instanceof Shape, shape.describe()];
  });
}

class Shape {
  constructor(name) {
    this.name = name;
  }
  area() {
    return 0;
  }
  describe() {
    return `${this.name} with area ${this.area().toFixed(2)}`;
  }
}

class Circle extends Shape {
  constructor(radius) {
    super("circle");
    this.radius = radius;
  }
  area() {
    return Math.PI * this.radius ** 2;
  }
}

class Rectangle extends Shape {
  constructor(width, height) {
    super("rectangle");
    this.width = width;
    this.height = height;
  }
  area() {
    return this.width * this.height;
  }
}
\end{lstlisting}
{\small \textit{Verified: passes all 2 test cases on the card in the app's code runner; the starter code fails them.}}\par
\smallskip{\small\textbf{Why:} A subclass must call super() before it uses \code{this}. It inherits describe() and overrides area(), so Shape's describe() calls the right area for each shape.}\par
\cardsub{From the readings}
\readchap{12. How do classes bundle data with behaviour, and what is this?}{A class is a template for objects that share methods; \code{\#fields} are private; \code{this} is set by how a method is called, so a detached method or a \code{function} callback loses it.}
\readq{How does one class build on another?}
With \code{extends}: the child class gets the parent's methods and can replace (override) some. A child constructor must call \code{super(...)}---the parent's constructor---before it uses \code{this}. A method inherited from the parent that calls \code{this.area()} gets the child's version.
\Needspace{6\baselineskip}
\begin{lstlisting}[language=JavaScript]
class Animal {
  constructor(name) { this.name = name; }
  describe() { return `${this.name}: ${this.sound()}`; }
  sound() { return "..."; }
}
class Dog extends Animal {
  sound() { return "woof"; }
}
console.log(new Dog("Rex").describe());
\end{lstlisting}
{\small\textit{Prints:}~\code{Rex: woof}}\par

\readq{What does a class give you?}
A \textbf{class} describes objects; \code{new} makes an \textbf{instance}. In the \code{constructor}, \code{this} is the new instance, so the constructor stores its starting data on \code{this}. Methods are written once in the class and shared by every instance.
\Needspace{6\baselineskip}
\begin{lstlisting}[language=JavaScript]
class Point {
  constructor(x, y) { this.x = x; this.y = y; }
  sum() { return this.x + this.y; }
}
console.log(new Point(2, 3).sum());
\end{lstlisting}
{\small\textit{Prints:}~\code{5}}\par

\cardend
\deckhead{13 · Asynchronous JavaScript}{4}
\begin{corebox}\textbf{Core idea.} A promise stands for a result that isn't ready yet; \code{await} pauses only its own async function; promise callbacks run before timers; \code{Promise.all} waits for every promise and fails on the first rejection.\end{corebox}
\begin{sumbox}{Deck summary sheet --- 13 · Asynchronous JavaScript}
Promise callbacks (microtasks) run as soon as the current code finishes, before timers. \code{await} pauses only the async function it is in. Work starts when you create the promises; \code{Promise.all} just waits for them all and rejects on the first failure. \code{fetch} rejects only on network failure — check \code{response.ok} yourself.\par
\smallskip{\footnotesize Read more: \href{https://developer.mozilla.org/en-US/docs/Web/JavaScript/Guide/Using_promises}{MDN — Using promises}; \href{https://developer.mozilla.org/en-US/docs/Web/JavaScript/Reference/Global_Objects/Promise/all}{MDN — Promise.all()}; \href{https://developer.mozilla.org/en-US/docs/Web/API/Window/fetch}{MDN — Window: fetch()}}\end{sumbox}
\cardhead{Card 58 --- What does this print?}{predict · js13-event-loop}
\safemargin{\textbf{Terms for this card}\par\smallskip \textbf{Microtask:} a small job (such as a promise callback) run as soon as the current code finishes, before timers. Like the express lane at a checkout.}
\noindent{\small\textbf{Covered in:} \hyperref[ch:c13]{Ch.~13} of this guide; documentation: \href{https://developer.mozilla.org/en-US/docs/Web/JavaScript/Guide/Using_promises}{MDN — Using promises}.}\par
\cardsub{The card}
What does this print?\par
\begin{lstlisting}[language=JavaScript]
console.log("A");
setTimeout(() => console.log("B"), 0);
Promise.resolve().then(() => console.log("C"));
console.log("D");
\end{lstlisting}
\cardsub{How to solve it}
\begin{cardrows}
\item[Solve] Run the synchronous logs (A, D), then the promise callback (C), then the timer (B).
\item[Invariants] Microtasks run before the next timer.
\item[Memory] One timer and one promise callback waiting.
\item[Types] strings.
\item[Failure] Putting B before C because its delay is 0.
\item[Tests] Run it, then change one line and predict again.
\end{cardrows}
\trap A 0 ms timer still runs after promise callbacks.\par
\cardsub{Solution}
{\small\textbf{Prints:}}
\begin{lstlisting}
A
D
C
B
\end{lstlisting}
{\small\textit{Verified: Node v22.23.4 prints exactly this.}}\par
\smallskip{\small\textbf{Why:} Synchronous code runs first (A, D). Then the microtask queue — promise callbacks — is emptied (C). Timers run after that, even with a 0 ms delay (B).}\par
\cardsub{From the readings}
\readchap{13. How does JavaScript wait without stopping?}{A promise stands for a result that isn't ready yet; \code{await} pauses only its own async function; promise callbacks run before timers; \code{Promise.all} waits for every promise and fails on the first rejection.}
\readq{In what order do callbacks run?}
Synchronous code first. Promise callbacks are queued as \textbf{microtasks} and run as soon as the current code finishes; timer callbacks (\code{setTimeout}) wait in a separate queue and run after that, even with a 0 ms delay.
\Needspace{6\baselineskip}
\begin{lstlisting}[language=JavaScript]
setTimeout(() => console.log("timer"), 0);
Promise.resolve().then(() => console.log("promise"));
console.log("sync");
\end{lstlisting}
{\small\textit{Prints:}~\code{sync}\enspace\textcolor{black!45}{|}\enspace \code{promise}\enspace\textcolor{black!45}{|}\enspace \code{timer}}\par

\cardend
\cardhead{Card 59 --- Make loadNames(ids) concurrent}{fix · javascript · Medium}
\safemargin{\textbf{Terms for this card}\par\smallskip \textbf{Promise:} an object standing for a result that will arrive later, or a failure. Like a claim ticket for a coat that isn't ready yet.\par\smallskip \textbf{Async function:} a function declared with \code{async}; it always returns a promise and may use \code{await}. Like a cook who can step away from one dish while it simmers.}
\noindent{\small\textbf{Covered in:} \hyperref[ch:c13]{Ch.~13} of this guide; documentation: \href{https://developer.mozilla.org/en-US/docs/Web/JavaScript/Reference/Global_Objects/Promise/all}{MDN — Promise.all()}.}\par
\cardsub{The card}
\textbf{Make loadNames(ids) concurrent}\par
\code{loadNames(ids)} returns the right names but fetches one user at a time. Make all the \code{getUser} calls run at once with \code{Promise.all}, keeping the names in the order of \code{ids}. Don't change \code{check} or \code{getUser}.\par
\begin{lstlisting}[language=JavaScript]
// Test driver: don't change.
async function check(ids) {
  maxInFlight = 0;
  const names = await loadNames(ids);
  return { names, concurrent: ids.length < 2 || maxInFlight === ids.length };
}

async function loadNames(ids) {
  const names = [];
  for (const id of ids) {
    const user = await getUser(id);
    names.push(user.name);
  }
  return names;
}

// Fake API: don't change.
let inFlight = 0;
let maxInFlight = 0;
function getUser(id) {
  inFlight++;
  maxInFlight = Math.max(maxInFlight, inFlight);
  return new Promise((resolve) => {
    setTimeout(() => {
      inFlight--;
      resolve({ id, name: `user${id}` });
    }, 10);
  });
}
\end{lstlisting}
\cardsub{How to solve it}
\begin{cardrows}
\item[Solve] Start every request with \code{ids.map(getUser)} and \code{await Promise.all(...)}, then map users to names.
\item[Invariants] Names come back in the order of \code{ids}.
\item[Memory] One array of pending promises, then one of users.
\item[Types] array of numbers in; promise of array of strings out.
\item[Failure] \code{await} inside a loop runs the requests one at a time.
\item[Tests] Three ids (must overlap), none, one.
\end{cardrows}
\trap await in a loop is sequential.\par
\cardsub{Solution}
\begin{lstlisting}[language=JavaScript]
// Test driver: don't change.
async function check(ids) {
  maxInFlight = 0;
  const names = await loadNames(ids);
  return { names, concurrent: ids.length < 2 || maxInFlight === ids.length };
}

async function loadNames(ids) {
  const users = await Promise.all(ids.map((id) => getUser(id)));
  return users.map((user) => user.name);
}

// Fake API: don't change.
let inFlight = 0;
let maxInFlight = 0;
function getUser(id) {
  inFlight++;
  maxInFlight = Math.max(maxInFlight, inFlight);
  return new Promise((resolve) => {
    setTimeout(() => {
      inFlight--;
      resolve({ id, name: `user${id}` });
    }, 10);
  });
}
\end{lstlisting}
{\small \textit{Verified: passes all 3 test cases on the card in the app's code runner; the starter code fails them.}}\par
\smallskip{\small\textbf{Why:} \code{await} in a loop waits for each request before starting the next. Start them all with map, then await Promise.all, which keeps results in input order.}\par
\cardsub{From the readings}
\readchap{13. How does JavaScript wait without stopping?}{A promise stands for a result that isn't ready yet; \code{await} pauses only its own async function; promise callbacks run before timers; \code{Promise.all} waits for every promise and fails on the first rejection.}
\readq{How do you run several requests at once?}
Start them all, then wait for all. A promise's work starts when the promise is created; \code{Promise.all} only waits. It fulfills with an array of the values in input order, or rejects with the first rejection. \code{await} inside a loop waits for each request before starting the next, so the requests run one at a time. \code{Promise.allSettled} always fulfills, with a status for each promise.
\Needspace{6\baselineskip}
\begin{lstlisting}[language=JavaScript]
const wait = (ms, v) => new Promise((r) => setTimeout(() => r(v), ms));
const start = Date.now();
const both = await Promise.all([wait(100, "a"), wait(100, "b")]);
console.log(both.join(""), Date.now() - start < 190);
\end{lstlisting}
{\small\textit{Prints:}~\code{ab true}}\par

\readq{What is a promise, and what does await do?}
A \textbf{promise} is either still pending, fulfilled with a value, or rejected with an error. Inside an \textbf{async function}, \code{await p} pauses that function until \code{p} settles, then gives its value (or throws its error). The rest of the program keeps running meanwhile.
\Needspace{6\baselineskip}
\begin{lstlisting}[language=JavaScript]
async function slowDouble(n) {
  await null;
  return n * 2;
}
slowDouble(4).then((v) => console.log("got", v));
console.log("first");
\end{lstlisting}
{\small\textit{Prints:}~\code{first}\enspace\textcolor{black!45}{|}\enspace \code{got 8}}\par

\cardend
\cardhead{Card 60 --- Fix getJson(url)}{fix · javascript · Medium}
\safemargin{\textbf{Terms for this card}\par\smallskip \textbf{Exception:} an error object thrown to signal that something went wrong. Like a fire alarm: work stops and everyone heads for the nearest exit.}
\noindent{\small\textbf{Covered in:} \hyperref[ch:c13]{Ch.~13} of this guide; documentation: \href{https://developer.mozilla.org/en-US/docs/Web/API/Window/fetch}{MDN — Window: fetch()}.}\par
\cardsub{The card}
\textbf{Fix getJson(url)}\par
\code{getJson(url)} returns the error page's body when the server answers 404 or 500. Throw \code{new Error(}HTTP \$\{status\}\code{)} (e.g. \code{\textquotedbl{}HTTP 404\textquotedbl{}}) when the response isn't ok. Don't change the fake \code{fetch}.\par
\begin{lstlisting}[language=JavaScript]
async function getJson(url) {
  const response = await fetch(url);
  return response.json();
}

// Fake server for the tests: don't change.
async function fetch(url) {
  const routes = {
    "/users/1": { status: 200, body: { id: 1, name: "Ada" } },
    "/down": { status: 500, body: { error: "server error" } },
  };
  const route = routes[url] ?? { status: 404, body: { error: "not found" } };
  return {
    ok: route.status >= 200 && route.status < 300,
    status: route.status,
    json: async () => route.body,
  };
}
\end{lstlisting}
\cardsub{How to solve it}
\begin{cardrows}
\item[Solve] After \code{await fetch(url)}, throw \code{new Error(}HTTP \$\{response.status\}\code{)} unless \code{response.ok}.
\item[Invariants] Only 200–299 responses reach \code{json()}.
\item[Memory] One response object.
\item[Types] string in; promise of the body out, or a rejection.
\item[Failure] Checking only for 404 lets 500 through.
\item[Tests] An OK URL, a missing one (404) and a failing one (500).
\end{cardrows}
\trap fetch resolves on HTTP errors; check response.ok.\par
\cardsub{Solution}
\begin{lstlisting}[language=JavaScript]
async function getJson(url) {
  const response = await fetch(url);
  if (!response.ok) throw new Error(`HTTP ${response.status}`);
  return response.json();
}

// Fake server for the tests: don't change.
async function fetch(url) {
  const routes = {
    "/users/1": { status: 200, body: { id: 1, name: "Ada" } },
    "/down": { status: 500, body: { error: "server error" } },
  };
  const route = routes[url] ?? { status: 404, body: { error: "not found" } };
  return {
    ok: route.status >= 200 && route.status < 300,
    status: route.status,
    json: async () => route.body,
  };
}
\end{lstlisting}
{\small \textit{Verified: passes all 3 test cases on the card in the app's code runner; the starter code fails them.}}\par
\smallskip{\small\textbf{Why:} An HTTP error is still a successful fetch: the promise resolves with the response. \code{response.ok} is true only for 200–299, so check it before reading the body.}\par
\cardsub{From the readings}
\readchap{13. How does JavaScript wait without stopping?}{A promise stands for a result that isn't ready yet; \code{await} pauses only its own async function; promise callbacks run before timers; \code{Promise.all} waits for every promise and fails on the first rejection.}
\readq{Why must you check response.ok?}
Because \code{fetch} rejects only when the request itself fails (a bad URL, no network). A 404 or 500 answer still \emph{fulfills} with a response. \code{response.ok} is true only for status codes 200--299, so check it and throw yourself.
\Needspace{6\baselineskip}
\begin{lstlisting}[language=JavaScript]
const response = await fetch(url);
if (!response.ok) throw new Error(`HTTP ${response.status}`);
const data = await response.json();
\end{lstlisting}

\readchap{10. What should a program do when something goes wrong?}{\code{throw} stops the function and travels up the calls until a \code{catch} handles it; catch only where you can do something useful, and \code{finally} runs whichever way the block exits.}
\readq{What happens when you throw?}
The function stops and an \textbf{exception} travels up the call stack to the first \code{catch} that waits for it; with no \code{catch}, the program ends. Throw the most specific built-in error: \code{RangeError} for a value outside the allowed range. An error has a \code{name} and a \code{message}.
\Needspace{6\baselineskip}
\begin{lstlisting}[language=JavaScript]
function checkAge(age) {
  if (age < 0) throw new RangeError("negative age");
  return age;
}
try {
  checkAge(-1);
} catch (error) {
  console.log(`${error.name}: ${error.message}`);
}
\end{lstlisting}
{\small\textit{Prints:}~\code{RangeError: negative age}}\par

\cardend
\cardhead{Card 61 --- What does this print?}{predict · js13-all-vs-allsettled}
\noindent{\small\textbf{Covered in:} \hyperref[ch:c13]{Ch.~13} of this guide; documentation: \href{https://developer.mozilla.org/en-US/docs/Web/JavaScript/Reference/Global_Objects/Promise/all}{MDN — Promise.all()}.}\par
\cardsub{The card}
What does this print?\par
\begin{lstlisting}[language=JavaScript]
const ok = Promise.resolve("ok");
const fail = Promise.reject(new Error("boom"));
try {
  await Promise.all([ok, fail]);
} catch (error) {
  console.log("all:", error.message);
}
const results = await Promise.allSettled([ok, fail]);
console.log("allSettled:", results.map((r) => r.status).join(","));
\end{lstlisting}
\cardsub{How to solve it}
\begin{cardrows}
\item[Solve] \code{Promise.all} rejects with the first rejection ("boom"); \code{allSettled} lists a status per promise.
\item[Invariants] allSettled always fulfills.
\item[Memory] Two settled promises.
\item[Types] promises, strings.
\item[Failure] Expecting \code{Promise.all} to return the successful value.
\item[Tests] Run it, then change one line and predict again.
\end{cardrows}
\trap Promise.all fails as a whole when one promise fails.\par
\cardsub{Solution}
{\small\textbf{Prints:}}
\begin{lstlisting}
all: boom
allSettled: fulfilled,rejected
\end{lstlisting}
{\small\textit{Verified: Node v22.23.4 prints exactly this.}}\par
\smallskip{\small\textbf{Why:} Promise.all rejects with the first rejection. Promise.allSettled always fulfills, with a status for each promise. Neither one starts the work: the promises were already running when created.}\par
\cardsub{From the readings}
\readchap{13. How does JavaScript wait without stopping?}{A promise stands for a result that isn't ready yet; \code{await} pauses only its own async function; promise callbacks run before timers; \code{Promise.all} waits for every promise and fails on the first rejection.}
\readq{How do you run several requests at once?}
Start them all, then wait for all. A promise's work starts when the promise is created; \code{Promise.all} only waits. It fulfills with an array of the values in input order, or rejects with the first rejection. \code{await} inside a loop waits for each request before starting the next, so the requests run one at a time. \code{Promise.allSettled} always fulfills, with a status for each promise.
\Needspace{6\baselineskip}
\begin{lstlisting}[language=JavaScript]
const wait = (ms, v) => new Promise((r) => setTimeout(() => r(v), ms));
const start = Date.now();
const both = await Promise.all([wait(100, "a"), wait(100, "b")]);
console.log(both.join(""), Date.now() - start < 190);
\end{lstlisting}
{\small\textit{Prints:}~\code{ab true}}\par

\cardend
\deckhead{14 · Modules}{3}
\begin{corebox}\textbf{Core idea.} Each file is a module with its own scope: \code{export} shares a name, \code{import} borrows it---named imports in braces, the default import without---and a module runs once, so its importers share its state.\end{corebox}
\begin{sumbox}{Deck summary sheet --- 14 · Modules}
Each file is a module with its own scope. \code{export} makes a name available; \code{import \{ name \}} takes a named export (renamed with \code{as}), \code{import name} takes the default export. A module runs once, so every importer shares its state.\par
\smallskip{\footnotesize Read more: \href{https://developer.mozilla.org/en-US/docs/Web/JavaScript/Guide/Modules}{MDN — JavaScript modules}}\end{sumbox}
\cardhead{Card 62 --- What does this print?}{predict · js14-default-and-named}
\safemargin{\textbf{Terms for this card}\par\smallskip \textbf{Module:} a file with its own scope that shares only what it exports. Like a room with labelled hatches: only what is passed through a hatch leaves the room.\par\smallskip \textbf{Named export:} a value exported under its own name and imported in braces, \code{import \{ PI \}}. A hatch with a nameplate.\par\smallskip \textbf{Default export:} the one unnamed export of a module, imported without braces under any name you choose. The main hatch.}
\noindent{\small\textbf{Covered in:} \hyperref[ch:c14]{Ch.~14} of this guide; documentation: \href{https://developer.mozilla.org/en-US/docs/Web/JavaScript/Guide/Modules}{MDN — JavaScript modules}.}\par
\cardsub{The card}
What does running \code{main.js} print?\par
\code{math.js}\par
\begin{lstlisting}[language=JavaScript]
export default function add(a, b) {
  return a + b;
}
export const PI = 3.14;
\end{lstlisting}
\code{main.js}\par
\begin{lstlisting}[language=JavaScript]
import plus, { PI } from "./math.js";
console.log(plus(2, 3), PI);
\end{lstlisting}
\cardsub{How to solve it}
\begin{cardrows}
\item[Solve] \code{plus} is the default export (any name works); \code{PI} is a named export in braces.
\item[Invariants] There is at most one default export per module.
\item[Memory] One module evaluated once.
\item[Types] function and number.
\item[Failure] Thinking the default import must be called \code{add}.
\item[Tests] Run it, then change one line and predict again.
\end{cardrows}
\trap Default imports take any name; named imports must match.\par
\cardsub{Solution}
{\small\textbf{Prints:}}
\begin{lstlisting}
5 3.14
\end{lstlisting}
{\small\textit{Verified: Node v22.23.4 prints exactly this.}}\par
\smallskip{\small\textbf{Why:} The default export is imported without braces, under any name you choose. Named exports go in braces and must match the exported name.}\par
\cardsub{From the readings}
\readchap{14. How do separate files share code?}{Each file is a module with its own scope: \code{export} shares a name, \code{import} borrows it---named imports in braces, the default import without---and a module runs once, so its importers share its state.}
\readq{How do you share a function between files?}
Export it from one \textbf{module} and import it in another. A \textbf{named export} is imported in braces under its exact name (rename with \code{as}). A module can have one \textbf{default export}, imported without braces under any name.
\Needspace{6\baselineskip}
\begin{lstlisting}[language=JavaScript]
// shapes.js
export const sides = { triangle: 3 };
export default function area(w, h) {
  return w * h;
}
// main.js
import area, { sides } from "./shapes.js";
console.log(area(2, 3), sides.triangle);
\end{lstlisting}
{\small\textit{Prints:}~\code{6 3}}\par

\cardend
\cardhead{Card 63 --- What does this print?}{predict · js14-missing-named-export}
\safemargin{\textbf{Terms for this card}\par\smallskip \textbf{Module:} a file with its own scope that shares only what it exports. Like a room with labelled hatches: only what is passed through a hatch leaves the room.\par\smallskip \textbf{Named export:} a value exported under its own name and imported in braces, \code{import \{ PI \}}. A hatch with a nameplate.\par\smallskip \textbf{Default export:} the one unnamed export of a module, imported without braces under any name you choose. The main hatch.}
\noindent{\small\textbf{Covered in:} \hyperref[ch:c14]{Ch.~14} of this guide; documentation: \href{https://developer.mozilla.org/en-US/docs/Web/JavaScript/Reference/Statements/import}{MDN — import}.}\par
\cardsub{The card}
What does running \code{main.js} print?\par
\code{math.js}\par
\begin{lstlisting}[language=JavaScript]
export default function add(a, b) {
  return a + b;
}
\end{lstlisting}
\code{main.js}\par
\begin{lstlisting}[language=JavaScript]
import { add } from "./math.js";
console.log(add(2, 3));
\end{lstlisting}
\cardsub{How to solve it}
\begin{cardrows}
\item[Solve] \code{add} is a default export, so \code{\{ add \}} asks for a named export that does not exist: a SyntaxError before anything runs.
\item[Invariants] Imports are checked before the module's code runs.
\item[Memory] Nothing runs.
\item[Types] SyntaxError.
\item[Failure] Expecting it to print 5, or undefined.
\item[Tests] Run it, then change one line and predict again.
\end{cardrows}
\trap A default export is not a named export called add.\par
\cardsub{Solution}
{\small\textbf{Prints:}}
\begin{lstlisting}
SyntaxError: The requested module './math.js' does not provide an export named 'add'
\end{lstlisting}
{\small\textit{Verified: Node v22.23.4 prints exactly this.}}\par
\smallskip{\small\textbf{Why:} add is the default export, not a named one. Imports are checked before any code runs, so nothing is printed. Write \code{import add from \textquotedbl{}./math.js\textquotedbl{}}.}\par
\cardsub{From the readings}
\readchap{14. How do separate files share code?}{Each file is a module with its own scope: \code{export} shares a name, \code{import} borrows it---named imports in braces, the default import without---and a module runs once, so its importers share its state.}
\readq{What happens if you import a name that isn't exported?}
Nothing runs. Imports are checked before any code in the importing file runs, so asking for a named export that doesn't exist---for example asking for \code{\{ add \}} when \code{add} is the default export---fails at once with a \code{SyntaxError}.
\readq{How do you share a function between files?}
Export it from one \textbf{module} and import it in another. A \textbf{named export} is imported in braces under its exact name (rename with \code{as}). A module can have one \textbf{default export}, imported without braces under any name.
\Needspace{6\baselineskip}
\begin{lstlisting}[language=JavaScript]
// shapes.js
export const sides = { triangle: 3 };
export default function area(w, h) {
  return w * h;
}
// main.js
import area, { sides } from "./shapes.js";
console.log(area(2, 3), sides.triangle);
\end{lstlisting}
{\small\textit{Prints:}~\code{6 3}}\par

\cardend
\cardhead{Card 64 --- What does this print?}{predict · js14-shared-state}
\noindent{\small\textbf{Covered in:} \hyperref[ch:c14]{Ch.~14} of this guide; documentation: \href{https://developer.mozilla.org/en-US/docs/Web/JavaScript/Reference/Statements/import}{MDN — import}.}\par
\cardsub{The card}
What does running \code{main.js} print?\par
\code{counter.js}\par
\begin{lstlisting}[language=JavaScript]
let count = 0;
export function increment() {
  count += 1;
  return count;
}
\end{lstlisting}
\code{main.js}\par
\begin{lstlisting}[language=JavaScript]
import { increment as next } from "./counter.js";
import { increment } from "./counter.js";
next();
next();
console.log(increment());
\end{lstlisting}
\cardsub{How to solve it}
\begin{cardrows}
\item[Solve] Both imports are the same function with the same \code{count}: two calls make 2, the third returns 3.
\item[Invariants] A module runs once; its variables are shared.
\item[Memory] One \code{count} variable in counter.js.
\item[Types] number.
\item[Failure] Answering 1 (thinking each import gets its own copy).
\item[Tests] Run it, then change one line and predict again.
\end{cardrows}
\trap Renaming an import doesn't copy the module.\par
\cardsub{Solution}
{\small\textbf{Prints:}}
\begin{lstlisting}
3
\end{lstlisting}
{\small\textit{Verified: Node v22.23.4 prints exactly this.}}\par
\smallskip{\small\textbf{Why:} A module is evaluated once. Both imports are the same function, closing over the same count.}\par
\cardsub{From the readings}
\readchap{14. How do separate files share code?}{Each file is a module with its own scope: \code{export} shares a name, \code{import} borrows it---named imports in braces, the default import without---and a module runs once, so its importers share its state.}
\readq{If two files import the same module, do they get two copies?}
No. A module is evaluated once, and every import is a live, read-only view of its exports. Two imports of the same function share the same variables behind it.
\Needspace{6\baselineskip}
\begin{lstlisting}[language=JavaScript]
// log.js
export const lines = [];
export function log(x) {
  lines.push(x);
}
// main.js
import { log, lines } from "./log.js";
import { log as record } from "./log.js";
log("a");
record("b");
console.log(lines.join(","));
\end{lstlisting}
{\small\textit{Prints:}~\code{a,b}}\par

\cardend
\addtocontents{toc}{\vspace{6pt}}
\Needspace{14\baselineskip}\phantomsection\section*{Appendix: Audit Trail}\addcontentsline{toc}{section}{Appendix: Audit Trail}
\subsection*{Sources}
{\footnotesize\begin{tabularx}{\linewidth}{@{}p{1.6in}p{0.9in}X@{}}\toprule
\textbf{Source} & \textbf{How much read} & \textbf{Used for}\\\midrule
Course deck: \path{backend/app/content/js_ts_course/part1.py}, \path{quotes.py} & In full & Cards, tests, solutions, quotes, deck summaries\\
MDN JavaScript Guide (source text): Loops and iteration; Functions; Closures; Keyed collections; Using promises; JavaScript modules; Using classes; Control flow and error handling; Data types and structures & Searched by keyword; relevant sections read & Chapter facts; card quotes\\
MDN reference pages (source text): typeof, const, let, if...else, strict equality, logical OR and AND, nullish coalescing, Number.isInteger, parseInt, Array.isArray, map, filter, reduce, forEach, sort, reverse, String slice, split, trim, toUpperCase, addition, remainder, for, for...of, for...in, return, rest parameters, spread, destructuring, optional chaining, in, Set, Map, Object.fromEntries, try...catch, throw, Error, RangeError, SyntaxError, JSON.parse, JSON.stringify, extends, super, arrow functions, bind, this, Promise.all, import, Window.fetch & Opening paragraphs read; searched by keyword & Chapter facts; card quotes\\
MDN Glossary: Hoisting; Recursion & In full & Chapter facts; card quotes\\
\bottomrule\end{tabularx}}

{\footnotesize\textbf{Not read:} the ECMAScript specification; MDN pages beyond those listed (for example Expressions and operators, Iterators and generators, Working with objects); the TypeScript documentation (used in Parts 2--4, not here).}
\subsection*{Claim ledger}
{\footnotesize Each card's explanation rests on the documentation excerpt shown with it; the table lists, per source page, the cards whose claim it backs. Code behaviour claimed in the chapters and solutions was established by running it (see Checks run).}
{\footnotesize\begin{tabularx}{\linewidth}{@{}p{2.2in}X@{}}\toprule
\textbf{Source (all quoted, checked word for word)} & \textbf{Cards whose claim it backs}\\\midrule
MDN — Addition (+) (\href{https://developer.mozilla.org/en-US/docs/Web/JavaScript/Reference/Operators/Addition}{page}) & js04-plus-coercion\\
MDN — Array.isArray() (\href{https://developer.mozilla.org/en-US/docs/Web/JavaScript/Reference/Global_Objects/Array/isArray}{page}) & js01-describe-type\\
MDN — Array.prototype.filter() (\href{https://developer.mozilla.org/en-US/docs/Web/JavaScript/Reference/Global_Objects/Array/filter}{page}) & js02-count-truthy\\
MDN — Array.prototype.forEach() (\href{https://developer.mozilla.org/en-US/docs/Web/JavaScript/Reference/Global_Objects/Array/forEach}{page}) & js06-foreach-returns-undefined\\
MDN — Array.prototype.map() (\href{https://developer.mozilla.org/en-US/docs/Web/JavaScript/Reference/Global_Objects/Array/map}{page}) & js06-squares-of-evens\\
MDN — Array.prototype.reduce() (\href{https://developer.mozilla.org/en-US/docs/Web/JavaScript/Reference/Global_Objects/Array/reduce}{page}) & js06-group-by-category\\
MDN — Array.prototype.reduce(), exceptions: TypeError (\href{https://developer.mozilla.org/en-US/docs/Web/JavaScript/Reference/Global_Objects/Array/reduce}{page}) & js06-average\\
MDN — Array.prototype.reverse() (\href{https://developer.mozilla.org/en-US/docs/Web/JavaScript/Reference/Global_Objects/Array/reverse}{page}) & js04-palindrome\\
MDN — Array.prototype.sort() (\href{https://developer.mozilla.org/en-US/docs/Web/JavaScript/Reference/Global_Objects/Array/sort}{page}) & js06-sort-default\\
MDN — Arrow function expressions (\href{https://developer.mozilla.org/en-US/docs/Web/JavaScript/Reference/Functions/Arrow_functions}{page}) & js12-this-in-callback\\
MDN — Closures (\href{https://developer.mozilla.org/en-US/docs/Web/JavaScript/Guide/Closures}{page}) & js03-closure-in-loop, js05-make-counter, js05-once\\
MDN — Control flow and error handling (\href{https://developer.mozilla.org/en-US/docs/Web/JavaScript/Guide/Control_flow_and_error_handling}{page}) & js10-parse-age\\
MDN — Destructuring (\href{https://developer.mozilla.org/en-US/docs/Web/JavaScript/Reference/Operators/Destructuring}{page}) & js07-format-user\\
MDN — Error (\href{https://developer.mozilla.org/en-US/docs/Web/JavaScript/Reference/Global_Objects/Error}{page}) & js10-custom-error\\
MDN — Function.prototype.bind() (\href{https://developer.mozilla.org/en-US/docs/Web/JavaScript/Reference/Global_Objects/Function/bind}{page}) & js12-lost-this\\
MDN — Functions (\href{https://developer.mozilla.org/en-US/docs/Web/JavaScript/Guide/Functions}{page}) & js01-pass-by-sharing, js09-count-files, js09-trace\\
MDN — Hoisting (\href{https://developer.mozilla.org/en-US/docs/Glossary/Hoisting}{page}) & js05-hoisting\\
MDN — JSON.parse() (\href{https://developer.mozilla.org/en-US/docs/Web/JavaScript/Reference/Global_Objects/JSON/parse}{page}) & js10-dont-swallow, js11-load-settings\\
MDN — JSON.stringify() (\href{https://developer.mozilla.org/en-US/docs/Web/JavaScript/Reference/Global_Objects/JSON/stringify}{page}) & js11-stringify-drops, js11-pretty, js11-map-to-json\\
MDN — JavaScript data types (\href{https://developer.mozilla.org/en-US/docs/Web/JavaScript/Guide/Data_structures}{page}) & js01-typeof, js04-immutable\\
MDN — JavaScript modules (\href{https://developer.mozilla.org/en-US/docs/Web/JavaScript/Guide/Modules}{page}) & js14-default-and-named\\
MDN — Keyed collections (\href{https://developer.mozilla.org/en-US/docs/Web/JavaScript/Guide/Keyed_collections}{page}) & js08-map-keys\\
\bottomrule\end{tabularx}}
{\footnotesize\begin{tabularx}{\linewidth}{@{}p{2.2in}X@{}}\toprule
\textbf{Source (all quoted, checked word for word)} & \textbf{Cards whose claim it backs}\\\midrule
MDN — Logical OR (||) (\href{https://developer.mozilla.org/en-US/docs/Web/JavaScript/Reference/Operators/Logical_OR}{page}) & js02-logical-operands\\
MDN — Map (\href{https://developer.mozilla.org/en-US/docs/Web/JavaScript/Reference/Global_Objects/Map}{page}) & js08-most-frequent-word\\
MDN — Nullish coalescing operator (??) (\href{https://developer.mozilla.org/en-US/docs/Web/JavaScript/Reference/Operators/Nullish_coalescing}{page}) & js02-port-default\\
MDN — Number.isInteger() (\href{https://developer.mozilla.org/en-US/docs/Web/JavaScript/Reference/Global_Objects/Number/isInteger}{page}) & js01-parse-count\\
MDN — Object.fromEntries() (\href{https://developer.mozilla.org/en-US/docs/Web/JavaScript/Reference/Global_Objects/Object/fromEntries}{page}) & js08-group-by-map\\
MDN — Optional chaining (?.) (\href{https://developer.mozilla.org/en-US/docs/Web/JavaScript/Reference/Operators/Optional_chaining}{page}) & js07-optional-chaining\\
MDN — Promise.all() (\href{https://developer.mozilla.org/en-US/docs/Web/JavaScript/Reference/Global_Objects/Promise/all}{page}) & js13-load-names, js13-all-vs-allsettled\\
MDN — Recursion (\href{https://developer.mozilla.org/en-US/docs/Glossary/Recursion}{page}) & js09-factorial, js09-sum-nested\\
MDN — Remainder (\%) (\href{https://developer.mozilla.org/en-US/docs/Web/JavaScript/Reference/Operators/Remainder}{page}) & js03-fizzbuzz\\
MDN — Rest parameters (\href{https://developer.mozilla.org/en-US/docs/Web/JavaScript/Reference/Functions/rest_parameters}{page}) & js05-sum-rest\\
MDN — Set (\href{https://developer.mozilla.org/en-US/docs/Web/JavaScript/Reference/Global_Objects/Set}{page}) & js08-unique\\
MDN — Spread syntax (\href{https://developer.mozilla.org/en-US/docs/Web/JavaScript/Reference/Operators/Spread_syntax}{page}) & js06-shared-reference, js07-shallow-spread, js07-with-defaults\\
MDN — Strict equality (===) (\href{https://developer.mozilla.org/en-US/docs/Web/JavaScript/Reference/Operators/Strict_equality}{page}) & js02-equality\\
MDN — String.prototype.slice() (\href{https://developer.mozilla.org/en-US/docs/Web/JavaScript/Reference/Global_Objects/String/slice}{page}) & js04-capitalize\\
MDN — String.prototype.split() (\href{https://developer.mozilla.org/en-US/docs/Web/JavaScript/Reference/Global_Objects/String/split}{page}) & js04-initials\\
MDN — Using classes (\href{https://developer.mozilla.org/en-US/docs/Web/JavaScript/Guide/Using_classes}{page}) & js12-bank-account\\
MDN — Using promises (\href{https://developer.mozilla.org/en-US/docs/Web/JavaScript/Guide/Using_promises}{page}) & js13-event-loop\\
MDN — Window: fetch() (\href{https://developer.mozilla.org/en-US/docs/Web/API/Window/fetch}{page}) & js13-fetch-ok\\
MDN — const (\href{https://developer.mozilla.org/en-US/docs/Web/JavaScript/Reference/Statements/const}{page}) & js01-const\\
MDN — extends (\href{https://developer.mozilla.org/en-US/docs/Web/JavaScript/Reference/Classes/extends}{page}) & js12-inheritance\\
MDN — for (\href{https://developer.mozilla.org/en-US/docs/Web/JavaScript/Reference/Statements/for}{page}) & js03-sum-to\\
MDN — for...in (\href{https://developer.mozilla.org/en-US/docs/Web/JavaScript/Reference/Statements/for...in}{page}) & js03-for-in-vs-of\\
\bottomrule\end{tabularx}}
{\footnotesize\begin{tabularx}{\linewidth}{@{}p{2.2in}X@{}}\toprule
\textbf{Source (all quoted, checked word for word)} & \textbf{Cards whose claim it backs}\\\midrule
MDN — for...of (\href{https://developer.mozilla.org/en-US/docs/Web/JavaScript/Reference/Statements/for...of}{page}) & js03-count-vowels\\
MDN — if...else (\href{https://developer.mozilla.org/en-US/docs/Web/JavaScript/Reference/Statements/if...else}{page}) & js02-letter-grade\\
MDN — import (\href{https://developer.mozilla.org/en-US/docs/Web/JavaScript/Reference/Statements/import}{page}) & js14-missing-named-export, js14-shared-state\\
MDN — in (\href{https://developer.mozilla.org/en-US/docs/Web/JavaScript/Reference/Operators/in}{page}) & js07-pick\\
MDN — let (\href{https://developer.mozilla.org/en-US/docs/Web/JavaScript/Reference/Statements/let}{page}) & js01-tdz\\
MDN — return (\href{https://developer.mozilla.org/en-US/docs/Web/JavaScript/Reference/Statements/return}{page}) & js03-first-index-over\\
MDN — try...catch (\href{https://developer.mozilla.org/en-US/docs/Web/JavaScript/Reference/Statements/try...catch}{page}) & js10-finally\\
\bottomrule\end{tabularx}}
{\footnotesize\textbf{Authored, not sourced:} every analogy in a term definition (47 terms), every rule of thumb, every Think-it-through box, the chapter prose that connects the sourced facts, the How-to-solve rows and traps, and all chapter code snippets (their behaviour was run, see below). The cards themselves (tasks, tests, solutions) were written for this course earlier in this session and are verified, not sourced.}
\subsection*{Flags}
{\footnotesize\begin{itemize}
\item Outside the material: how Node prints values (for example \code{undefined} and error names) was established by running Node, not taken from the documentation.
\item Conflicts: none found between the deck and the documentation.
\item Simplification: Ch.~1 says a \code{var} exists from the top of its function; at the top level of a script it belongs to the global scope instead. The material's examples are inside functions or modules.
\item The js06-average quote is one item of MDN's list of exceptions (``Thrown if\ldots''); its subject, TypeError, is named in the source title instead.
\item Analogies are authored and simplify: a reference is not literally an address, and a closure keeps live variables that can still change, not a frozen snapshot as the snow-globe image suggests.
\item Doubtful premise: the deck's name says ``Zero to Mastery''; this Part alone covers the fundamentals.
\end{itemize}}
\subsection*{Checks run}
{\footnotesize\begin{itemize}
\item Code cards: 38 reference answers run in the app's own code runner against the cards' 156 test cases: all pass. Every starter fails. 0 TypeScript answers also type-checked (strict) with their type tests: all pass.
\item Predict cards: 26 snippets run in Node v22.23.4: each prints exactly the card's answer.
\item Chapter snippets: 47 run in Node v22.23.4; every \textit{Prints} line is the captured output. 0 checked with tsc 5.9.3.
\item Quotes: every card's quote found word for word (ignoring markdown) in the documentation source text (MDN content and TypeScript-Website repositories, fetched on 2026-11-16).
\item The course's own test suite (\code{tests/test\_js\_ts\_course.py}) passes, including planted wrong answers that each card must reject.
\item Deck files: \path{study-guides/js-ts/decks/part-1-javascript-fundamentals/} holds the 14 decks of this Part as JSON (a plain array and a wrapped copy each): 38 code cards and 26 qa cards, each with its verified quote. They are written by \path{export\_decks.py} from the same course data the checks above ran on. All 27 deck files of the course were imported through the app's \code{POST /api/import/json} into a fresh test database: no warnings, every quote became a citation, and every code card's answer passed all its test cases (and its type check) through the imported card's run-code endpoint.
\item Not run: the cards in your own app with the local 3B grading model (grading of typed answers is not verified here).
\item Not run: Node 24 on Windows; checks used Node v22.23.4 on Linux.
\item Not checked: the live web pages (the quotes were checked against the pages' source files, which can lag or lead the site).
\end{itemize}}
\subsection*{Deviations}
{\footnotesize\begin{itemize}
\item Summary sheets: the course has no separate summary sheets, so each deck's own documentation (the summary stored with the deck in the app) is copied in its box instead.
\item \emph{From the readings} copies only the chapter sections each card drills (listed per card), not every section of its chapters, to keep the guide usable.
\item The card decks were not generated anew: the guide uses the 105-card deck built earlier in this session (the same material), exported as JSON deck files alongside this guide.
\item The skill requires a quote on every card; 68 cards had none, so verified quotes were added to the course itself (\code{quotes.py}), where the app shows them too.
\item Read-only cards (what does this print? / does it compile?) are typed-answer cards in the app; their \emph{Solution} shows the verified output or compiler verdict, and their How-to-solve rows read: Solve = how to trace it, Tests = how to check it by running.
\item Chapters follow the deck's subdeck order rather than a strict stasis order, so card and chapter numbers match the app.
\item Chapter 14's snippets need two files; each listing shows both files, which were run as separate modules.
\item One chapter snippet (Ch.~13, checking response.ok) was not run: it needs a live server. It is the same check as card js13-fetch-ok, which was run against a fake server.
\end{itemize}}
\subsection*{Revision log}
{\footnotesize\begin{itemize}
\item 2026-11-16: first version.
\end{itemize}}
\addtocontents{toc}{\vspace{6pt}}
\phantomsection\section*{Index of Terms}\addcontentsline{toc}{section}{Index of Terms}
{\small\begin{multicols}{2}
\ix{accumulator}{accumulator}{Ch.~6}
\ix{array}{array}{Ch.~6}
\ix{arrow}{arrow function}{Ch.~5}
\ix{async}{async function}{Ch.~13}
\ix{basecase}{base case}{Ch.~9}
\ix{callstack}{call stack}{Ch.~9}
\ix{callback}{callback}{Ch.~6}
\ix{class}{class}{Ch.~12}
\ix{closure}{closure}{Ch.~3}
\ix{coercion}{coercion}{Ch.~2}
\ix{defaultexport}{default export}{Ch.~14}
\ix{destructuring}{Destructuring}{Ch.~7}
\ix{exception}{exception}{Ch.~10}
\ix{function}{function}{Ch.~5}
\ix{hoisting}{hoisted}{Ch.~5}
\ix{immutable}{immutable}{Ch.~4}
\ix{inplace}{in place}{Ch.~6}
\ix{instance}{instance}{Ch.~12}
\ix{iteration}{iteration}{Ch.~3}
\ix{json}{JSON}{Ch.~11}
\ix{loop}{loop}{Ch.~3}
\ix{map}{Map}{Ch.~8}
\ix{method}{method}{Ch.~4}
\ix{microtask}{microtasks}{Ch.~13}
\ix{module}{module}{Ch.~14}
\ix{namedexport}{named export}{Ch.~14}
\ix{nullish}{nullish coalescing}{Ch.~2}
\ix{object}{object}{Ch.~1}
\ix{offbyone}{off-by-one error}{Ch.~3}
\ix{optional}{optional chaining}{Ch.~7}
\ix{parameter}{parameters}{Ch.~5}
\ix{primitive}{primitive}{Ch.~1}
\ix{privatefield}{private field}{Ch.~12}
\ix{promise}{promise}{Ch.~13}
\ix{property}{property}{Ch.~7}
\ix{prototype}{prototype chain}{Ch.~8}
\ix{recursion}{Recursion}{Ch.~9}
\ix{reference}{reference}{Ch.~1}
\ix{rest}{rest parameter}{Ch.~5}
\ix{set}{Set}{Ch.~8}
\ix{shallow}{shallow copy}{Ch.~7}
\ix{strict}{strict equality}{Ch.~2}
\ix{tdz}{temporal dead zone}{Ch.~1}
\ix{truthy}{truthy}{Ch.~2}
\ix{typeof}{typeof}{Ch.~1}
\ix{value}{value}{Ch.~1}
\ix{variable}{variable}{Ch.~1}
\end{multicols}}
\end{document}
